\documentclass[english]{article}
\usepackage[T1]{fontenc}
\usepackage[utf8]{inputenc}
\usepackage{amsmath}
\usepackage{amssymb}

\makeatletter
\usepackage[a4paper,margin=15mm]{geometry}
\allowdisplaybreaks[1]
\setlength{\emergencystretch}{1em}

\usepackage{needspace}
\usepackage[unicode,hidelinks,linktoc=all,bookmarksopen=true]{hyperref}
\renewcommand{\l@section}[2]{\par\addvspace{7pt}\noindent\parbox[t]{\dimexpr\linewidth-14mm\relax}{\raggedright\normalfont #1}\hfill\makebox[10mm][r]{\normalfont #2}\par}

\makeatother

\usepackage{babel}
\begin{document}
\vspace*{10mm}
\begin{center}
{\huge\textbf{Detailed Derivations for the}}{\huge\par}
\par\end{center}

\begin{center}
{\LARGE\textbf{``Comment on `Casimir Energy in Kerr Space-time'\thinspace{}''}}{\LARGE\par}
\par\end{center}

\begin{center}
{\Large Shaohui Yao}{\Large\par}
\par\end{center}

\begin{center}
{\large November 6, 2026}{\large\par}
\par\end{center}

\begin{center}
{\large Version 1.0.0}{\large\par}
\par\end{center}

\begin{center}\rule{55mm}{0.4pt}\end{center}\vspace{4mm}

\begin{center}\begin{minipage}{\linewidth}\fontsize{11}{14}\selectfont\tableofcontents\end{minipage}\end{center}\clearpage

\Needspace{110pt}\par\addvspace{20pt}\noindent\begin{minipage}{\linewidth}\begingroup\fontsize{12}{15}\selectfont\setlength{\abovedisplayskip}{0pt}\setlength{\belowdisplayskip}{0pt}\setlength{\abovedisplayshortskip}{0pt}\setlength{\belowdisplayshortskip}{0pt}\phantomsection\addcontentsline{toc}{section}{I. Four-Velocity, Normalization Factor, and Comoving Coordinates}
\[
\begin{gathered}\text{\bfseries I. Four-Velocity, Normalization Factor,}\\
\text{\bfseries and Comoving Coordinates}
\end{gathered}
\]
\endgroup\end{minipage}\par\nopagebreak[4]\vspace{20pt}\nopagebreak[4]

In the coordinate system $\left(t,r,\theta,\phi\right)$ (the order
of $\theta,\phi$ is reversed relative to that in the paper), the
four-velocity is 
\[
w^{\mu}=\left(\frac{dt}{d\tau},\frac{dr}{d\tau},\frac{d\theta}{d\tau},\frac{d\phi}{d\tau}\right)
\]
 On the equatorial plane, $\theta=\frac{\pi}{2},\frac{d\theta}{d\tau}=0$,
and for a circular orbit, $r=r_{0},\frac{dr}{d\tau}=0$. Hence, 
\[
w^{\mu}=\left(\frac{dt}{d\tau},0,0,\frac{d\phi}{d\tau}\right)
\]

The angular velocity $\Omega=\frac{d\phi}{dt}$ is the angular velocity
measured by an observer at infinity, using coordinate time. Thus,
\begin{align*}
\frac{d\phi}{d\tau} & =\frac{d\phi}{dt}\cdot\frac{dt}{d\tau}\\
 & =\Omega\cdot\frac{dt}{d\tau}
\end{align*}

We obtain 
\[
w^{\mu}=\frac{dt}{d\tau}\left(1,0,0,\Omega\right)
\]

$\delta^{\mu}_{\!t}=\left(1,0,0,0\right)$ and $\delta^{\mu}_{\!\phi}=\left(0,0,0,1\right)$,
where $\delta$ is the Kronecker delta. Hence, 
\[
w^{\mu}=\frac{dt}{d\tau}\left(\delta^{\mu}_{\!t}+\Omega\delta^{\mu}_{\!\phi}\right)
\]

The normalization condition is $w^{\mu}w_{\mu}=1$. (The paper uses
the metric signature $(+,-,-,-)$, which is opposite to the convention
commonly used in general relativity; the normalization condition is
therefore not $w^{\mu}w_{\mu}=-1$.)

\[
g_{\mu\rho}w^{\rho}=w_{\mu}
\]

\begin{align*}
w_{t} & =g_{tt}w^{t}+g_{tr}w^{r}+g_{t\theta}w^{\theta}+g_{t\phi}w^{\phi}\\
 & =g_{tt}\frac{dt}{d\tau}+0+0+\Omega g_{t\phi}\frac{dt}{d\tau}\\
 & =\frac{dt}{d\tau}\left(g_{tt}+\Omega g_{t\phi}\right)
\end{align*}

\begin{align*}
w_{r} & =g_{rt}w^{t}+g_{rr}w^{r}+g_{r\theta}w^{\theta}+g_{r\phi}w^{\phi}\\
 & =\frac{dt}{d\tau}\left(g_{rt}+\Omega g_{r\phi}\right)
\end{align*}

\[
w_{\theta}=\frac{dt}{d\tau}\left(g_{\theta t}+\Omega g_{\theta\phi}\right)
\]

\[
w_{\phi}=\frac{dt}{d\tau}\left(g_{\phi t}+\Omega g_{\phi\phi}\right)
\]

\begin{align*}
w^{\mu}w_{\mu} & =w^{t}w_{t}+w^{r}w_{r}+w^{\theta}w_{\theta}+w^{\phi}w_{\phi}\\
 & =\frac{dt}{d\tau}\cdot\frac{dt}{d\tau}\left(g_{tt}+\Omega g_{t\phi}\right)+0+0+\Omega\frac{dt}{d\tau}\cdot\frac{dt}{d\tau}\left(g_{\phi t}+\Omega g_{\phi\phi}\right)\\
 & =\frac{dt}{d\tau}\cdot\frac{dt}{d\tau}\left[g_{tt}+\Omega g_{t\phi}+\Omega\left(g_{\phi t}+\Omega g_{\phi\phi}\right)\right]\\
 & =\left(\frac{dt}{d\tau}\right)^{2}\left(g_{tt}+\Omega g_{t\phi}+\Omega g_{\phi t}+\Omega^{2}g_{\phi\phi}\right)
\end{align*}

\begin{align*}
\left(\frac{dt}{d\tau}\right)^{2}\left(g_{tt}+\Omega g_{t\phi}+\Omega g_{\phi t}+\Omega^{2}g_{\phi\phi}\right) & =1\\
\Rightarrow\frac{dt}{d\tau} & =\sqrt{\frac{1}{g_{tt}+\Omega g_{t\phi}+\Omega g_{\phi t}+\Omega^{2}g_{\phi\phi}}}\\
 & =\left(g_{tt}+\Omega g_{t\phi}+\Omega g_{\phi t}+\Omega^{2}g_{\phi\phi}\right)^{-1/2}\\
 & =\left(g_{tt}+2\Omega g_{t\phi}+\Omega^{2}g_{\phi\phi}\right)^{-1/2}
\end{align*}

Here, $g_{\phi t}=g_{t\phi}$. Defining $C\left(\Omega\right):=\left(g_{tt}+2\Omega g_{t\phi}+\Omega^{2}g_{\phi\phi}\right)^{-1/2}=dt/d\tau$,
we have 
\begin{align*}
w^{\mu} & =\frac{dt}{d\tau}\left(\delta^{\mu}_{\!t}+\Omega\delta^{\mu}_{\!\phi}\right)\\
 & =\left(g_{tt}+2\Omega g_{t\phi}+\Omega^{2}g_{\phi\phi}\right)^{-1/2}\left(\delta^{\mu}_{\!t}+\Omega\delta^{\mu}_{\!\phi}\right)\\
 & =C\left(\Omega\right)\left(\delta^{\mu}_{\!t}+\Omega\delta^{\mu}_{\!\phi}\right)
\end{align*}

Solving $g_{tt}+2\Omega g_{t\phi}+\Omega^{2}g_{\phi\phi}=0$ (using
$x=\frac{-b\pm\sqrt{b^{2}-4ac}}{2a}$ for the quadratic equation $ax^{2}+bx+c=0$)
gives 
\begin{align*}
\Omega^{'}_{\pm} & =\frac{-2g_{t\phi}\pm\sqrt{4g^{2}_{t\phi}-4g_{\phi\phi}g_{tt}}}{2g_{\phi\phi}}\\
 & =-\frac{g_{t\phi}}{g_{\phi\phi}}\pm\frac{\sqrt{g^{2}_{t\phi}-g_{\phi\phi}g_{tt}}}{g_{\phi\phi}}
\end{align*}

The line element is:

\begin{align*}
ds^{2} & =\left(1-\frac{2Mr}{\Sigma}\right)dt^{2}+2\frac{2Mar}{\Sigma}\text{sin}^{2}\theta dtd\phi-\frac{\Sigma}{\Delta}dr^{2}-\Sigma d\theta^{2}-\frac{A}{\Sigma}\text{sin}^{2}\theta d\phi^{2}\\
 & =g_{tt}dt^{2}+2g_{t\phi}dtd\phi+g_{rr}dr^{2}+g_{\theta\theta}d\theta^{2}+g_{\phi\phi}d\phi^{2}
\end{align*}

where 
\begin{align*}
\Sigma & =r^{2}+a^{2}\text{cos}^{2}\theta,\\
\Delta & =r^{2}+a^{2}-2Mr,\\
A & =\left(r^{2}+a^{2}\right)^{2}-a^{2}\Delta\text{sin}^{2}\theta\\
 & =\left(r^{2}+a^{2}\right)\Sigma+2Mra^{2}\text{sin}^{2}\theta,\\
a & =\frac{J}{M}.
\end{align*}

The metric matrix is 
\[
g_{\mu\nu}=\begin{pmatrix}g_{tt} & g_{tr} & g_{t\theta} & g_{t\phi}\\
g_{rt} & g_{rr} & g_{r\theta} & g_{r\phi}\\
g_{\theta t} & g_{\theta r} & g_{\theta\theta} & g_{\theta\phi}\\
g_{\phi t} & g_{\phi r} & g_{\phi\theta} & g_{\phi\phi}
\end{pmatrix}
\]

Here, $g_{t\phi}=g_{\phi t}=\frac{2Mar}{\Sigma}\text{sin}^{2}\theta$,
and 
\begin{align*}
g_{t\phi}dtd\phi+g_{\phi t}d\phi dt & =\frac{2Mar}{\Sigma}\text{sin}^{2}\theta dtd\phi+\frac{2Mar}{\Sigma}\text{sin}^{2}\theta d\phi dt\\
 & =\frac{4Mar}{\Sigma}\text{sin}^{2}\theta dtd\phi
\end{align*}

The coordinate transformation is:

\[
\begin{cases}
t'=t\\
r'=r\\
\theta'=\theta\\
\phi'=\phi-\Omega t
\end{cases}
\]

The inverse transformation is:

\[
\begin{cases}
t=t'\\
r=r'\\
\theta=\theta'\\
\phi=\phi'+\Omega t'
\end{cases}
\]

This gives

\begin{align*}
ds^{2} & =g_{tt}dt^{2}+2g_{t\phi}dtd\phi+g_{rr}dr^{2}+g_{\theta\theta}d\theta^{2}+g_{\phi\phi}d\phi^{2}\\
 & =g_{tt}dt'^{2}+2g_{t\phi}dt'd\left(\phi^{'}+\Omega t'\right)+g_{rr}dr'^{2}+g_{\theta\theta}d\theta'^{2}+g_{\phi\phi}d\left(\phi^{'}+\Omega t'\right)d\left(\phi^{'}+\Omega t'\right)\\
 & =g_{tt}dt'^{2}+2g_{t\phi}dt'\left(d\phi^{'}+\Omega dt'\right)+g_{rr}dr'^{2}+g_{\theta\theta}d\theta'^{2}+g_{\phi\phi}\left(d\phi^{'}+\Omega dt'\right)\left(d\phi^{'}+\Omega dt'\right)\\
 & =g_{tt}dt'^{2}+2g_{t\phi}\left(dt'd\phi^{'}+\Omega dt'^{2}\right)+g_{rr}dr'^{2}+g_{\theta\theta}d\theta'^{2}+g_{\phi\phi}\left(d\phi^{'2}+\Omega d\phi^{'}dt'+\Omega dt'd\phi^{'}+\Omega^{2}dt'^{2}\right)\\
 & =g_{tt}dt'^{2}+2g_{t\phi}\left(dt'd\phi^{'}+\Omega dt'^{2}\right)+g_{rr}dr'^{2}+g_{\theta\theta}d\theta'^{2}+g_{\phi\phi}\left(d\phi^{'2}+2\Omega dt'd\phi^{'}+\Omega^{2}dt'^{2}\right)\\
 & =\left(g_{tt}+2\Omega g_{t\phi}+\Omega^{2}g_{\phi\phi}\right)dt'^{2}+2\left(g_{t\phi}+\Omega g_{\phi\phi}\right)dt'd\phi^{'}+g_{rr}dr'^{2}+g_{\theta\theta}d\theta'^{2}+g_{\phi\phi}d\phi^{'2}\\
 & =g^{'}_{\mu\nu}dx^{'\mu}dx^{'\nu}
\end{align*}

\[
\begin{cases}
g^{'}_{tt}=\left(g_{tt}+2\Omega g_{t\phi}+\Omega^{2}g_{\phi\phi}\right)\\
g^{'}_{t\phi}=g^{'}_{\phi t}=\left(g_{t\phi}+\Omega g_{\phi\phi}\right)\\
g^{'}_{rr}=g_{rr}\\
g^{'}_{\theta\theta}=g_{\theta\theta}\\
g^{'}_{\phi\phi}=g_{\phi\phi}
\end{cases}
\]

When the angular velocity $\Omega$ of an observer moving along a
circular equatorial orbit satisfies $g_{t\phi}+\Omega g_{\phi\phi}=0$,
or equivalently $\Omega=-\frac{g_{t\phi}}{g_{\phi\phi}}$, we have
$2\left(g_{t\phi}+\Omega g_{\phi\phi}\right)dt'd\phi^{'}=0$. The
$dt'd\phi^{'}$ cross term is eliminated, and there is no frame-dragging
effect in this frame (the local reference frame of the observer on
the circular equatorial orbit). Such an observer is also called a
zero-angular-momentum observer (ZAMO). We define this angular velocity
as 
\[
\omega_{d}=-\frac{g_{t\phi}}{g_{\phi\phi}}
\]

\begin{align*}
\omega_{d} & =-\frac{g_{t\phi}}{g_{\phi\phi}}\\
 & =-\frac{\frac{2Mar}{\Sigma}\text{sin}^{2}\theta}{-\frac{A}{\Sigma}\text{sin}^{2}\theta}\\
 & =\frac{2Mar}{A}
\end{align*}
 

This corresponds to the first term in $\Omega_{\pm}$. For the second
term, on an equatorial orbit $\theta=\pi/2$, with the metric components
treated as constants (see the detailed analysis below), we have 
\begin{align*}
\left.\frac{\sqrt{g^{2}_{t\phi}-g_{\phi\phi}g_{tt}}}{g_{\phi\phi}}\right|_{\theta=\pi/2} & =\left.\frac{\sqrt{\left(\frac{2Mar}{\Sigma}\text{sin}^{2}\theta\right)^{2}-\left(-\frac{A}{\Sigma}\text{sin}^{2}\theta\right)\left(1-\frac{2Mr}{\Sigma}\right)}}{-\frac{A}{\Sigma}\text{sin}^{2}\theta}\right|_{\theta=\pi/2}\\
 & =-\frac{\Sigma\sqrt{\Delta}}{A}
\end{align*}

where

\begin{align*}
 & \biggl(\frac{2Mar}{\Sigma}\text{sin}^{2}\theta\biggr)^{2}-\biggl(-\frac{A}{\Sigma}\text{sin}^{2}\theta\biggr)\biggl(1-\frac{2Mr}{\Sigma}\biggr)\bigg|_{\theta=\pi/2}=\biggl(\frac{2Mar}{\Sigma}\text{sin}^{2}\theta\biggr)^{2}-\biggl(-\frac{A}{\Sigma}\text{sin}^{2}\theta\biggr)\biggl(1-\frac{2Mr}{\Sigma}\biggr)\\
 & =\biggl(\frac{2Mar}{\Sigma}\text{sin}^{2}\theta\biggr)^{2}+\frac{1}{\Sigma}\biggl[\biggl(r^{2}+a^{2}\biggr)\Sigma+2Mra^{2}\text{sin}^{2}\theta\biggr]\text{sin}^{2}\theta\biggl(1-\frac{2Mr}{\Sigma}\biggr)\\
 & =\biggl(\frac{2Mar}{\Sigma}\text{sin}^{2}\theta\biggr)^{2}+\biggl[\biggl(r^{2}+a^{2}\biggr)+\frac{2Mra^{2}}{\Sigma}\text{sin}^{2}\theta\biggr]\text{sin}^{2}\theta\biggl(1-\frac{2Mr}{\Sigma}\biggr)\\
 & =\biggl(\frac{2Mar}{\Sigma}\text{sin}^{2}\theta\biggr)^{2}+\biggl(r^{2}+a^{2}\biggr)\text{sin}^{2}\theta\biggl(1-\frac{2Mr}{\Sigma}\biggr)+\biggl[\frac{2Mra^{2}}{\Sigma}\biggl(\text{sin}^{2}\theta\biggr)^{2}-\biggl(\frac{2Mar}{\Sigma}\text{sin}^{2}\theta\biggr)^{2}\biggr]\\
 & =\biggl(r^{2}+a^{2}\biggr)\text{sin}^{2}\theta\biggl(1-\frac{2Mr}{\Sigma}\biggr)+\frac{2Mra^{2}}{\Sigma}\biggl(\text{sin}^{2}\theta\biggr)^{2}\\
 & =\biggl(r^{2}+a^{2}\biggr)\text{sin}^{2}\theta-\biggl(r^{2}+a^{2}\biggr)\text{sin}^{2}\theta\cdot\frac{2Mr}{\Sigma}+\frac{2Mra^{2}}{\Sigma}\biggl(\text{sin}^{2}\theta\biggr)^{2}\\
 & =\biggl[\biggl(r^{2}+a^{2}\biggr)\text{sin}^{2}\theta-r^{2}\frac{2Mr}{\Sigma}\text{sin}^{2}\theta-\frac{2Mra^{2}}{\Sigma}\text{sin}^{2}\theta+\frac{2Mra^{2}}{\Sigma}\biggl(\text{sin}^{2}\theta\biggr)^{2}\biggr]\bigg|_{\theta=\pi/2}\\
 & =\biggl(r^{2}+a^{2}\biggr)-r^{2}\frac{2Mr}{\Sigma|_{\theta=\pi/2}}-\frac{2Mra^{2}}{\Sigma|_{\theta=\pi/2}}+\frac{2Mra^{2}}{\Sigma|_{\theta=\pi/2}}\\
 & =r^{2}+a^{2}-2Mr\\
 & =\Delta
\end{align*}

Here, $\Sigma|_{\theta=\pi/2}=r^{2}+a^{2}\text{cos}^{2}\theta\Big|_{\theta=\pi/2}=r^{2}$.

Method 2 (preferred):

\begin{align*}
g^{2}_{t\phi}-g_{\phi\phi}g_{tt} & =\left(\frac{2Mar}{\Sigma}\text{sin}^{2}\theta\right)^{2}-\left(-\frac{A}{\Sigma}\text{sin}^{2}\theta\right)\left(1-\frac{2Mr}{\Sigma}\right)\\
 & =\left(\frac{2Mar}{\Sigma}\text{sin}^{2}\theta\right)^{2}-\left(-\frac{A}{\Sigma}\text{sin}^{2}\theta\right)\left(\frac{\Sigma}{\Sigma}-\frac{2Mr}{\Sigma}\right)\\
 & =\frac{\text{sin}^{2}\theta}{\Sigma^{2}}\left[\left(2Mar\right)^{2}\text{sin}^{2}\theta+A\left(\Sigma-2Mr\right)\right]\\
 & =\Delta\text{sin}^{2}\theta
\end{align*}

where 
\begin{align*}
\left(2Mar\right)^{2}\text{sin}^{2}\theta+A\left(\Sigma-2Mr\right) & =\left(2Mar\right)^{2}\text{sin}^{2}\theta+A\left(\Sigma-2Mr\right)\\
 & =\left(2Mar\right)^{2}\text{sin}^{2}\theta+\left[\left(r^{2}+a^{2}\right)\Sigma+2Mra^{2}\text{sin}^{2}\theta\right]\left(\Sigma-2Mr\right)\\
 & =\left(2Mar\right)^{2}\text{sin}^{2}\theta+\left(r^{2}+a^{2}\right)\Sigma\left(\Sigma-2Mr\right)+2Mra^{2}\text{sin}^{2}\theta\left(\Sigma-2Mr\right)\\
 & =\left(r^{2}+a^{2}\right)\Sigma\left(\Sigma-2Mr\right)+\left(2Mar\right)^{2}\text{sin}^{2}\theta+2Mra^{2}\text{sin}^{2}\theta\left(\Sigma-2Mr\right)\\
 & =\left(r^{2}+a^{2}\right)\Sigma\left(\Sigma-2Mr\right)+2Mra^{2}\text{sin}^{2}\theta\left[2Mr+\left(\Sigma-2Mr\right)\right]\\
 & =\Sigma\left[\left(r^{2}+a^{2}\right)\left(\Sigma-2Mr\right)+2Mra^{2}\text{sin}^{2}\theta\right]\\
 & =\Sigma\left[\left(r^{2}+a^{2}\right)\left(\Delta-a^{2}\sin^{2}\theta\right)+2Mra^{2}\text{sin}^{2}\theta\right]\\
 & =\Sigma\left[\left(r^{2}+a^{2}\right)\Delta-\left(r^{2}+a^{2}\right)a^{2}\sin^{2}\theta+2Mra^{2}\text{sin}^{2}\theta\right]\\
 & =\Sigma\left[\left(r^{2}+a^{2}\right)\Delta-\left(r^{2}+a^{2}-2Mr\right)a^{2}\sin^{2}\theta\right]\\
 & =\Sigma\left[\left(r^{2}+a^{2}\right)\Delta-\Delta a^{2}\sin^{2}\theta\right]\\
 & =\Sigma\Delta\left[r^{2}+a^{2}\left(1-\sin^{2}\theta\right)\right]\\
 & =\Sigma\Delta\left(r^{2}+a^{2}\cos^{2}\theta\right)\\
 & =\Sigma^{2}\Delta
\end{align*}

where

\begin{align*}
\Sigma-2Mr & =r^{2}+a^{2}\cos^{2}\theta-2Mr\\
 & =r^{2}+a^{2}\left(1-\sin^{2}\theta\right)-2Mr\\
 & =r^{2}+a^{2}-2Mr-a^{2}\sin^{2}\theta\\
 & =\Delta-a^{2}\sin^{2}\theta
\end{align*}

This gives 
\begin{align*}
\left.\frac{\sqrt{g^{2}_{t\phi}-g_{\phi\phi}g_{tt}}}{g_{\phi\phi}}\right|_{\theta=\pi/2} & =\left.\frac{\sqrt{\left(\frac{2Mar}{\Sigma}\text{sin}^{2}\theta\right)^{2}-\left(-\frac{A}{\Sigma}\text{sin}^{2}\theta\right)\left(1-\frac{2Mr}{\Sigma}\right)}}{-\frac{A}{\Sigma}\text{sin}^{2}\theta}\right|_{\theta=\pi/2}\\
 & =\left.\frac{\sqrt{\Delta\text{sin}^{2}\theta}}{-\frac{A}{\Sigma}\text{sin}^{2}\theta}\right|_{\theta=\pi/2}\\
 & =-\frac{\Sigma\sqrt{\Delta}}{A}
\end{align*}

We obtain 
\begin{align*}
\Omega^{'}_{\pm} & =-\frac{g_{t\phi}}{g_{\phi\phi}}\pm\frac{\sqrt{g^{2}_{t\phi}-g_{\phi\phi}g_{tt}}}{g_{\phi\phi}}\\
 & =\omega_{d}\mp\frac{\Sigma\sqrt{\Delta}}{A}
\end{align*}

Define 
\[
\Omega_{\pm}=\omega_{d}\pm\frac{\Sigma\sqrt{\Delta}}{A}
\]

The condition $r>r_{+}$ is imposed because the orbit must lie outside
the outer horizon of the black hole.

\begin{align*}
\left.g_{tt}+2\Omega g_{t\phi}+\Omega^{2}g_{\phi\phi}\right|_{\theta=\pi/2} & =g_{\phi\phi}\left(\frac{g_{tt}}{g_{\phi\phi}}+2\Omega\frac{g_{t\phi}}{g_{\phi\phi}}+\Omega^{2}\right)\\
 & =g_{\phi\phi}\left(\frac{g_{tt}}{g_{\phi\phi}}-2\Omega\omega_{d}+\Omega^{2}\right)\\
 & =g_{\phi\phi}\left[\frac{g_{tt}}{g_{\phi\phi}}+\left(\Omega-\omega_{d}\right)^{2}-\omega^{2}_{d}\right]\\
 & =\frac{g_{\phi\phi}g_{tt}}{g_{\phi\phi}}+g_{\phi\phi}\left(\Omega-\omega_{d}\right)^{2}-g_{\phi\phi}\left(-\frac{g_{t\phi}}{g_{\phi\phi}}\right)^{2}\\
 & =g_{\phi\phi}\left(\Omega-\omega_{d}\right)^{2}+\frac{g_{\phi\phi}g_{tt}-g^{2}_{t\phi}}{g_{\phi\phi}}\\
 & =-\frac{A}{\Sigma}\text{sin}^{2}\theta\left(\Omega-\omega_{d}\right)^{2}+\frac{-\Delta\sin^{2}\theta}{-\frac{A}{\Sigma}\text{sin}^{2}\theta}\\
 & =-\frac{A}{\Sigma}\text{sin}^{2}\theta\left(\Omega-\omega_{d}\right)^{2}+\frac{\Delta\Sigma}{A}\\
 & =\left.\frac{\Delta\Sigma}{A}\left[1-\frac{A^{2}}{\Delta\Sigma^{2}}\text{sin}^{2}\theta\left(\Omega-\omega_{d}\right)^{2}\right]\right|_{\theta=\pi/2}\\
 & =\frac{\Delta\Sigma}{A}\left[1-\frac{A^{2}}{\Delta\Sigma^{2}}\left(\Omega-\omega_{d}\right)^{2}\right]
\end{align*}

\begin{align*}
C\left(\Omega\right) & =\left(g_{tt}+2\Omega g_{t\phi}+\Omega^{2}g_{\phi\phi}\right)^{-1/2}\\
 & =\left(\left.g_{tt}+2\Omega g_{t\phi}+\Omega^{2}g_{\phi\phi}\right|_{\theta=\pi/2}\right)^{-1/2}\\
 & =\left\{ \frac{\Delta\Sigma}{A}\left[1-\frac{A^{2}}{\Delta\Sigma^{2}}\left(\Omega-\omega_{d}\right)^{2}\right]\right\} ^{-1/2}
\end{align*}
 Alternatively, 
\begin{align*}
C\left(\Omega\right) & =\left\{ \frac{\Delta\Sigma|_{\theta=\pi/2}}{A}\left[1-\frac{A^{2}\text{sin}^{2}\theta|_{\theta=\pi/2}}{\Delta\Sigma|^{2}_{\theta=\pi/2}}\left(\Omega-\omega_{d}\right)^{2}\right]\right\} ^{-1/2}\\
 & =\left\{ \frac{\Delta r^{2}}{A}\left[1-\frac{A^{2}}{\Delta r^{4}}\left(\Omega-\omega_{d}\right)^{2}\right]\right\} ^{-1/2}
\end{align*}

For Eq. (16),

\[
\begin{cases}
t'=t\\
r'=r\\
\theta'=\theta\\
\phi'=\phi-\Omega t
\end{cases}
\]

\begin{align*}
w^{'\mu} & =\left(\frac{dt'}{d\tau},\frac{dr'}{d\tau},\frac{d\theta'}{d\tau},\frac{d\phi'}{d\tau}\right)\\
 & =\left(\frac{dt}{d\tau},\frac{dr}{d\tau},\frac{d\theta}{d\tau},\frac{d\left(\phi-\Omega t\right)}{d\tau}\right)\\
 & =\left(\frac{dt}{d\tau},\frac{dr}{d\tau},\frac{d\theta}{d\tau},\frac{d\phi}{d\tau}-\Omega\frac{dt}{d\tau}\right)\\
 & =\left(\frac{dt}{d\tau},0,0,\Omega\frac{dt}{d\tau}-\Omega\frac{dt}{d\tau}\right)\\
 & =C\left(\Omega\right)\left(1,0,0,0\right)\\
 & =C\left(\Omega\right)\delta^{\mu}_{\!t}
\end{align*}

Here, $\frac{d\phi}{d\tau}=\Omega\frac{dt}{d\tau}$ and $\frac{dt}{d\tau}=C\left(\Omega\right)$.
A scalar, such as  the inner product of the four-velocity $w^{\mu}w_{\mu}=1$
, is invariant under coordinate transformations. The transformed four-velocity
$w^{'\mu}$ therefore does not need to be renormalized. Specifically,
as shown above, the four-velocity becomes 
\[
w^{'\mu}=\frac{dt'}{d\tau}\delta^{\mu}_{\!t}
\]

The normalization condition $w^{'\mu}w^{'}_{\mu}=1$ gives 
\[
\frac{dt'}{d\tau}=\left(g^{'}_{tt}\right)^{-1/2}
\]

We also have the result derived above: 
\[
g^{'}_{tt}=g_{tt}+2\Omega g_{t\phi}+\Omega^{2}g_{\phi\phi}
\]

Hence, 
\begin{align*}
\frac{dt'}{d\tau} & =\left(g_{tt}+2\Omega g_{t\phi}+\Omega^{2}g_{\phi\phi}\right)^{-1/2}\\
 & =C\left(\Omega\right)\\
 & =\frac{dt}{d\tau}
\end{align*}

The arc length is $\Delta l=r\Delta\theta$, so $dx=rd\phi'$ and
$dz=rd\theta'=rd\theta$. We obtain 
\[
\begin{cases}
d\phi'=\frac{1}{r}dx\\
d\theta'=\frac{1}{r}dz\\
dr'=dy\\
dt'=dt
\end{cases}
\]

\begin{align*}
ds^{2} & =g^{'}_{\mu\nu}dx^{'\mu}dx^{'\nu}\\
 & =\left(g_{tt}+2\Omega g_{t\phi}+\Omega^{2}g_{\phi\phi}\right)dt'^{2}+2\left(g_{t\phi}+\Omega g_{\phi\phi}\right)dt'd\phi^{'}+g_{rr}dr'^{2}+g_{\theta\theta}d\theta'^{2}+g_{\phi\phi}d\phi^{'2}\\
 & =\left(g_{tt}+2\Omega g_{t\phi}+\Omega^{2}g_{\phi\phi}\right)dt^{2}+2\left(g_{t\phi}+\Omega g_{\phi\phi}\right)dt\left(\frac{1}{r}dx\right)+g_{rr}dy^{2}+g_{\theta\theta}\left(\frac{1}{r}dz\right)^{2}+g_{\phi\phi}\left(\frac{1}{r}dx\right)^{2}\\
 & =\left(g_{tt}+2\Omega g_{t\phi}+\Omega^{2}g_{\phi\phi}\right)dt^{2}+2g_{\phi\phi}\left(\frac{g_{t\phi}}{g_{\phi\phi}}+\Omega\right)dt\left(\frac{1}{r}dx\right)-\frac{\Sigma}{\Delta}dy^{2}-\Sigma\left(\frac{1}{r}dz\right)^{2}-\frac{A}{\Sigma}\text{sin}^{2}\theta\left(\frac{1}{r}dx\right)^{2}\\
 & =\left(g_{tt}+2\Omega g_{t\phi}+\Omega^{2}g_{\phi\phi}\right)dt^{2}+2\left(-\frac{A}{\Sigma}\text{sin}^{2}\theta\right)\left(-\omega_{d}+\Omega\right)dt\left(\frac{1}{r}dx\right)-\frac{\Sigma}{\Delta}dy^{2}-\Sigma\left(\frac{1}{r}dz\right)^{2}-\frac{A}{\Sigma}\text{sin}^{2}\theta\left(\frac{1}{r}dx\right)^{2}
\end{align*}
 When $\theta=\pi/2$, we have $\Sigma=r^{2}$, and 
\begin{align*}
ds^{2} & =C^{-2}\left(\Omega\right)dt^{2}-2\frac{A}{r^{3}}\left(\Omega-\omega_{d}\right)dtdx-\frac{r^{2}}{\Delta}dy^{2}-dz^{2}-\frac{A}{r^{4}}dx^{2}\\
 & =\hat{g}_{tt}dt^{2}+2\hat{g}_{tx}dtdx+\hat{g}_{yy}dy^{2}+\hat{g}_{zz}dz^{2}+\hat{g}_{xx}dx^{2}
\end{align*}

\begin{align*}
\hat{g}_{\mu\nu} & =\begin{pmatrix}C^{-2}\left(\Omega\right) & -\frac{A}{r^{3}}\left(\Omega-\omega_{d}\right) & 0 & 0\\
-\frac{A}{r^{3}}\left(\Omega-\omega_{d}\right) & -\frac{A}{r^{4}} & 0 & 0\\
0 & 0 & -\frac{r^{2}}{\Delta} & 0\\
0 & 0 & 0 & -1
\end{pmatrix}\\
 & =\begin{pmatrix}\hat{g}_{tt} & \hat{g}_{tx} & 0 & 0\\
\hat{g}_{xt} & \hat{g}_{xx} & 0 & 0\\
0 & 0 & \hat{g}_{yy} & 0\\
0 & 0 & 0 & \hat{g}_{zz}
\end{pmatrix}\\
 & =\begin{pmatrix}A & O\\
O & B
\end{pmatrix}
\end{align*}

Here, $\hat{g}_{tx}=\hat{g}_{xt}$.

\begin{align*}
\hat{g}_{tt}\hat{g}_{xx}-\hat{g}^{2}_{tx} & =C^{-2}\left(\Omega\right)\left(-\frac{A}{r^{4}}\right)-\left[\frac{A}{r^{3}}\left(\Omega-\omega_{d}\right)\right]^{2}\\
 & =\frac{\Delta\Sigma}{A}\left[1-\frac{A^{2}}{\Delta\Sigma^{2}}\left(\Omega-\omega_{d}\right)^{2}\right]\left(-\frac{A}{r^{4}}\right)-\left[\frac{A}{r^{3}}\left(\Omega-\omega_{d}\right)\right]^{2}\\
 & =\frac{\Delta r^{2}}{A}\left[1-\frac{A^{2}}{\Delta r^{4}}\left(\Omega-\omega_{d}\right)^{2}\right]\left(-\frac{A}{r^{4}}\right)-\left[\frac{A}{r^{3}}\left(\Omega-\omega_{d}\right)\right]^{2}\\
 & =-\frac{\Delta}{r^{2}}+\frac{A^{2}}{r^{6}}\left(\Omega-\omega_{d}\right)^{2}-\frac{A^{2}}{r^{6}}\left(\Omega-\omega_{d}\right)^{2}\\
 & =-\frac{\Delta}{r^{2}}
\end{align*}
 Here, for an equatorial orbit, $\theta=\pi/2$ and $\Sigma=r^{2}$.

\begin{align*}
\det\left(\hat{g}_{\mu\nu}\right) & =\begin{vmatrix}\hat{g}_{tt} & \hat{g}_{tx} & 0 & 0\\
\hat{g}_{xt} & \hat{g}_{xx} & 0 & 0\\
0 & 0 & \hat{g}_{yy} & 0\\
0 & 0 & 0 & \hat{g}_{zz}
\end{vmatrix}\\
 & =\begin{vmatrix}A & O\\
O & B
\end{vmatrix}\\
 & =\det\left(A\right)\cdot\det\left(B\right)\\
 & =\left(\hat{g}_{tt}\hat{g}_{xx}-\hat{g}^{2}_{tx}\right)\times\hat{g}_{yy}\times\hat{g}_{zz}\\
 & =-\frac{\Delta}{r^{2}}\times\left(-\frac{r^{2}}{\Delta}\right)\times\left(-1\right)\\
 & =-1
\end{align*}

This is equal to the determinant of the Minkowski metric, $\det\left(\eta_{\mu\nu}\right)=-1$,
for an equatorial orbit ($\theta=\pi/2$).

Inverting the metric:

\[
\hat{g}^{\mu\alpha}\hat{g}_{\alpha\nu}=\delta^{\mu}_{\!\nu}
\]

Equivalently, in matrix notation, 
\[
\hat{g}^{\mu\nu}=\left(\hat{g}_{\mu\nu}\right)^{-1}
\]

\[
\left(\hat{g}_{\mu\nu}\right)^{-1}=\begin{pmatrix}A^{-1} & O\\
O & B^{-1}
\end{pmatrix}
\]

\begin{align*}
\begin{pmatrix}\hat{g}_{tt} & \hat{g}_{tx}\\
\hat{g}_{xt} & \hat{g}_{xx}
\end{pmatrix}\begin{pmatrix}a & b\\
c & d
\end{pmatrix} & =\begin{pmatrix}1 & 0\\
0 & 1
\end{pmatrix}\\
\begin{pmatrix}a\hat{g}_{tt}+c\hat{g}_{tx} & b\hat{g}_{tt}+d\hat{g}_{tx}\\
a\hat{g}_{xt}+c\hat{g}_{xx} & b\hat{g}_{xt}+d\hat{g}_{xx}
\end{pmatrix} & =\begin{pmatrix}1 & 0\\
0 & 1
\end{pmatrix}
\end{align*}

\[
\begin{cases}
a\hat{g}_{tt}+c\hat{g}_{tx}=1\\
b\hat{g}_{tt}+d\hat{g}_{tx}=0\\
a\hat{g}_{xt}+c\hat{g}_{xx}=0\\
b\hat{g}_{xt}+d\hat{g}_{xx}=1
\end{cases}
\]

\begin{align*}
\left(a\hat{g}_{tt}+c\hat{g}_{tx}\right)+\left(a\hat{g}_{xt}+c\hat{g}_{xx}\right) & =1\\
a\left(\hat{g}_{tt}+\hat{g}_{xt}\right)+\frac{1-a\hat{g}_{tt}}{\hat{g}_{tx}}\left(\hat{g}_{tx}+\hat{g}_{xx}\right) & =1\\
a\left(\hat{g}_{tt}+\hat{g}_{xt}\right)+\frac{1}{\hat{g}_{tx}}\left(\hat{g}_{tx}+\hat{g}_{xx}\right)-\frac{a\hat{g}_{tt}}{\hat{g}_{tx}}\left(\hat{g}_{tx}+\hat{g}_{xx}\right) & =1\\
a\left(\hat{g}_{tt}+\hat{g}_{xt}\right)+\left(1+\frac{\hat{g}_{xx}}{\hat{g}_{tx}}\right)-a\hat{g}_{tt}\left(1+\frac{\hat{g}_{xx}}{\hat{g}_{tx}}\right) & =1\\
a\left[\left(\hat{g}_{tt}+\hat{g}_{xt}\right)-\hat{g}_{tt}\left(1+\frac{\hat{g}_{xx}}{\hat{g}_{tx}}\right)\right] & =1-\left(1+\frac{\hat{g}_{xx}}{\hat{g}_{tx}}\right)\\
a & =\frac{1-\left(1+\frac{\hat{g}_{xx}}{\hat{g}_{tx}}\right)}{\left(\hat{g}_{tt}+\hat{g}_{xt}\right)-\hat{g}_{tt}\left(1+\frac{\hat{g}_{xx}}{\hat{g}_{tx}}\right)}\\
a & =\frac{-\hat{g}_{xx}}{\hat{g}_{tx}\hat{g}_{xt}-\hat{g}_{tt}\hat{g}_{xx}}
\end{align*}

\begin{align*}
c & =\frac{-a\hat{g}_{xt}}{\hat{g}_{xx}}\\
 & =-\left(\frac{-\hat{g}_{xx}}{\hat{g}_{tx}\hat{g}_{xt}-\hat{g}_{tt}\hat{g}_{xx}}\right)\frac{\hat{g}_{xt}}{\hat{g}_{xx}}\\
 & =\frac{\hat{g}_{xt}}{\hat{g}_{tx}\hat{g}_{xt}-\hat{g}_{tt}\hat{g}_{xx}}
\end{align*}

\begin{align*}
\left(b\hat{g}_{tt}+d\hat{g}_{tx}\right)+\left(b\hat{g}_{xt}+d\hat{g}_{xx}\right) & =1\\
b\left(\hat{g}_{tt}+\hat{g}_{xt}\right)+\frac{-b\hat{g}_{tt}}{\hat{g}_{tx}}\left(\hat{g}_{tx}+\hat{g}_{xx}\right) & =1\\
b\left(\hat{g}_{tt}+\hat{g}_{xt}\right)-b\left(\hat{g}_{tt}+\frac{\hat{g}_{tt}\hat{g}_{xx}}{\hat{g}_{tx}}\right) & =1\\
b\left(\hat{g}_{xt}-\frac{\hat{g}_{tt}\hat{g}_{xx}}{\hat{g}_{tx}}\right) & =1\\
b & =\frac{\hat{g}_{tx}}{\hat{g}_{tx}\hat{g}_{xt}-\hat{g}_{tt}\hat{g}_{xx}}
\end{align*}

\begin{align*}
d & =\frac{-b\hat{g}_{tt}}{\hat{g}_{tx}}\\
 & =\frac{-\hat{g}_{tx}}{\hat{g}_{tx}\hat{g}_{xt}-\hat{g}_{tt}\hat{g}_{xx}}\frac{\hat{g}_{tt}}{\hat{g}_{tx}}\\
 & =\frac{-\hat{g}_{tt}}{\hat{g}_{tx}\hat{g}_{xt}-\hat{g}_{tt}\hat{g}_{xx}}
\end{align*}

\begin{align*}
A^{-1} & =\begin{pmatrix}a & b\\
c & d
\end{pmatrix}\\
 & =\frac{1}{\hat{g}_{tx}\hat{g}_{xt}-\hat{g}_{tt}\hat{g}_{xx}}\begin{pmatrix}-\hat{g}_{xx} & \hat{g}_{tx}\\
\hat{g}_{xt} & -\hat{g}_{tt}
\end{pmatrix}\\
 & =\frac{1}{\hat{g}^{2}_{tx}-\hat{g}_{tt}\hat{g}_{xx}}\begin{pmatrix}-\hat{g}_{xx} & \hat{g}_{tx}\\
\hat{g}_{tx} & -\hat{g}_{tt}
\end{pmatrix}
\end{align*}

\begin{align*}
\begin{pmatrix}\hat{g}_{yy} & 0\\
0 & \hat{g}_{zz}
\end{pmatrix}\begin{pmatrix}\alpha & \beta\\
\gamma & \delta
\end{pmatrix} & =\begin{pmatrix}1 & 0\\
0 & 1
\end{pmatrix}\\
\begin{pmatrix}\hat{g}_{yy}\alpha & \hat{g}_{yy}\beta\\
\hat{g}_{zz}\gamma & \hat{g}_{zz}\delta
\end{pmatrix} & =\begin{pmatrix}1 & 0\\
0 & 1
\end{pmatrix}
\end{align*}

\[
\begin{cases}
\alpha=\frac{1}{\hat{g}_{yy}}\\
\beta=0\\
\gamma=0\\
\delta=\frac{1}{\hat{g}_{zz}}
\end{cases}
\]

\begin{align*}
B^{-1} & =\begin{pmatrix}\alpha & \beta\\
\gamma & \delta
\end{pmatrix}\\
 & =\begin{pmatrix}\frac{1}{\hat{g}_{yy}} & 0\\
0 & \frac{1}{\hat{g}_{zz}}
\end{pmatrix}
\end{align*}

We obtain the inverse metric matrix:

\begin{align*}
\left(\hat{g}_{\mu\nu}\right)^{-1} & =\begin{pmatrix}A^{-1} & O\\
O & B^{-1}
\end{pmatrix}\\
 & =\begin{pmatrix}\frac{-\hat{g}_{xx}}{\hat{g}^{2}_{tx}-\hat{g}_{tt}\hat{g}_{xx}} & \frac{\hat{g}_{tx}}{\hat{g}^{2}_{tx}-\hat{g}_{tt}\hat{g}_{xx}} & 0 & 0\\
\frac{\hat{g}_{tx}}{\hat{g}^{2}_{tx}-\hat{g}_{tt}\hat{g}_{xx}} & \frac{-\hat{g}_{tt}}{\hat{g}^{2}_{tx}-\hat{g}_{tt}\hat{g}_{xx}} & 0 & 0\\
0 & 0 & \frac{1}{\hat{g}_{yy}} & 0\\
0 & 0 & 0 & \frac{1}{\hat{g}_{zz}}
\end{pmatrix}\\
 & =\begin{pmatrix}\frac{\frac{A}{r^{4}}}{\frac{\Delta}{r^{2}}} & \frac{-\frac{A}{r^{3}}\left(\Omega-\omega_{d}\right)}{\frac{\Delta}{r^{2}}} & 0 & 0\\
\frac{-\frac{A}{r^{3}}\left(\Omega-\omega_{d}\right)}{\frac{\Delta}{r^{2}}} & \frac{-C^{-2}\left(\Omega\right)}{\frac{\Delta}{r^{2}}} & 0 & 0\\
0 & 0 & \frac{1}{-\frac{r^{2}}{\Delta}} & 0\\
0 & 0 & 0 & \frac{1}{-1}
\end{pmatrix}\\
 & =\begin{pmatrix}\frac{A}{\Delta r^{2}} & \frac{A}{\Delta r}\left(\omega_{d}-\Omega\right) & 0 & 0\\
\frac{A}{\Delta r}\left(\omega_{d}-\Omega\right) & -\frac{r^{2}}{\Delta}C^{-2}\left(\Omega\right) & 0 & 0\\
0 & 0 & -\frac{\Delta}{r^{2}} & 0\\
0 & 0 & 0 & -1
\end{pmatrix}
\end{align*}

When the metric components are approximated as constants, the d'Alembertian
is 
\begin{align*}
\partial^{2}_{s} & =\hat{g}^{\mu\nu}\partial_{\mu}\partial_{\nu}\\
 & =\frac{A}{\Delta r^{2}}\partial^{2}_{t}+2\frac{A}{\Delta r}\left(\omega_{d}-\Omega\right)\partial_{t}\partial_{x}-\frac{r^{2}}{\Delta}C^{-2}\left(\Omega\right)\partial^{2}_{x}-\frac{\Delta}{r^{2}}\partial^{2}_{y}-\partial^{2}_{z}
\end{align*}
 where $\partial_{\mu}\partial_{\nu}=\left(\frac{\partial}{\partial x^{\mu}}\right)\left(\frac{\partial}{\partial x^{\nu}}\right)$.

The components of the four-velocity transform contravariantly under
a coordinate transformation. This follows because the coordinate differentials
$dx^{\mu}$ transform contravariantly, whereas the proper-time interval
$d\tau$ is a scalar. Under the transformation from the old coordinates
$x^{'\mu}$ to the new coordinates $\hat{x}^{\alpha}$, the chain
rule gives 
\[
d\hat{x}^{\alpha}=\frac{\partial\hat{x}^{\alpha}}{\partial x^{'\mu}}dx^{'\mu}
\]
 This is the transformation law relating coordinate differentials
in the two systems. Dividing both sides by $d\tau$ gives 
\[
\frac{d\hat{x}^{\alpha}}{d\tau}=\frac{\partial\hat{x}^{\alpha}}{\partial x^{'\mu}}\frac{dx^{'\mu}}{d\tau}
\]
 or equivalently, 
\[
\hat{w}^{\alpha}=\frac{\partial\hat{x}^{\alpha}}{\partial x^{'\mu}}w^{'\mu}
\]

From 
\[
\begin{cases}
d\phi'=\frac{1}{r}dx\\
d\theta'=\frac{1}{r}dz\\
dr'=dy\\
dt'=dt
\end{cases}
\]
 we obtain 
\[
\begin{cases}
\frac{\partial x}{\partial\phi'}=r\\
\frac{\partial z}{\partial\theta'}=r\\
\frac{\partial y}{\partial r'}=1\\
\frac{\partial t}{\partial t'}=1
\end{cases}
\]

\begin{align*}
\hat{w}^{\hat{t}} & =\frac{\partial\hat{t}}{\partial x^{'\mu}}w^{'\mu}\\
 & =\frac{\partial\hat{t}}{\partial t'}w^{'t^{'}}+\frac{\partial\hat{t}}{\partial r^{'}}w^{'r^{'}}+\frac{\partial\hat{t}}{\partial\theta^{'}}w^{'\theta^{'}}+\frac{\partial\hat{t}}{\partial\phi^{'}}w^{'\phi^{'}}\\
 & =\frac{\partial t}{\partial t}C\left(\Omega\right)+0\cdot0+0\cdot0+0\cdot0\\
 & =C\left(\Omega\right)
\end{align*}
 Here, $\hat{t}=t^{'}=t$. Since $\hat{t}=t^{'}$ does not depend
on $r^{'},\theta^{'}$ or $\phi^{'}$, we have $\partial\hat{t}/\partial r^{'}=\partial\hat{t}/\partial\theta^{'}=\partial\hat{t}/\partial\phi^{'}=0$.
\begin{align*}
\hat{w}^{x} & =\frac{\partial x}{\partial x^{'\mu}}w^{'\mu}\\
 & =\frac{\partial x}{\partial t'}w^{'t^{'}}+\frac{\partial x}{\partial r^{'}}w^{'r^{'}}+\frac{\partial x}{\partial\theta^{'}}w^{'\theta^{'}}+\frac{\partial x}{\partial\phi^{'}}w^{'\phi^{'}}\\
 & =0\cdot w^{'t^{'}}+\frac{\partial x}{\partial r^{'}}\cdot0+\frac{\partial x}{\partial\theta^{'}}\cdot0+r\cdot0\\
 & =0
\end{align*}
 Similarly, 
\[
\hat{w}^{y}=\hat{w}^{z}=0
\]
 Thus, 
\[
\hat{w}^{\mu}=w^{'\mu}
\]

\Needspace{110pt}\par\addvspace{20pt}\noindent\begin{minipage}{\linewidth}\begingroup\fontsize{12}{15}\selectfont\setlength{\abovedisplayskip}{0pt}\setlength{\belowdisplayskip}{0pt}\setlength{\abovedisplayshortskip}{0pt}\setlength{\belowdisplayshortskip}{0pt}\phantomsection\addcontentsline{toc}{section}{II. The Constant-Metric Approximation}
\[
\text{\bfseries II. The Constant-Metric Approximation}
\]
\endgroup\end{minipage}\par\nopagebreak[4]\vspace{20pt}\nopagebreak[4]

Taylor expansion of a multivariable function:

\[
f\left(\boldsymbol{x}\right)=f\left(\boldsymbol{x}_{0}\right)+\nabla f\left(\boldsymbol{x}_{0}\right)\cdot\left(\boldsymbol{x}-\boldsymbol{x}_{0}\right)+\mathcal{O}\left(\|\boldsymbol{x}-\boldsymbol{x}_{0}\|^{2}\right)
\]

For example, let $\boldsymbol{x}=\left(a,b,c\right)$ and $\boldsymbol{x}_{0}=\left(a_{0},b_{0},c_{0}\right)$.
Then 
\[
\nabla f\left(\boldsymbol{x}_{0}\right)=\left(\frac{\partial f\left(\boldsymbol{x}_{0}\right)}{\partial x^{1}},\frac{\partial f\left(\boldsymbol{x}_{0}\right)}{\partial x^{2}},\frac{\partial f\left(\boldsymbol{x}_{0}\right)}{\partial x^{3}}\right)
\]
 {}
\[
\left[\nabla f\left(\boldsymbol{x}_{0}\right)\right]_{i}=\partial_{i}f\left(\boldsymbol{x}_{0}\right)=\frac{\partial f\left(a,b,c\right)}{\partial x^{i}}\Bigg|_{\boldsymbol{x}=\boldsymbol{x}_{0}}
\]
 and 
\[
\|\boldsymbol{x}-\boldsymbol{x}_{0}\|=\sqrt{\left(a-a_{0}\right)^{2}+\left(b-b_{0}\right)^{2}+\left(c-c_{0}\right)^{2}}
\]
 Define 
\[
\delta\boldsymbol{x}=\boldsymbol{x}-\boldsymbol{x}_{0}=\left(a-a_{0},b-b_{0},c-c_{0}\right)
\]
 Then 
\begin{align*}
\nabla f\left(\boldsymbol{x}_{0}\right)\cdot\left(\boldsymbol{x}-\boldsymbol{x}_{0}\right) & =\frac{\partial f\left(a,b,c\right)}{\partial x^{i}}\Bigg|_{\boldsymbol{x}=\boldsymbol{x}_{0}}\cdot\delta x^{i}\\
 & =\frac{\partial f\left(a,b,c\right)}{\partial a}\Bigg|_{\boldsymbol{x}=\boldsymbol{x}_{0}}\left(a-a_{0}\right)+\frac{\partial f\left(a,b,c\right)}{\partial b}\Bigg|_{\boldsymbol{x}=\boldsymbol{x}_{0}}\left(b-b_{0}\right)+\frac{\partial f\left(a,b,c\right)}{\partial c}\Bigg|_{\boldsymbol{x}=\boldsymbol{x}_{0}}\left(c-c_{0}\right)
\end{align*}

If terms through second order are retained, the expansion is 
\[
f\left(\boldsymbol{x}\right)=f\left(\boldsymbol{x}_{0}\right)+\nabla f\left(\boldsymbol{x}_{0}\right)\cdot\left(\boldsymbol{x}-\boldsymbol{x}_{0}\right)+\frac{1}{2}\left(\boldsymbol{x}-\boldsymbol{x}_{0}\right)^{T}H\left(\boldsymbol{x}_{0}\right)\left(\boldsymbol{x}-\boldsymbol{x}_{0}\right)+\mathcal{O}\left(\|\boldsymbol{x}-\boldsymbol{x}_{0}\|^{3}\right)
\]
 where $H$ is the Hessian matrix, whose entries are 
\[
H_{ij}=\frac{\partial^{2}f}{\partial x_{i}\partial x_{j}}
\]

Equivalently, 
\begin{align*}
 & f\biggl(a,b,c\biggr)=f\biggl(a,b,c\biggr)\Bigg|_{\boldsymbol{x}=\boldsymbol{x}_{0}}+\sum^{3}_{i=1}\frac{\partial f\biggl(a,b,c\biggr)}{\partial x^{i}}\Bigg|_{\boldsymbol{x}=\boldsymbol{x}_{0}}\biggl(x^{i}-x^{i}_{0}\biggr)+\frac{1}{2}\sum^{3}_{i,j=1}\frac{\partial^{2}f\biggl(a,b,c\biggr)}{\partial x^{i}\partial x^{j}}\Bigg|_{\boldsymbol{x}=\boldsymbol{x}_{0}}\\
 & \qquad\biggl(x^{i}-x^{i}_{0}\biggr)\biggl(x^{j}-x^{j}_{0}\biggr)+\mathcal{O}\biggl(\biggl(\sqrt{\sum^{3}_{i=1}\biggl(x^{i}-x^{i}_{0}\biggr)^{2}}\biggr)^{3}\biggr)
\end{align*}
 By Taylor's theorem with the mean-value form of the remainder, $f\left(a,b,c\right)$
can also be written as 
\[
f\left(a,b,c\right)=f\left(a,b,c\right)\Bigg|_{\boldsymbol{x}=\boldsymbol{x}_{0}}+\sum^{3}_{i=1}\frac{\partial f\left(a,b,c\right)}{\partial x^{i}}\Bigg|_{\boldsymbol{x}=\boldsymbol{x}_{0}}\left(x^{i}-x^{i}_{0}\right)+\mathcal{R}
\]
 where 
\[
\mathcal{R}=\frac{1}{2}\sum^{3}_{i,j=1}\frac{\partial^{2}f\left(a,b,c\right)}{\partial x^{i}\partial x^{j}}\Bigg|_{\boldsymbol{x}=\boldsymbol{\xi}}\left(x^{i}-x^{i}_{0}\right)\left(x^{j}-x^{j}_{0}\right)
\]
 Here, $\boldsymbol{\xi}=\left(\xi_{1},\xi_{2},\xi_{3}\right)$, and
$\xi_{1},\ \xi_{2},\ \xi_{3}$ lie between $a_{0}$ and $a$, $b_{0}$
and $b$, and $c_{0}$ and $c$, respectively (or, equivalently, $\boldsymbol{\xi}=\boldsymbol{x}_{0}+s\left(\boldsymbol{x}-\boldsymbol{x}_{0}\right),\ 0<s<1$).

The Taylor expansion of the metric components $\hat{g}_{\mu\nu}$
about $\boldsymbol{x}_{0}$ is 
\begin{align*}
\hat{g}_{\mu\nu}\left(t,x,y,z\right) & =\hat{g}_{\mu\nu}\left(t,x,y,z\right)\big|_{\boldsymbol{x}=\boldsymbol{x}_{0}}+\sum^{3}_{i=1}\frac{\partial\hat{g}_{\mu\nu}\left(t,x,y,z\right)}{\partial x^{i}}\Bigg|_{\boldsymbol{x}=\boldsymbol{x}_{0}}\left(x^{i}-x^{i}_{0}\right)+\mathcal{O}\left(\|\boldsymbol{x}-\boldsymbol{x}_{0}\|^{2}\right)\\
 & =\hat{g}_{\mu\nu}\left(t,x_{0},y_{0},z_{0}\right)+\sum^{3}_{i=1}\frac{\partial\hat{g}_{\mu\nu}\left(t,x,y,z\right)}{\partial x^{i}}\Bigg|_{\boldsymbol{x}=\boldsymbol{x}_{0}}\cdot\delta x^{i}+\mathcal{O}\left(\|\delta\boldsymbol{x}\|^{2}\right)
\end{align*}
 or 
\begin{align*}
 & \hat{g}_{\mu\nu}\biggl(t,x,y,z\biggr)=\hat{g}_{\mu\nu}\biggl(t,x,y,z\biggr)\big|_{\boldsymbol{x}=\boldsymbol{x}_{0}}+\sum^{3}_{i=1}\frac{\partial\hat{g}_{\mu\nu}\biggl(t,x,y,z\biggr)}{\partial x^{i}}\Bigg|_{\boldsymbol{x}=\boldsymbol{x}_{0}}\biggl(x^{i}-x^{i}_{0}\biggr)+\frac{1}{2}\sum^{3}_{i,j=1}\\
 & \qquad\frac{\partial^{2}\hat{g}_{\mu\nu}\biggl(t,x,y,z\biggr)}{\partial x^{i}\partial x^{j}}\Bigg|_{\boldsymbol{x}=\boldsymbol{\xi}}\biggl(x^{i}-x^{i}_{0}\biggr)\biggl(x^{j}-x^{j}_{0}\biggr)
\end{align*}
 where $\boldsymbol{\xi}=\left(\xi_{1},\xi_{2},\xi_{3}\right)$, and
$\xi_{1},\ \xi_{2},\ \xi_{3}$ lie between $x_{0}$ and $x$, $y_{0}$
and $y$, and $z_{0}$ and $z$, respectively (or, equivalently, $\boldsymbol{\xi}=\boldsymbol{x}_{0}+s\left(\boldsymbol{x}-\boldsymbol{x}_{0}\right),\ 0<s<1$).
Here, $\boldsymbol{x}=\left(x,y,z\right)$ and $\boldsymbol{x}_{0}=\left(x_{0},y_{0},z_{0}\right)$.
Define $\delta x^{i}=x^{i}-x^{i}_{0}$, or equivalently $x^{i}=x^{i}_{0}+\delta x^{i}$;
explicitly, $\delta x=x-x_{0},\delta y=y-y_{0},\delta z=z-z_{0}$.
Also define $\delta\boldsymbol{x}=\boldsymbol{x}-\boldsymbol{x}_{0}$.
Let $\boldsymbol{x}_{0}=\left(x_{0},y_{0},z_{0}\right)$ correspond
to $\boldsymbol{x}_{0}=\left(r_{0},\theta_{0},\phi_{0}\right)$ in
the coordinate system $\left(t,r,\theta,\phi\right)$. Under the coordinate
transformations 
\[
\begin{cases}
t'=t\\
r'=r\\
\theta'=\theta\\
\phi'=\phi-\Omega t
\end{cases}
\]
 and 
\[
\begin{cases}
d\bar{t}=dt'=dt\\
dx=r_{0}d\phi'=r_{0}d\left(\phi-\Omega t\right)=r_{0}d\phi-\Omega r_{0}dt\\
dy=dr'=dr\\
dz=r_{0}d\theta'=r_{0}d\theta
\end{cases}
\]
 the variations near the point $\boldsymbol{x}_{0}=\left(r_{0},\theta_{0},\phi_{0}\right)$
are 
\[
\begin{cases}
\delta x=\left(r_{0}\phi-\Omega r_{0}t\right)-\left(r_{0}\phi_{0}-\Omega r_{0}t\right)=r_{0}\left(\phi-\phi_{0}\right)\\
\delta y=r-r_{0}\\
\delta z=r_{0}\left(\theta-\theta_{0}\right)
\end{cases}
\]
 Here, the angular velocity $\Omega$ of the Casimir cavity is treated
as a constant. Applying the coordinate transformation 
\[
\begin{cases}
d\phi'=\frac{1}{r_{0}}dx\\
d\theta'=\frac{1}{r_{0}}dz\\
dr'=dy\\
dt'=dt
\end{cases}
\]
 based at the point $\boldsymbol{x}_{0}=\left(r_{0},\theta_{0},\phi_{0}\right)$
gives 
\begin{align*}
ds^{2} & =g^{'}_{\mu\nu}dx^{'\mu}dx^{'\nu}\\
 & =\left(g_{tt}+2\Omega g_{t\phi}+\Omega^{2}g_{\phi\phi}\right)dt'^{2}+2\left(g_{t\phi}+\Omega g_{\phi\phi}\right)dt'd\phi^{'}+g_{rr}dr'^{2}+g_{\theta\theta}d\theta'^{2}+g_{\phi\phi}d\phi^{'2}\\
 & =\left(g_{tt}+2\Omega g_{t\phi}+\Omega^{2}g_{\phi\phi}\right)dt^{2}+2\left(g_{t\phi}+\Omega g_{\phi\phi}\right)dt\left(\frac{1}{r_{0}}dx\right)+g_{rr}dy^{2}+g_{\theta\theta}\left(\frac{1}{r_{0}}dz\right)^{2}+g_{\phi\phi}\left(\frac{1}{r_{0}}dx\right)^{2}\\
 & =\left(g_{tt}+2\Omega g_{t\phi}+\Omega^{2}g_{\phi\phi}\right)dt^{2}+2\left(-\frac{A}{\Sigma}\text{sin}^{2}\theta\right)\left(-\omega_{d}+\Omega\right)dt\left(\frac{1}{r_{0}}dx\right)-\frac{\Sigma}{\Delta}dy^{2}-\Sigma\left(\frac{1}{r_{0}}dz\right)^{2}-\frac{A}{\Sigma}\text{sin}^{2}\theta\left(\frac{1}{r_{0}}dx\right)^{2}\\
 & =\left(g_{tt}+2\Omega g_{t\phi}+\Omega^{2}g_{\phi\phi}\right)dt^{2}-2\frac{A}{r_{0}\Sigma}\text{sin}^{2}\theta\left(\Omega-\omega_{d}\right)dtdx-\frac{\Sigma}{\Delta}dy^{2}-\frac{\Sigma}{r^{2}_{0}}dz^{2}-\frac{A}{r^{2}_{0}\Sigma}\text{sin}^{2}\theta dx^{2}\\
 & =C^{-2}\left(\Omega;r,\theta\right)dt^{2}-2\frac{A}{r_{0}\Sigma}\text{sin}^{2}\theta\left(\Omega-\omega_{d}\right)dtdx-\frac{A}{r^{2}_{0}\Sigma}\text{sin}^{2}\theta dx^{2}-\frac{\Sigma}{\Delta}dy^{2}-\frac{\Sigma}{r^{2}_{0}}dz^{2}\\
 & =\hat{g}_{tt}dt^{2}+2\hat{g}_{tx}dtdx+\hat{g}_{xx}dx^{2}+\hat{g}_{yy}dy^{2}+\hat{g}_{zz}dz^{2}
\end{align*}
 where 
\[
C\left(\Omega;r,\theta\right)=\left\{ \frac{\Delta\Sigma}{A}\left[1-\frac{A^{2}\text{sin}^{2}\theta}{\Delta\Sigma^{2}}\left(\Omega-\omega_{d}\right)^{2}\right]\right\} ^{-1/2}
\]
 The coordinate $r$ is a function of $x,y,z$.

\begin{align*}
\hat{g}_{\mu\nu} & =\begin{pmatrix}C^{-2}\left(\Omega;r,\theta\right) & -\frac{A}{r_{0}\Sigma}\text{sin}^{2}\theta\left(\Omega-\omega_{d}\right) & 0 & 0\\
-\frac{A}{r_{0}\Sigma}\text{sin}^{2}\theta\left(\Omega-\omega_{d}\right) & -\frac{A}{r^{2}_{0}\Sigma}\text{sin}^{2}\theta & 0 & 0\\
0 & 0 & -\frac{\Sigma}{\Delta} & 0\\
0 & 0 & 0 & -\frac{\Sigma}{r^{2}_{0}}
\end{pmatrix}\\
 & =\begin{pmatrix}\hat{g}_{tt} & \hat{g}_{tx} & 0 & 0\\
\hat{g}_{xt} & \hat{g}_{xx} & 0 & 0\\
0 & 0 & \hat{g}_{yy} & 0\\
0 & 0 & 0 & \hat{g}_{zz}
\end{pmatrix}
\end{align*}

Here, $\hat{g}_{tx}=\hat{g}_{xt}$. 
\begin{align*}
 & \sum^{3}_{i=1}\frac{\partial\hat{g}_{\mu\nu}\biggl(t,x,y,z\biggr)}{\partial x^{i}}\Bigg|_{\boldsymbol{x}=\boldsymbol{x}_{0}}\cdot\delta x^{i}\\
 & =\begin{pmatrix}\begin{array}{c}
\frac{\partial\hat{g}_{tt}}{\partial x}\big|_{\boldsymbol{x}=\boldsymbol{x}_{0}}\delta x\\
+\frac{\partial\hat{g}_{tt}}{\partial y}\big|_{\boldsymbol{x}=\boldsymbol{x}_{0}}\delta y\\
+\frac{\partial\hat{g}_{tt}}{\partial z}\big|_{\boldsymbol{x}=\boldsymbol{x}_{0}}\delta z
\end{array} & \begin{array}{c}
\frac{\partial\hat{g}_{tx}}{\partial x}\big|_{\boldsymbol{x}=\boldsymbol{x}_{0}}\delta x\\
+\frac{\partial\hat{g}_{tx}}{\partial y}\big|_{\boldsymbol{x}=\boldsymbol{x}_{0}}\delta y\\
+\frac{\partial\hat{g}_{tx}}{\partial z}\big|_{\boldsymbol{x}=\boldsymbol{x}_{0}}\delta z
\end{array} & 0 & 0\\
\begin{array}{c}
\frac{\partial\hat{g}_{xt}}{\partial x}\big|_{\boldsymbol{x}=\boldsymbol{x}_{0}}\delta x\\
+\frac{\partial\hat{g}_{xt}}{\partial y}\big|_{\boldsymbol{x}=\boldsymbol{x}_{0}}\delta y\\
+\frac{\partial\hat{g}_{xt}}{\partial z}\big|_{\boldsymbol{x}=\boldsymbol{x}_{0}}\delta z
\end{array} & \begin{array}{c}
\frac{\partial\hat{g}_{xx}}{\partial x}\big|_{\boldsymbol{x}=\boldsymbol{x}_{0}}\delta x\\
+\frac{\partial\hat{g}_{xx}}{\partial y}\big|_{\boldsymbol{x}=\boldsymbol{x}_{0}}\delta y\\
+\frac{\partial\hat{g}_{xx}}{\partial z}\big|_{\boldsymbol{x}=\boldsymbol{x}_{0}}\delta z
\end{array} & 0 & 0\\
0 & 0 & \begin{array}{c}
\frac{\partial\hat{g}_{yy}}{\partial x}\big|_{\boldsymbol{x}=\boldsymbol{x}_{0}}\delta x\\
+\frac{\partial\hat{g}_{yy}}{\partial y}\big|_{\boldsymbol{x}=\boldsymbol{x}_{0}}\delta y\\
+\frac{\partial\hat{g}_{yy}}{\partial z}\big|_{\boldsymbol{x}=\boldsymbol{x}_{0}}\delta z
\end{array} & 0\\
0 & 0 & 0 & \begin{array}{c}
\frac{\partial\hat{g}_{zz}}{\partial x}\big|_{\boldsymbol{x}=\boldsymbol{x}_{0}}\delta x\\
+\frac{\partial\hat{g}_{zz}}{\partial y}\big|_{\boldsymbol{x}=\boldsymbol{x}_{0}}\delta y\\
+\frac{\partial\hat{g}_{zz}}{\partial z}\big|_{\boldsymbol{x}=\boldsymbol{x}_{0}}\delta z
\end{array}
\end{pmatrix}\\
 & =\begin{pmatrix}\begin{array}{c}
\frac{\partial C^{-2}\biggl(\Omega;r,\theta\biggr)}{\partial\phi^{'}}\\
\qquad\frac{\partial\phi^{'}}{\partial x}\big|_{\boldsymbol{x}=\boldsymbol{x}_{0}}\delta x\\
+\frac{\partial C^{-2}\biggl(\Omega;r,\theta\biggr)}{\partial r^{'}}\\
\qquad\frac{\partial r^{'}}{\partial y}\big|_{\boldsymbol{x}=\boldsymbol{x}_{0}}\delta y\\
+\frac{\partial C^{-2}\biggl(\Omega;r,\theta\biggr)}{\partial\theta^{'}}\\
\qquad\frac{\partial\theta^{'}}{\partial z}\big|_{\boldsymbol{x}=\boldsymbol{x}_{0}}\delta z
\end{array} & \begin{array}{c}
\frac{\partial\biggl[-\frac{A}{r_{0}\Sigma}\text{sin}^{2}\theta\biggl(\Omega-\omega_{d}\biggr)\biggr]}{\partial\phi^{'}}\\
\qquad\frac{\partial\phi^{'}}{\partial x}\big|_{\boldsymbol{x}=\boldsymbol{x}_{0}}\delta x\\
+\frac{\partial\biggl[-\frac{A}{r_{0}\Sigma}\text{sin}^{2}\theta\biggl(\Omega-\omega_{d}\biggr)\biggr]}{\partial r^{'}}\\
\qquad\frac{\partial r^{'}}{\partial y}\big|_{\boldsymbol{x}=\boldsymbol{x}_{0}}\delta y\\
+\frac{\partial\biggl[-\frac{A}{r_{0}\Sigma}\text{sin}^{2}\theta\biggl(\Omega-\omega_{d}\biggr)\biggr]}{\partial\theta^{'}}\\
\qquad\frac{\partial\theta^{'}}{\partial z}\big|_{\boldsymbol{x}=\boldsymbol{x}_{0}}\delta z
\end{array} & 0 & 0\\
\begin{array}{c}
\frac{\partial\biggl[-\frac{A}{r_{0}\Sigma}\text{sin}^{2}\theta\biggl(\Omega-\omega_{d}\biggr)\biggr]}{\partial\phi^{'}}\\
\qquad\frac{\partial\phi^{'}}{\partial x}\big|_{\boldsymbol{x}=\boldsymbol{x}_{0}}\delta x\\
+\frac{\partial\biggl[-\frac{A}{r_{0}\Sigma}\text{sin}^{2}\theta\biggl(\Omega-\omega_{d}\biggr)\biggr]}{\partial r^{'}}\\
\qquad\frac{\partial r^{'}}{\partial y}\big|_{\boldsymbol{x}=\boldsymbol{x}_{0}}\delta y\\
+\frac{\partial\biggl[-\frac{A}{r_{0}\Sigma}\text{sin}^{2}\theta\biggl(\Omega-\omega_{d}\biggr)\biggr]}{\partial\theta^{'}}\\
\qquad\frac{\partial\theta^{'}}{\partial z}\big|_{\boldsymbol{x}=\boldsymbol{x}_{0}}\delta z
\end{array} & \begin{array}{c}
\frac{\partial\biggl[-\frac{A}{r^{2}_{0}\Sigma}\text{sin}^{2}\theta\biggr]}{\partial\phi^{'}}\\
\qquad\frac{\partial\phi^{'}}{\partial x}\big|_{\boldsymbol{x}=\boldsymbol{x}_{0}}\delta x\\
+\frac{\partial\biggl[-\frac{A}{r^{2}_{0}\Sigma}\text{sin}^{2}\theta\biggr]}{\partial r^{'}}\\
\qquad\frac{\partial r^{'}}{\partial y}\big|_{\boldsymbol{x}=\boldsymbol{x}_{0}}\delta y\\
+\frac{\partial\biggl[-\frac{A}{r^{2}_{0}\Sigma}\text{sin}^{2}\theta\biggr]}{\partial\theta^{'}}\\
\qquad\frac{\partial\theta^{'}}{\partial z}\big|_{\boldsymbol{x}=\boldsymbol{x}_{0}}\delta z
\end{array} & 0 & 0\\
0 & 0 & \begin{array}{c}
\frac{\partial\biggl[-\frac{\Sigma}{\Delta}\biggr]}{\partial\phi^{'}}\\
\qquad\frac{\partial\phi^{'}}{\partial x}\big|_{\boldsymbol{x}=\boldsymbol{x}_{0}}\delta x\\
+\frac{\partial\biggl[-\frac{\Sigma}{\Delta}\biggr]}{\partial r^{'}}\\
\qquad\frac{\partial r^{'}}{\partial y}\big|_{\boldsymbol{x}=\boldsymbol{x}_{0}}\delta y\\
+\frac{\partial\biggl[-\frac{\Sigma}{\Delta}\biggr]}{\partial\theta^{'}}\\
\qquad\frac{\partial\theta^{'}}{\partial z}\big|_{\boldsymbol{x}=\boldsymbol{x}_{0}}\delta z
\end{array} & 0\\
0 & 0 & 0 & \begin{array}{c}
\frac{\partial\biggl[-\frac{\Sigma}{r^{2}_{0}}\biggr]}{\partial\phi^{'}}\\
\qquad\frac{\partial\phi^{'}}{\partial x}\big|_{\boldsymbol{x}=\boldsymbol{x}_{0}}\delta x\\
+\frac{\partial\biggl[-\frac{\Sigma}{r^{2}_{0}}\biggr]}{\partial r^{'}}\\
\qquad\frac{\partial r^{'}}{\partial y}\big|_{\boldsymbol{x}=\boldsymbol{x}_{0}}\delta y\\
+\frac{\partial\biggl[-\frac{\Sigma}{r^{2}_{0}}\biggr]}{\partial\theta^{'}}\\
\qquad\frac{\partial\theta^{'}}{\partial z}\big|_{\boldsymbol{x}=\boldsymbol{x}_{0}}\delta z
\end{array}
\end{pmatrix}
\end{align*}

Here, 
\begin{align*}
\frac{\partial f\left(t,r,\theta,\phi\right)}{\partial\phi^{'}} & =\frac{\partial f\left(t,r,\theta,\phi'+\Omega t\right)}{\partial\phi^{'}}\\
 & =\frac{\partial f\left(t,r,\theta,\phi^{'}\left(\phi,t\right)+\Omega t\right)}{\partial\phi^{'}}\\
 & =\frac{\partial f}{\partial\phi}\frac{\partial\phi}{\partial\phi^{'}}+\frac{\partial f}{\partial t}\frac{\partial t}{\partial\phi^{'}}\hfill\left(\text{chain rule}\right)\\
 & =\frac{\partial f}{\partial\phi}\frac{\partial\left(\phi'+\Omega t\right)}{\partial\phi^{'}}+\frac{\partial f}{\partial t}\frac{\partial t}{\partial\phi^{'}}\\
 & =\frac{\partial f}{\partial\phi}\hfill\left(\text{t is an independent variable, \ensuremath{\frac{\partial t}{\partial\phi^{'}}=0}}\right)
\end{align*}

(An additional, redundant intermediate step would give 
\begin{align*}
\frac{\partial f\left(t,r,\theta,\phi\right)}{\partial\left(\phi-\Omega t\right)}\frac{\partial\left(\phi-\Omega t\right)}{\partial\phi^{'}}\frac{\partial\phi^{'}}{\partial x} & =\frac{\partial f\left(t,r,\theta,\psi+\Omega t\right)}{\partial\psi}\frac{\partial\psi}{\partial\left(\phi-\Omega t\right)}\frac{\partial\left(\phi-\Omega t\right)}{\partial\phi^{'}}\frac{\partial\phi^{'}}{\partial x}\\
 & =\frac{\partial f\left(t,r,\theta,\psi\left(\phi,t\right)+\Omega t\right)}{\partial\psi}\frac{\partial\left(\phi-\Omega t\right)}{\partial\left(\phi-\Omega t\right)}\frac{\partial\left(\phi-\Omega t\right)}{\partial\left(\phi-\Omega t\right)}\frac{\partial\phi^{'}}{\partial x}\\
 & =\left(\frac{\partial f}{\partial\phi}\frac{\partial\phi}{\partial\psi}+\frac{\partial f}{\partial t}\frac{\partial t}{\partial\psi}\right)\frac{\partial\phi^{'}}{\partial x}\\
 & =\left(\frac{\partial f}{\partial\phi}\frac{\partial\left(\psi+\Omega t\right)}{\partial\psi}+\frac{\partial f}{\partial t}\frac{\partial t}{\partial\psi}\right)\frac{\partial\phi^{'}}{\partial x}\\
 & =\frac{\partial f}{\partial\phi}\frac{\partial\phi^{'}}{\partial x}
\end{align*}
)

We obtain

\begin{align*}
\frac{\partial C^{-2}\left(\Omega;r,\theta\right)}{\partial\phi^{'}}\frac{\partial\phi^{'}}{\partial x} & =\frac{\partial\left\{ \frac{\Delta\Sigma}{A}\left[1-\frac{A^{2}\text{sin}^{2}\theta}{\Delta\Sigma^{2}}\left(\Omega-\omega_{d}\right)^{2}\right]\right\} }{\partial\phi^{'}}\cdot\frac{1}{r_{0}}\\
 & =\frac{\partial\left\{ \frac{\Delta\Sigma}{A}\left[1-\frac{A^{2}\text{sin}^{2}\theta}{\Delta\Sigma^{2}}\left(\Omega-\omega_{d}\right)^{2}\right]\right\} }{\partial\phi}\cdot\frac{1}{r_{0}}=0\cdot\frac{1}{r_{0}}=0
\end{align*}
 In other words, since none of the metric components $\hat{g}_{\mu\nu}$
depends on $\phi$, for all $\mu,\nu$ we have 
\[
\frac{\partial\hat{g}_{\mu\nu}}{\partial x}=\frac{\partial\hat{g}_{\mu\nu}}{\partial\phi^{'}}\frac{\partial\phi^{'}}{\partial x}=\frac{\partial\hat{g}_{\mu\nu}}{\partial\phi}\frac{\partial\phi^{'}}{\partial x}=0
\]
 Introduce the notation 
\[
\mathcal{D}f\equiv\frac{\partial f}{\partial r^{'}}\frac{\partial r^{'}}{\partial y}\delta y+\frac{\partial f}{\partial\theta^{'}}\frac{\partial\theta^{'}}{\partial z}\delta z=\frac{\partial f}{\partial r}\delta y+\frac{1}{r_{0}}\frac{\partial f}{\partial\theta}\delta z
\]
 {}
\[
Q\equiv\frac{A\sin^{2}\theta}{\Sigma}
\]
 The following partial derivatives were obtained using a Python program.
The subscript $0$ denotes evaluation at $\boldsymbol{x}=\boldsymbol{x}_{0}$.
Write $s_{0}=\sin\theta_{0}$, $c_{0}=\cos\theta_{0}$, $\Sigma_{0}=r^{2}_{0}+a^{2}c^{2}_{0}$,
$\Delta_{0}=r^{2}_{0}-2Mr_{0}+a^{2}$, and $U_{0}=1-a\Omega s^{2}_{0}$.

\[
\left.\partial_{x}\hat{g}\right|_{0}=0
\]
\[
R_{0}=\left.\partial_{y}\hat{g}\right|_{0}=\begin{pmatrix}-\frac{2MU^{2}_{0}}{\Sigma_{0}}+\frac{4MU^{2}_{0}r^{2}_{0}}{\Sigma^{2}_{0}}-2\Omega^{2}r_{0}s^{2}_{0} & \frac{2s^{2}_{0}\left(\frac{MU_{0}a}{\Sigma_{0}}-\frac{2MU_{0}ar^{2}_{0}}{\Sigma^{2}_{0}}-\Omega r_{0}\right)}{r_{0}} & 0 & 0\\
\frac{2s^{2}_{0}\left(\frac{MU_{0}a}{\Sigma_{0}}-\frac{2MU_{0}ar^{2}_{0}}{\Sigma^{2}_{0}}-\Omega r_{0}\right)}{r_{0}} & -\frac{2s^{2}_{0}\left(\frac{Ma^{2}s^{2}_{0}}{\Sigma_{0}}-\frac{2Ma^{2}r^{2}_{0}s^{2}_{0}}{\Sigma^{2}_{0}}+r_{0}\right)}{r^{2}_{0}} & 0 & 0\\
0 & 0 & -\frac{2\left(r_{0}+\frac{\Sigma_{0}\left(M-r_{0}\right)}{\Delta_{0}}\right)}{\Delta_{0}} & 0\\
0 & 0 & 0 & -\frac{2}{r_{0}}
\end{pmatrix}
\]

\begin{align*}
 & P_{0}=\biggl.\partial_{z}\hat{g}\biggr|_{0}=\\
 & \begin{pmatrix}\frac{\begin{gathered}2c_{0}s_{0}\biggl(\frac{4M\Omega U_{0}ar_{0}}{\Sigma_{0}}\\
-\frac{2MU^{2}_{0}a^{2}r_{0}}{\Sigma^{2}_{0}}-\Omega^{2}\\
\biggl(a^{2}+r^{2}_{0}\biggr)\biggr)
\end{gathered}
}{r_{0}} & \frac{\begin{gathered}2c_{0}s_{0}\biggl(\frac{2MU_{0}ar_{0}}{\Sigma_{0}}\\
+\frac{2Ma^{2}r_{0}s^{2}_{0}\biggl(-\Omega+\frac{U_{0}a}{\Sigma_{0}}\biggr)}{\Sigma_{0}}\\
-\Omega\biggl(a^{2}+r^{2}_{0}\biggr)\biggr)
\end{gathered}
}{r^{2}_{0}} & 0 & 0\\
\frac{\begin{gathered}2c_{0}s_{0}\biggl(\frac{2MU_{0}ar_{0}}{\Sigma_{0}}\\
+\frac{2Ma^{2}r_{0}s^{2}_{0}\biggl(-\Omega+\frac{U_{0}a}{\Sigma_{0}}\biggr)}{\Sigma_{0}}\\
-\Omega\biggl(a^{2}+r^{2}_{0}\biggr)\biggr)
\end{gathered}
}{r^{2}_{0}} & -\frac{\begin{gathered}2c_{0}s_{0}\\
\biggl(\frac{2Ma^{2}r_{0}s^{2}_{0}\biggl(1+\frac{a^{2}s^{2}_{0}}{\Sigma_{0}}\biggr)}{\Sigma_{0}}\\
+\frac{2Ma^{2}r_{0}s^{2}_{0}}{\Sigma_{0}}+a^{2}\\
+r^{2}_{0}\biggr)
\end{gathered}
}{r^{3}_{0}} & 0 & 0\\
0 & 0 & \frac{2a^{2}c_{0}s_{0}}{\Delta_{0}r_{0}} & 0\\
0 & 0 & 0 & \frac{2a^{2}c_{0}s_{0}}{r^{3}_{0}}
\end{pmatrix}
\end{align*}

\begin{align*}
 & \hat{M}=R_{0}\,\delta y+P_{0}\,\delta z=\\
 & \begin{pmatrix}\begin{gathered}\frac{\begin{gathered}2c_{0}\delta zs_{0}\biggl(\frac{4M\Omega U_{0}ar_{0}}{\Sigma_{0}}\\
-\frac{2MU^{2}_{0}a^{2}r_{0}}{\Sigma^{2}_{0}}-\Omega^{2}\biggl(a^{2}\\
+r^{2}_{0}\biggr)\biggr)
\end{gathered}
}{r_{0}}\\
+2\delta y\biggl(-\frac{MU^{2}_{0}}{\Sigma_{0}}+\frac{2MU^{2}_{0}r^{2}_{0}}{\Sigma^{2}_{0}}\\
-\Omega^{2}r_{0}s^{2}_{0}\biggr)
\end{gathered}
 & \begin{gathered}\frac{\begin{gathered}2c_{0}\delta zs_{0}\biggl(\frac{2MU_{0}ar_{0}}{\Sigma_{0}}\\
+\frac{2Ma^{2}r_{0}s^{2}_{0}\biggl(-\Omega+\frac{U_{0}a}{\Sigma_{0}}\biggr)}{\Sigma_{0}}\\
-\Omega\biggl(a^{2}+r^{2}_{0}\biggr)\biggr)
\end{gathered}
}{r^{2}_{0}}\\
+\frac{\begin{gathered}2\delta ys^{2}_{0}\biggl(\frac{MU_{0}a}{\Sigma_{0}}\\
-\frac{2MU_{0}ar^{2}_{0}}{\Sigma^{2}_{0}}-\Omega r_{0}\biggr)
\end{gathered}
}{r_{0}}
\end{gathered}
 & 0 & 0\\
\begin{gathered}\frac{\begin{gathered}2c_{0}\delta zs_{0}\biggl(\frac{2MU_{0}ar_{0}}{\Sigma_{0}}\\
+\frac{2Ma^{2}r_{0}s^{2}_{0}\biggl(-\Omega+\frac{U_{0}a}{\Sigma_{0}}\biggr)}{\Sigma_{0}}\\
-\Omega\biggl(a^{2}+r^{2}_{0}\biggr)\biggr)
\end{gathered}
}{r^{2}_{0}}\\
+\frac{\begin{gathered}2\delta ys^{2}_{0}\biggl(\frac{MU_{0}a}{\Sigma_{0}}\\
-\frac{2MU_{0}ar^{2}_{0}}{\Sigma^{2}_{0}}-\Omega r_{0}\biggr)
\end{gathered}
}{r_{0}}
\end{gathered}
 & \begin{gathered}-\frac{\begin{gathered}2c_{0}\delta zs_{0}\\
\biggl(\frac{2Ma^{2}r_{0}s^{2}_{0}\biggl(1+\frac{a^{2}s^{2}_{0}}{\Sigma_{0}}\biggr)}{\Sigma_{0}}\\
+\frac{2Ma^{2}r_{0}s^{2}_{0}}{\Sigma_{0}}+a^{2}+r^{2}_{0}\\
\biggr)
\end{gathered}
}{r^{3}_{0}}\\
-\frac{\begin{gathered}2\delta ys^{2}_{0}\biggl(\frac{Ma^{2}s^{2}_{0}}{\Sigma_{0}}\\
-\frac{2Ma^{2}r^{2}_{0}s^{2}_{0}}{\Sigma^{2}_{0}}+r_{0}\biggr)
\end{gathered}
}{r^{2}_{0}}
\end{gathered}
 & 0 & 0\\
0 & 0 & \begin{gathered}\frac{2a^{2}c_{0}\delta zs_{0}}{\Delta_{0}r_{0}}\\
-\frac{2\delta y\biggl(r_{0}+\frac{\Sigma_{0}\biggl(M-r_{0}\biggr)}{\Delta_{0}}\biggr)}{\Delta_{0}}
\end{gathered}
 & 0\\
0 & 0 & 0 & \frac{2a^{2}c_{0}\delta zs_{0}}{r^{3}_{0}}-\frac{2\delta y}{r_{0}}
\end{pmatrix}.
\end{align*}

When $\theta_{0}=\pi/2$, 
\[
\left.\hat{M}\right|_{\theta_{0}=\pi/2}=\delta y\,\begin{pmatrix}2\left(\frac{M\left(-\Omega a+1\right)^{2}}{r^{2}_{0}}-\Omega^{2}r_{0}\right) & -\frac{2\left(\frac{Ma\left(-\Omega a+1\right)}{r^{2}_{0}}+\Omega r_{0}\right)}{r_{0}} & 0 & 0\\
-\frac{2\left(\frac{Ma\left(-\Omega a+1\right)}{r^{2}_{0}}+\Omega r_{0}\right)}{r_{0}} & -\frac{2\left(-\frac{Ma^{2}}{r^{2}_{0}}+r_{0}\right)}{r^{2}_{0}} & 0 & 0\\
0 & 0 & -\frac{2r_{0}\left(1+\frac{r_{0}\left(M-r_{0}\right)}{\Delta_{0}}\right)}{\Delta_{0}} & 0\\
0 & 0 & 0 & -\frac{2}{r_{0}}
\end{pmatrix}
\]

(This is because, when $\theta_{0}=\pi/2$, the factor $c_{0}=\cos\theta_{0}$
appearing in every term of $P_{0}=\left.\partial_{z}\hat{g}\right|_{0}$
is $0$.) The center of the cavity has fixed spatial coordinates in
the frame at rest relative to the cavity. Thus, along its central
trajectory, 
\[
dx=dy=dz=0
\]
 and 
\[
ds^{2}=\left.C^{-2}\left(\Omega;r_{0},\theta_{0}\right)\right|_{\theta_{0}=\pi/2}dt^{2}
\]
 {}
\begin{align*}
C^{-2}\left(\Omega;r,\theta\right) & =g_{tt}+2\Omega g_{t\phi}+\Omega^{2}g_{\phi\phi}\\
 & =\left(1-\frac{2Mr}{\Sigma}\right)+2\Omega\frac{2Mar}{\Sigma}\text{sin}^{2}\theta+\Omega^{2}\left[-\left(r^{2}+a^{2}+\frac{2Mra^{2}\text{sin}^{2}\theta}{\Sigma}\right)\text{sin}^{2}\theta\right]\\
 & =1-\frac{2Mr}{\Sigma}+\frac{2Mr}{\Sigma}\left(2a\Omega\text{sin}^{2}\theta\right)-\Omega^{2}\left(r^{2}+a^{2}\right)\text{sin}^{2}\theta-\frac{2Mr}{\Sigma}\left(\Omega^{2}a^{2}\text{sin}^{4}\theta\right)\\
 & =1-\Omega^{2}\left(r^{2}+a^{2}\right)\text{sin}^{2}\theta-\frac{2Mr}{\Sigma}\left(1-2a\Omega\text{sin}^{2}\theta+\Omega^{2}a^{2}\text{sin}^{4}\theta\right)
\end{align*}
 {}
\begin{align*}
g_{\phi\phi}=-\frac{A}{\Sigma}\text{sin}^{2}\theta & =-\frac{\left(r^{2}+a^{2}\right)\Sigma+2Mra^{2}\text{sin}^{2}\theta}{\Sigma}\text{sin}^{2}\theta\\
 & =-\left(r^{2}+a^{2}+\frac{2Mra^{2}\text{sin}^{2}\theta}{\Sigma}\right)\text{sin}^{2}\theta
\end{align*}
 where we have used 
\[
A=\left(r^{2}+a^{2}\right)\Sigma+2Mra^{2}\text{sin}^{2}\theta
\]
 Taking the cavity center to be at $r_{0},\theta_{0}=\pi/2$, we obtain
\begin{align*}
\left.C^{-2}\left(\Omega;r_{0},\theta_{0}\right)\right|_{\theta_{0}=\pi/2} & =1-\Omega^{2}\left(r^{2}_{0}+a^{2}\right)-\frac{2Mr_{0}}{\Sigma|_{\theta_{0}=\pi/2}}\left(1-2a\Omega+\Omega^{2}a^{2}\right)\\
 & =1-\Omega^{2}\left(r^{2}_{0}+a^{2}\right)-\frac{2M}{r_{0}}\left(1-a\Omega\right)^{2}
\end{align*}
 where 
\[
\Sigma|_{\theta_{0}=\pi/2}=r^{2}_{0}+a^{2}\text{cos}^{2}\left(\frac{\pi}{2}\right)=r^{2}_{0}
\]
 Define 
\[
u=\frac{M}{r_{0}},j=\frac{a}{r_{0}},v=\Omega r_{0},d_{0}=\frac{\Delta_{0}}{r^{2}_{0}}
\]
 where 
\[
\Delta_{0}=r^{2}_{0}-2Mr_{0}+a^{2}
\]
 Then 
\begin{align*}
\left.C^{-2}\left(\Omega;r_{0},\theta_{0}\right)\right|_{\theta_{0}=\pi/2} & =1-\Omega^{2}\left(r^{2}_{0}+a^{2}\right)-\frac{2M}{r_{0}}\left(1-a\Omega\right)^{2}\\
 & =1-\left(1+\frac{a^{2}}{r^{2}_{0}}\right)\Omega^{2}r^{2}_{0}-2\frac{M}{r_{0}}\left(1-\frac{a}{r_{0}}\cdot\Omega r_{0}\right)^{2}\\
 & =1-\left(1+j^{2}\right)v^{2}-2u\left(1-jv\right)^{2}
\end{align*}
 A massive observer must follow a timelike trajectory, so 
\[
ds^{2}=\left.C^{-2}\left(\Omega;r_{0},\theta_{0}\right)\right|_{\theta_{0}=\pi/2}dt^{2}>0
\]
 Hence, 
\[
\left.C^{-2}\left(\Omega;r_{0},\theta_{0}\right)\right|_{\theta_{0}=\pi/2}=1-\left(1+j^{2}\right)v^{2}-2u\left(1-jv\right)^{2}>0
\]
 which gives 
\[
\left(1+j^{2}\right)v^{2}+2u\left(1-jv\right)^{2}<1
\]
 Since $u>0,\left(1-jv\right)^{2}\geq0$, we obtain 
\[
\left(1+j^{2}\right)v^{2}\leq\left(1+j^{2}\right)v^{2}+2u\left(1-jv\right)^{2}<1
\]
 Using $1+j^{2}\geq1$ also gives 
\[
v^{2}<\frac{1}{1+j^{2}}\leq1
\]
 and hence 
\[
|v|<\frac{1}{\sqrt{1+j^{2}}}\leq1
\]
 Thus, 
\[
|v|\leq1,
\]
 or equivalently, 
\[
|\Omega r_{0}|\leq1.
\]
 From 
\[
r_{0}>r_{+}\geq M>0
\]
 we have 
\[
0<\frac{M}{r_{0}}<1
\]
 From 
\[
|a|\leq M,r_{0}>0
\]
 we obtain 
\[
\frac{|a|}{r_{0}}\leq\frac{M}{r_{0}}<1.
\]
 Using 
\begin{align*}
d_{0} & =\frac{\Delta_{0}}{r^{2}_{0}}=\frac{r^{2}_{0}-2Mr_{0}+a^{2}}{r^{2}_{0}}\\
 & =1-\frac{2M}{r_{0}}+\frac{a^{2}}{r^{2}_{0}}
\end{align*}
 and 
\[
\frac{a^{2}}{r^{2}_{0}}\leq\frac{M^{2}}{r^{2}_{0}}<\frac{M}{r_{0}}
\]
 and recalling that $0<M/r_{0}<1$, we obtain 
\[
d_{0}=1-\frac{2M}{r_{0}}+\frac{a^{2}}{r^{2}_{0}}<1-\frac{M}{r_{0}}<1
\]
 Outside the outer event horizon, 
\[
\Delta_{0}>0
\]
 Therefore, 
\[
0<d_{0}<1
\]
 (The outer event horizon of a Kerr black hole is $r_{+}=M+\sqrt{M^{2}-a^{2}}$.
Requiring the radicand to be nonnegative gives the maximum magnitude
of the specific angular momentum $a=J/M$ as $|a|_{\max}=M$, in which
case the outer event horizon is $r_{+}=M$. When $a=0$, we have $r_{+}=2M$.)

Let $L_{\mathcal{C}}$ denote the maximum extent of the entire finite
parallel-plate cavity $\mathcal{C}$, including both the plate separation
and the finite transverse dimensions. Define 
\[
L_{\mathcal{C}}=\sup_{\boldsymbol{x}\in\mathcal{C}}\max_{i}|\delta x^{i}|,
\]
 and 
\[
\varepsilon=\frac{L_{\mathcal{C}}}{r_{0}}.
\]
 Writing $U_{0}=1-a\Omega$, we have 
\[
\left.\hat{M}\right|_{\theta_{0}=\pi/2}=\delta y\,R_{0}=\delta y\,\begin{pmatrix}2\frac{MU^{2}_{0}}{r^{2}_{0}}-2\Omega^{2}r_{0} & -2\frac{MaU_{0}}{r^{3}_{0}}-2\Omega & 0 & 0\\
-2\frac{MaU_{0}}{r^{3}_{0}}-2\Omega & 2\frac{Ma^{2}}{r^{4}_{0}}-2\frac{1}{r_{0}} & 0 & 0\\
0 & 0 & \frac{2r_{0}\left(Mr_{0}-a^{2}\right)}{\Delta^{2}_{0}} & 0\\
0 & 0 & 0 & -\frac{2}{r_{0}}
\end{pmatrix}
\]
 where 
\begin{align*}
-\frac{2r_{0}\left(1+\frac{r_{0}\left(M-r_{0}\right)}{\Delta_{0}}\right)}{\Delta_{0}} & =-\frac{2r_{0}\left[\Delta_{0}+r_{0}\left(M-r_{0}\right)\right]}{\Delta^{2}_{0}}\\
 & =-\frac{2r_{0}\left[\left(r^{2}_{0}-2Mr_{0}+a^{2}\right)+r_{0}\left(M-r_{0}\right)\right]}{\Delta^{2}_{0}}\\
 & =\frac{2r_{0}\left(Mr_{0}-a^{2}\right)}{\Delta^{2}_{0}}
\end{align*}
 {}
\[
|U_{0}|=|1-a\Omega|\leq1+|a\Omega|=1+\frac{|a|}{r_{0}}|\Omega|r_{0}<2
\]
 (Here we have used the triangle inequality $|x+y|\leq|x|+|y|$ and
the previously established bounds $|\Omega r_{0}|\leq1$ and $|a|/r_{0}<1$.)
Under the stated convention, $|U_{0}|<2$ also implies $|U_{0}|\leq2$.
Thus, 
\[
r_{0}\left|\left(R_{0}\right)_{00}\right|=\left|2\frac{MU^{2}_{0}}{r_{0}}-2\Omega^{2}r^{2}_{0}\right|\leq\left|2\frac{M}{r_{0}}U^{2}_{0}\right|+\left|2\Omega^{2}r^{2}_{0}\right|<8+2=10,
\]
 {}
\[
r_{0}\left|\left(R_{0}\right)_{01}\right|=r_{0}\left|\left(R_{0}\right)_{10}\right|=\left|-2\frac{MaU_{0}}{r^{2}_{0}}-2\Omega r_{0}\right|\leq\left|2\frac{M}{r_{0}}\frac{a}{r_{0}}U_{0}\right|+\left|2\Omega r_{0}\right|<4+2=6,
\]
 Using 
\[
0\leq\frac{M}{r_{0}}\frac{a^{2}}{r^{2}_{0}}<1\Rightarrow-2\leq2\frac{M}{r_{0}}\frac{a^{2}}{r^{2}_{0}}-2<0\Rightarrow2\geq2-2\frac{M}{r_{0}}\frac{a^{2}}{r^{2}_{0}}=\left|2\frac{Ma^{2}}{r^{3}_{0}}-2\right|>0
\]
 gives 
\[
r_{0}\left|\left(R_{0}\right)_{11}\right|=\left|2\frac{Ma^{2}}{r^{3}_{0}}-2\right|\leq2,
\]
 {}
\[
r_{0}\left|\left(R_{0}\right)_{33}\right|=2,
\]
 We obtain 
\[
\left|\hat{M}_{00}\right|=\frac{r_{0}\left|\left(R_{0}\right)_{00}\right|}{r_{0}}\,\left|\delta y\right|<\frac{10}{r_{0}}\cdot L_{\mathcal{C}}=10\varepsilon\leq10\varepsilon
\]
 (Under the stated convention, $\left|\hat{M}_{00}\right|<10\varepsilon$
naturally implies $\left|\hat{M}_{00}\right|\leq10\varepsilon$.)
Similarly, 
\[
\left|\hat{M}_{01}\right|=\left|\hat{M}_{10}\right|\leq6\varepsilon,
\]
 {}
\[
\left|\hat{M}_{11}\right|\leq2\varepsilon,
\]
 {}
\[
\left|\hat{M}_{33}\right|=\frac{r_{0}\left|\left(R_{0}\right)_{33}\right|}{r_{0}}\,\left|\delta y\right|\leq\frac{r_{0}\left|\left(R_{0}\right)_{33}\right|}{r_{0}}\cdot L_{\mathcal{C}}=2\varepsilon.
\]
 while 
\[
\left|\hat{M}_{22}\right|=\left|\left(R_{0}\right)_{22}\right|\,\left|\delta y\right|\leq\frac{2r_{0}\left(Mr_{0}-a^{2}\right)}{\Delta^{2}_{0}}\cdot L_{\mathcal{C}}\leq\frac{2Mr^{2}_{0}}{\Delta^{2}_{0}}L_{\mathcal{C}}=\frac{2M}{r_{0}}\frac{L_{\mathcal{C}}/r_{0}}{\left(\Delta_{0}/r^{2}_{0}\right)^{2}}\leq\frac{2\varepsilon}{d^{2}_{0}}.
\]
 Since $0<d_{0}<1$, we have $1/d^{2}_{0}>1$, and hence 
\[
\left|\hat{M}_{00}\right|\leq10\varepsilon<\frac{10\varepsilon}{d^{2}_{0}}\leq\frac{10\varepsilon}{d^{2}_{0}},
\]
 {}
\[
\left|\hat{M}_{01}\right|=\left|\hat{M}_{10}\right|<\frac{6\varepsilon}{d^{2}_{0}}\leq\frac{6\varepsilon}{d^{2}_{0}},
\]
 {}
\[
\left|\hat{M}_{11}\right|<\frac{2\varepsilon}{d^{2}_{0}}\leq\frac{2\varepsilon}{d^{2}_{0}},
\]
 {}
\[
\left|\hat{M}_{33}\right|\leq2\varepsilon<\frac{2\varepsilon}{d^{2}_{0}}\leq\frac{2\varepsilon}{d^{2}_{0}}.
\]
 This gives 
\[
\max_{\mu,\nu}\left|\hat{M}_{\mu\nu}\right|\leq\frac{10\varepsilon}{d^{2}_{0}}.
\]
 By Taylor's theorem with the mean-value form of the remainder, $\hat{g}_{\mu\nu}\left(t,x,y,z\right)$
can be written as 
\[
\hat{g}_{\mu\nu}\left(t,x,y,z\right)=\hat{g}_{\mu\nu}\left(t,x,y,z\right)\big|_{\boldsymbol{x}=\boldsymbol{x}_{0}}+\sum^{3}_{i=1}\frac{\partial\hat{g}_{\mu\nu}\left(t,x,y,z\right)}{\partial x^{i}}\Bigg|_{\boldsymbol{x}=\boldsymbol{x}_{0}}\left(x^{i}-x^{i}_{0}\right)+\mathcal{R}_{\mu\nu}
\]
 where 
\[
\mathcal{R}_{\mu\nu}=\frac{1}{2}\sum^{3}_{i,j=1}\frac{\partial^{2}\hat{g}_{\mu\nu}\left(t,x,y,z\right)}{\partial x^{i}\partial x^{j}}\Bigg|_{\boldsymbol{x}=\boldsymbol{\xi}}\left(x^{i}-x^{i}_{0}\right)\left(x^{j}-x^{j}_{0}\right)
\]
 Here, $\boldsymbol{\xi}=\left(\xi_{1},\xi_{2},\xi_{3}\right)$, and
$\xi_{1},\ \xi_{2},\ \xi_{3}$ lie between $x_{0}$ and $x$, $y_{0}$
and $y$, and $z_{0}$ and $z$, respectively (or, equivalently, $\boldsymbol{\xi}=\boldsymbol{x}_{0}+s\left(\boldsymbol{x}-\boldsymbol{x}_{0}\right),\ 0<s<1$).
Let 
\[
\bar{K}(\boldsymbol{\xi})=\frac{1}{2}\sum^{3}_{i,j=1}\max_{\mu,\nu}\left|\frac{\partial^{2}\hat{g}_{\mu\nu}\left(t,x,y,z\right)}{\partial x^{i}\partial x^{j}}\Bigg|_{\boldsymbol{x}=\boldsymbol{\xi}}\right|,
\]
 To specify the range of $\boldsymbol{\xi}$ over which $\bar{K}(\boldsymbol{\xi})$
attains its maximum, first fix a small closed neighborhood $\mathcal{N}$
outside the outer event horizon, and then consider a small cavity
contained in it. Let 
\[
\tilde{K}=\frac{1}{2}\sum^{3}_{i,j=1}\max_{\boldsymbol{\xi}\in\mathcal{N};\mu,\nu}\left|\frac{\partial^{2}\hat{g}_{\mu\nu}\left(t,x,y,z\right)}{\partial x^{i}\partial x^{j}}\Bigg|_{\boldsymbol{x}=\boldsymbol{\xi}}\right|,\ K=\tilde{K}r^{2}_{0}
\]
 Since 
\[
\left|x^{i}-x^{i}_{0}\right|\leq L_{\mathcal{C}}
\]
 every component $\left|\mathcal{R}_{\mu\nu}\right|$, including the
largest one, satisfies 
\[
\left|\mathcal{R}_{\mu\nu}\right|\leq\frac{1}{2}\sum^{3}_{i,j=1}\left|\frac{\partial^{2}\hat{g}_{\mu\nu}\left(t,x,y,z\right)}{\partial x^{i}\partial x^{j}}\Bigg|_{\boldsymbol{x}=\boldsymbol{\xi}}\right|\left|x^{i}-x^{i}_{0}\right|\left|x^{j}-x^{j}_{0}\right|\leq\tilde{K}L^{2}_{\mathcal{C}}=\tilde{K}r^{2}_{0}\frac{L^{2}_{\mathcal{C}}}{r^{2}_{0}}=K\varepsilon^{2}
\]
 The first $\leq$ follows because the absolute value of a sum does
not exceed the sum of the absolute values. Thus, 
\[
\max_{\mu,\nu}\left|\mathcal{R}_{\mu\nu}\right|\leq K\varepsilon^{2}
\]
 Consequently, for all $\mu,\nu$, 
\begin{align*}
\left|\hat{g}_{\mu\nu}\left(t,x,y,z\right)-\hat{g}_{\mu\nu}\left(t,x,y,z\right)\big|_{\boldsymbol{x}=\boldsymbol{x}_{0}}\right| & =\left|\hat{M}_{\mu\nu}+\mathcal{R}_{\mu\nu}\right|\\
 & \leq\left|\hat{M}_{\mu\nu}\right|+\left|\mathcal{R}_{\mu\nu}\right|\\
 & \leq\max_{\mu,\nu}\left|\hat{M}_{\mu\nu}\right|+\max_{\mu,\nu}\left|\mathcal{R}_{\mu\nu}\right|\\
 & \leq\frac{10\varepsilon}{d^{2}_{0}}+K\varepsilon^{2}
\end{align*}
 For a fixed orbit ($d_{0}$ and $K$ are finite constants) and sufficiently
small $\varepsilon$, we therefore have 
\[
\max_{\mu,\nu}\left|\hat{g}_{\mu\nu}\left(t,x,y,z\right)-\hat{g}_{\mu\nu}\left(t,x,y,z\right)\big|_{\boldsymbol{x}=\boldsymbol{x}_{0}}\right|\leq\frac{10\varepsilon}{d^{2}_{0}}+K\varepsilon^{2}\ll1.
\]
 In other words, the maximum variation of the metric components near
$\boldsymbol{x}_{0}$ satisfies 
\begin{align*}
 & \max_{\mu,\nu}|\delta\hat{g}_{\mu\nu}|=\max_{\mu,\nu}\biggl|\hat{g}_{\mu\nu}\biggl(t,x,y,z\biggr)-\hat{g}_{\mu\nu}\biggl(t,x,y,z\biggr)\big|_{\boldsymbol{x}=\boldsymbol{x}_{0}}\biggr|\ll1
\end{align*}
, so the metric components can be treated as constants throughout
the region enclosed by the parallel plates. The ratio of the bound
on the higher-order remainder to the bound on the first-order term
is 
\[
P=\frac{K\varepsilon^{2}}{10\varepsilon/d^{2}_{0}}=\frac{Kd^{2}_{0}}{10}\varepsilon=C\varepsilon,
\]
 where $C$ is a finite constant. As $\varepsilon\rightarrow0$, the
bound on the higher-order remainder decreases faster than the bound
on the first-order term. We now make a numerical estimate using the
first-order approximation. Let $M_{\text{BH}}$ be the black hole
mass in SI units and $M$ the black hole mass in geometrized units.
They are related by 
\[
M=\frac{GM_{\text{BH}}}{c^{2}},
\]
 For a black hole of one solar mass, 
\[
M=\frac{GM_{\odot}}{c^{2}}=\frac{6.674\times10^{-11}\ \text{m}^{3}\text{kg}^{-1}\text{s}^{-2}\times1.988\times10^{30}\ \text{kg}}{(299792458\ \text{m/s})^{2}}\approx1.476\ \text{km},
\]
 When the black hole mass is $6.775$ solar masses, 
\[
M\approx10\ \text{km}.
\]
 1. Taking $a=0.998M$, $r_{0}=4M=40\ \text{km}$, and $L_{\mathcal{C}}=1\ \text{mm}$,
we have $d_{0}=1-\left(2M/r_{0}\right)+\left(a^{2}/r^{2}_{0}\right)\approx0.562$
and $\varepsilon=L_{\mathcal{C}}/r_{0}=10^{-3}/\left(4\times10^{4}\right)=2.5\times10^{-8}$.
Hence, 
\[
\max_{\mu,\nu}|\delta\hat{g}^{(1)}_{\mu\nu}|\leq\frac{10\varepsilon}{d^{2}_{0}}\approx7.915\times10^{-7}\ll1,
\]
 The superscript $(1)$ denotes the first-order approximation to the
metric components. 2. Taking $a=0.998M$, $r_{0}=1.5M=15\ \text{km}$,
and $L_{\mathcal{C}}=1\ \text{mm}$, we have $d_{0}=1-\left(2M/r_{0}\right)+\left(a^{2}/r^{2}_{0}\right)\approx0.109$
and $\varepsilon=L_{\mathcal{C}}/r_{0}\approx6.667\times10^{-8}$.
Hence, 
\[
\max_{\mu,\nu}|\delta\hat{g}^{(1)}_{\mu\nu}|\leq\frac{10\varepsilon}{d^{2}_{0}}\approx5.611\times10^{-5}\ll1.
\]
 Thus, in this physical setting, with a fixed orbit and a sufficiently
small cavity, the metric components vary very little near $\boldsymbol{x}_{0}$
in the first-order approximation. Since the bound on the higher-order
remainder decreases faster than that on the first-order term as $\varepsilon\rightarrow0$,
the components of the unapproximated metric also vary very little
near $\boldsymbol{x}_{0}$. The metric components can therefore reasonably
be approximated as constants, using the zeroth-order approximation
with small corrections.

\Needspace{110pt}\par\addvspace{20pt}\noindent\begin{minipage}{\linewidth}\begingroup\fontsize{12}{15}\selectfont\setlength{\abovedisplayskip}{0pt}\setlength{\belowdisplayskip}{0pt}\setlength{\abovedisplayshortskip}{0pt}\setlength{\belowdisplayshortskip}{0pt}\phantomsection\addcontentsline{toc}{section}{III. The Klein-Gordon Equation and Quantum Field Mode Solutions}
\[
\begin{gathered}\text{\bfseries III. The Klein-Gordon Equation}\\
\text{\bfseries and Quantum Field Mode Solutions}
\end{gathered}
\]
\endgroup\end{minipage}\par\nopagebreak[4]\vspace{20pt}\nopagebreak[4]

The Klein-Gordon equation for a scalar field in curved spacetime is
given in Spacetime and Geometry: An Introduction to General Relativity
by Sean M. Carroll, published by Cambridge University Press. With
the metric signature $(-,+,+,+)$, which is opposite to that used
in the present paper, it reads 
\[
\Box\Psi-m^{2}\Psi-\xi R\Psi=0
\]
 For a massless scalar field in vacuum, the equation becomes 
\begin{align*}
 & \Box\Psi=0\\
 & \nabla_{\mu}\nabla^{\mu}\Psi=0\\
 & \nabla_{\mu}\biggl(g^{\mu\nu}\partial_{\nu}\Psi\biggr)=0\\
 & \qquad\text{\parbox{155mm}{For a scalar field \ensuremath{\Psi}, the covariant derivative \ensuremath{\nabla_{\mu}} and the ordinary derivative \ensuremath{\partial_{\mu}} coincide}}\\
 & \partial_{\mu}\biggl(g^{\mu\nu}\partial_{\nu}\Psi\biggr)+\varGamma^{\mu}_{\!\mu\alpha}\biggl(g^{\alpha\nu}\partial_{\nu}\Psi\biggr)=0\\
 & \qquad\text{\parbox{155mm}{For a vector field \ensuremath{V^{\mu}=g^{\mu\nu}\partial_{\nu}\Psi}, the covariant derivative \ensuremath{\nabla_{\mu}V^{\mu}=\partial_{\mu}V^{\mu}+\varGamma^{\mu}_{\!\mu\alpha}V^{\alpha}}}}\\
 & \partial_{\mu}\biggl(g^{\mu\nu}\partial_{\nu}\Psi\biggr)+\frac{1}{\sqrt{-g}}\partial_{\alpha}\biggl(\sqrt{-g}\biggr)\biggl(g^{\alpha\nu}\partial_{\nu}\Psi\biggr)=0\\
 & \qquad\text{\parbox{155mm}{It can be shown that \ensuremath{\varGamma^{\mu}_{\!\mu\alpha}}=\ensuremath{\frac{1}{\sqrt{-g}}\partial_{\alpha}\biggl(\sqrt{-g}\biggr)}; g=\ensuremath{\det\biggl(g_{\mu\nu}\biggr)}}}\\
 & \frac{1}{\sqrt{-g}}\biggl[\sqrt{-g}\partial_{\mu}\biggl(g^{\mu\nu}\partial_{\nu}\Psi\biggr)+\partial_{\mu}\biggl(\sqrt{-g}\biggr)\biggl(g^{\mu\nu}\partial_{\nu}\Psi\biggr)\biggr]=0\\
 & \frac{1}{\sqrt{-g}}\biggl[\partial_{\mu}\biggl(\sqrt{-g}g^{\mu\nu}\partial_{\nu}\Psi\biggr)\biggr]=0\\
 & \qquad\text{\parbox{155mm}{product rule for differentiation}}
\end{align*}
 Since the metric matrix is treated as constant, the factor $\sqrt{-\hat{g}}\hat{g}^{\mu\nu}$
(with $\sqrt{-\hat{g}}=1$) can be taken outside the derivative: $\frac{1}{\sqrt{-\hat{g}}}\left[\partial_{\mu}\left(\sqrt{-\hat{g}}\hat{g}^{\mu\nu}\partial_{\nu}\Psi\right)\right]=\frac{\sqrt{-\hat{g}}}{\sqrt{-\hat{g}}}\hat{g}^{\mu\nu}\partial_{\mu}\partial_{\nu}\Psi$.
Thus, we obtain the following equation. Here, $\Psi$ represents only
one mode of the field, corresponding to a discrete frequency $\omega_{n}$
and a continuous wave vector $\ensuremath{\vec{k}_{\perp}}$, rather
than the full field.

\begin{align*}
\hat{g}^{\mu\nu}\partial_{\mu}\partial_{\nu}\Psi & =0\\
\left(\frac{A}{\Delta r^{2}}\partial^{2}_{t}+2\frac{A}{\Delta r}\left(\omega_{d}-\Omega\right)\partial_{t}\partial_{x}-\frac{r^{2}}{\Delta}C^{-2}\left(\Omega\right)\partial^{2}_{x}-\frac{\Delta}{r^{2}}\partial^{2}_{y}-\partial^{2}_{z}\right)e^{-i\omega_{n}t}e^{i\vec{k}_{\bot}\cdot\vec{x}_{\bot}}\chi\left(x\right) & =0\\
\left[\frac{A}{\Delta r^{2}}\left(-\omega^{2}_{n}\right)+2\frac{A}{\Delta r}\left(\omega_{d}-\Omega\right)\left(-i\omega_{n}\right)\partial_{x}-\frac{r^{2}}{\Delta}C^{-2}\left(\Omega\right)\partial^{2}_{x}-\frac{\Delta}{r^{2}}\left(-k^{2}_{y}\right)-\left(-k^{2}_{z}\right)\right]e^{-i\omega_{n}t}e^{i\vec{k}_{\bot}\cdot\vec{x}_{\bot}}\chi\left(x\right) & =0\\
e^{-i\omega_{n}t}e^{i\vec{k}_{\bot}\cdot\vec{x}_{\bot}}\left[-\frac{r^{2}}{\Delta}C^{-2}\left(\Omega\right)\partial^{2}_{x}\chi+2i\omega_{n}\frac{A}{\Delta r}\left(\Omega-\omega_{d}\right)\partial_{x}\chi+\left(-\frac{A}{\Delta r^{2}}\omega^{2}_{n}+\frac{\Delta}{r^{2}}k^{2}_{y}+k^{2}_{z}\right)\chi\right] & =0
\end{align*}
 where 
\begin{align*}
\partial^{2}_{t}\Psi & =\partial^{2}_{t}e^{-i\omega_{n}t}e^{i\vec{k}_{\bot}\cdot\vec{x}_{\bot}}\chi\left(x\right)\\
 & =\left(-i\omega_{n}\right)^{2}e^{-i\omega_{n}t}e^{i\vec{k}_{\bot}\cdot\vec{x}_{\bot}}\chi\left(x\right)\\
 & =-\omega^{2}_{n}e^{-i\omega_{n}t}e^{i\vec{k}_{\bot}\cdot\vec{x}_{\bot}}\chi\left(x\right)\\
 & =-\omega^{2}_{n}\Psi
\end{align*}

\begin{align*}
\partial_{t}\partial_{x}\Psi & =\partial_{t}\partial_{x}e^{-i\omega_{n}t}e^{i\vec{k}_{\bot}\cdot\vec{x}_{\bot}}\chi\left(x\right)\\
 & =\left(-i\omega_{n}\right)e^{-i\omega_{n}t}e^{i\vec{k}_{\bot}\cdot\vec{x}_{\bot}}\partial_{x}\chi\left(x\right)\\
 & =\left(-i\omega_{n}\right)\partial_{x}\Psi
\end{align*}

\begin{align*}
\partial^{2}_{y}\Psi & =\partial^{2}_{y}e^{-i\omega_{n}t}e^{i\vec{k}_{\bot}\cdot\vec{x}_{\bot}}\chi\left(x\right)\\
 & =\partial^{2}_{y}e^{-i\omega_{n}t}e^{i\left(k_{y}y+k_{z}z\right)}\chi\left(x\right)\\
 & =\left(ik_{y}\right)^{2}e^{-i\omega_{n}t}e^{i\left(k_{y}y+k_{z}z\right)}\chi\left(x\right)\\
 & =-k^{2}_{y}\Psi
\end{align*}
 Similarly, 
\[
\partial^{2}_{z}\Psi=-k^{2}_{z}\Psi
\]

Seek a solution of 
\[
-\frac{r^{2}}{\Delta}C^{-2}\left(\Omega\right)\partial^{2}_{x}\chi+2i\omega_{n}\frac{A}{\Delta r}\left(\Omega-\omega_{d}\right)\partial_{x}\chi+\left(-\frac{A}{\Delta r^{2}}\omega^{2}_{n}+\frac{\Delta}{r^{2}}k^{2}_{y}+k^{2}_{z}\right)\chi=0
\]
 in the form 
\[
\chi=e^{\lambda x}
\]
 (corresponding to $\Psi=e^{-i\omega_{n}t}e^{i\vec{k}_{\bot}\cdot\vec{x}_{\bot}}e^{\lambda x}$).
Then 
\begin{align*}
-\frac{r^{2}}{\Delta}C^{-2}\left(\Omega\right)\partial^{2}_{x}e^{\lambda x}+2i\omega_{n}\frac{A}{\Delta r}\left(\Omega-\omega_{d}\right)\partial_{x}e^{\lambda x}+\left(-\frac{A}{\Delta r^{2}}\omega^{2}_{n}+\frac{\Delta}{r^{2}}k^{2}_{y}+k^{2}_{z}\right)e^{\lambda x} & =0\\
\left[-\frac{r^{2}}{\Delta}C^{-2}\left(\Omega\right)\lambda^{2}+2i\omega_{n}\frac{A}{\Delta r}\left(\Omega-\omega_{d}\right)\lambda+\left(-\frac{A}{\Delta r^{2}}\omega^{2}_{n}+\frac{\Delta}{r^{2}}k^{2}_{y}+k^{2}_{z}\right)\right]e^{\lambda x} & =0
\end{align*}

This requires the characteristic equation 
\[
-\frac{r^{2}}{\Delta}C^{-2}\left(\Omega\right)\lambda^{2}+2i\omega_{n}\frac{A}{\Delta r}\left(\Omega-\omega_{d}\right)\lambda-\frac{A}{\Delta r^{2}}\omega^{2}_{n}+\frac{\Delta}{r^{2}}k^{2}_{y}+k^{2}_{z}=0
\]
 For the quadratic equation $a\lambda^{2}+b\lambda+c=0$, the discriminant
is 
\begin{align*}
\tilde{\Delta} & =b^{2}-4ac\\
 & =\left[2i\omega_{n}\frac{A}{\Delta r}\left(\Omega-\omega_{d}\right)\right]^{2}-4\left[-\frac{r^{2}}{\Delta}C^{-2}\left(\Omega\right)\right]\left(-\frac{A}{\Delta r^{2}}\omega^{2}_{n}+\frac{\Delta}{r^{2}}k^{2}_{y}+k^{2}_{z}\right)\\
 & =\left[2i\omega_{n}\frac{A}{\Delta r}\left(\Omega-\omega_{d}\right)\right]^{2}+4\frac{r^{2}}{\Delta}C^{-2}\left(\Omega\right)\left(-\frac{A}{\Delta r^{2}}\omega^{2}_{n}+\frac{\Delta}{r^{2}}k^{2}_{y}+k^{2}_{z}\right)\\
 & =-\left[2\omega_{n}\frac{A}{\Delta r}\left(\Omega-\omega_{d}\right)\right]^{2}+4\frac{r^{2}}{\Delta}\left\{ \frac{\Delta r^{2}}{A}\left[1-\frac{A^{2}}{\Delta r^{4}}\left(\Omega-\omega_{d}\right)^{2}\right]\right\} \left(-\frac{A}{\Delta r^{2}}\omega^{2}_{n}+\frac{\Delta}{r^{2}}k^{2}_{y}+k^{2}_{z}\right)\\
 & =-\left[2\omega_{n}\frac{A}{\Delta r}\left(\Omega-\omega_{d}\right)\right]^{2}+4\left[\frac{r^{4}}{A}-\frac{A}{\Delta}\left(\Omega-\omega_{d}\right)^{2}\right]\left(-\frac{A}{\Delta r^{2}}\omega^{2}_{n}+\frac{\Delta}{r^{2}}k^{2}_{y}+k^{2}_{z}\right)\\
 & =-\left[2\omega_{n}\frac{A}{\Delta r}\left(\Omega-\omega_{d}\right)\right]^{2}+4\frac{r^{4}}{A}\left(-\frac{A}{\Delta r^{2}}\omega^{2}_{n}+\frac{\Delta}{r^{2}}k^{2}_{y}+k^{2}_{z}\right)+\left[2\omega_{n}\frac{A}{\Delta r}\left(\Omega-\omega_{d}\right)\right]^{2}-4\frac{A}{\Delta}\left(\Omega-\omega_{d}\right)^{2}\left(\frac{\Delta}{r^{2}}k^{2}_{y}+k^{2}_{z}\right)\\
 & =4\frac{r^{4}}{A}\left(-\frac{A}{\Delta r^{2}}\omega^{2}_{n}+\frac{\Delta}{r^{2}}k^{2}_{y}+k^{2}_{z}\right)-4\frac{A}{\Delta}\left(\Omega-\omega_{d}\right)^{2}\left(\frac{\Delta}{r^{2}}k^{2}_{y}+k^{2}_{z}\right)\\
 & =4\frac{r^{4}}{A}\left(-\frac{A}{\Delta r^{2}}\omega^{2}_{n}\right)+4\frac{r^{4}}{A}\left(\frac{\Delta}{r^{2}}k^{2}_{y}+k^{2}_{z}\right)-4\frac{A}{\Delta}\left(\Omega-\omega_{d}\right)^{2}\left(\frac{\Delta}{r^{2}}k^{2}_{y}+k^{2}_{z}\right)\\
 & =-4\frac{r^{2}}{\Delta}\omega^{2}_{n}+4\left[\frac{r^{4}}{A}-\frac{A}{\Delta}\left(\Omega-\omega_{d}\right)^{2}\right]\left(\frac{\Delta}{r^{2}}k^{2}_{y}+k^{2}_{z}\right)\\
 & =-4\frac{r^{2}}{\Delta}\omega^{2}_{n}+4\frac{r^{4}}{A}\left[1-\frac{A^{2}}{\Delta r^{4}}\left(\Omega-\omega_{d}\right)^{2}\right]\left(\frac{\Delta}{r^{2}}k^{2}_{y}+k^{2}_{z}\right)\\
 & =-4\frac{r^{2}}{\Delta}\omega^{2}_{n}+4\frac{r^{2}}{\Delta}\left\{ \frac{\Delta r^{2}}{A}\left[1-\frac{A^{2}}{\Delta r^{4}}\left(\Omega-\omega_{d}\right)^{2}\right]\right\} \left(\frac{\Delta}{r^{2}}k^{2}_{y}+k^{2}_{z}\right)\\
 & =-4\frac{r^{2}}{\Delta}\omega^{2}_{n}+4\frac{r^{2}}{\Delta}C^{-2}\left(\Omega\right)\left(\frac{\Delta}{r^{2}}k^{2}_{y}+k^{2}_{z}\right)
\end{align*}
 The roots of the characteristic equation are 
\begin{align*}
\lambda_{1,2} & =\frac{-b\pm\sqrt{\tilde{\Delta}}}{2a}\\
 & =\frac{-b\pm\sqrt{i^{2}\cdot i^{-2}\tilde{\Delta}}}{2a}\\
 & =\frac{-b\pm i\sqrt{\left(-1\right)\tilde{\Delta}}}{2a}\\
 & =\frac{-b}{2a}\pm i\frac{\sqrt{-\tilde{\Delta}}}{2a}\\
 & =\frac{-\left[2i\omega_{n}\frac{A}{\Delta r}\left(\Omega-\omega_{d}\right)\right]}{2\left[-\frac{r^{2}}{\Delta}C^{-2}\left(\Omega\right)\right]}\pm i\frac{\sqrt{-\left[-4\frac{r^{2}}{\Delta}\omega^{2}_{n}+4\frac{r^{2}}{\Delta}C^{-2}\left(\Omega\right)\left(\frac{\Delta}{r^{2}}k^{2}_{y}+k^{2}_{z}\right)\right]}}{2\left[-\frac{r^{2}}{\Delta}C^{-2}\left(\Omega\right)\right]}
\end{align*}
 Let 
\begin{align*}
\tilde{\alpha} & =-\frac{b}{2a}\\
 & =\frac{-\left[2i\omega_{n}\frac{A}{\Delta r}\left(\Omega-\omega_{d}\right)\right]}{2\left[-\frac{r^{2}}{\Delta}C^{-2}\left(\Omega\right)\right]}\\
 & =i\frac{\omega_{n}A\left(\Omega-\omega_{d}\right)}{r^{3}C^{-2}\left(\Omega\right)}\\
 & =-i\frac{\omega_{n}A\left(\omega_{d}-\Omega\right)}{r^{3}C^{-2}\left(\Omega\right)}
\end{align*}
 {}
\begin{align*}
\tilde{\beta} & =\frac{\sqrt{-\tilde{\Delta}}}{2a}\\
 & =\frac{\sqrt{-\left[-4\frac{r^{2}}{\Delta}\omega^{2}_{n}+4\frac{r^{2}}{\Delta}C^{-2}\left(\Omega\right)\left(\frac{\Delta}{r^{2}}k^{2}_{y}+k^{2}_{z}\right)\right]}}{2\left[-\frac{r^{2}}{\Delta}C^{-2}\left(\Omega\right)\right]}
\end{align*}
 Then $\lambda_{1,2}=\tilde{\alpha}\pm i\tilde{\beta}$, and the general
solution is 
\begin{align*}
\chi & =\tilde{C}_{1}e^{\lambda_{1}x}+\tilde{C}_{2}e^{\lambda_{2}x}\\
 & =\tilde{C}_{1}e^{\left(\tilde{\alpha}+i\tilde{\beta}\right)x}+\tilde{C}_{2}e^{\left(\tilde{\alpha}-i\tilde{\beta}\right)x}\\
 & =\tilde{C}_{1}e^{\tilde{\alpha}x}e^{i\tilde{\beta}x}+\tilde{C}_{2}e^{\tilde{\alpha}x}e^{-i\tilde{\beta}x}\\
 & =e^{\tilde{\alpha}x}\left[\tilde{C}_{1}\left(\cos\tilde{\beta}x+i\sin\tilde{\beta}x\right)+\tilde{C}_{2}\left(\cos\tilde{\beta}x-i\sin\tilde{\beta}x\right)\right]\\
 & =e^{\tilde{\alpha}x}\left[\left(\tilde{C}_{1}+\tilde{C}_{2}\right)\cos\tilde{\beta}x+i\left(\tilde{C}_{1}-\tilde{C}_{2}\right)\sin\tilde{\beta}x\right]\\
 & =e^{\tilde{\alpha}x}\left[C_{1}\cos\tilde{\beta}x+C_{2}\sin\tilde{\beta}x\right]
\end{align*}
 where $\tilde{C}_{1}$ and $\tilde{C}_{2}$ are arbitrary complex
constants. Define $C_{1}=\tilde{C}_{1}+\tilde{C}_{2}$ and $C_{2}=i\left(\tilde{C}_{1}-\tilde{C}_{2}\right)$.
We may choose $\tilde{C}_{1}=\left(C_{1}-iC_{2}\right)/2$ and $\tilde{C}_{2}=\left(C_{1}+iC_{2}\right)/2$
so that $C_{1}$ and $C_{2}$ are arbitrary real constants. The Dirichlet
boundary conditions 
\[
\Psi\left(t,x=0,y,z\right)=\Psi\left(t,x=L,y,z\right)=0
\]
 give 
\begin{align*}
\chi\left(x=0\right) & =e^{\tilde{\alpha}\cdot0}\left(C_{1}\cos\tilde{\beta}\cdot0+C_{2}\sin\tilde{\beta}\cdot0\right)\\
 & =C_{1}=0
\end{align*}
 {}
\begin{align*}
\chi\left(x=L\right) & =e^{\tilde{\alpha}\cdot L}\left(0\cdot\cos\tilde{\beta}\cdot L+C_{2}\sin\tilde{\beta}\cdot L\right)\\
 & =C_{2}e^{-i\frac{\omega_{n}A\left(\omega_{d}-\Omega\right)}{r^{3}C^{-2}\left(\Omega\right)}\cdot L}\sin\tilde{\beta}L\\
 & =0
\end{align*}
 To obtain a nontrivial solution, that is, a solution that is not
identically zero, we require 
\[
\sin\tilde{\beta}L=0
\]
 Hence, 
\begin{align*}
\tilde{\beta}L & =n\pi\\
\Rightarrow & \tilde{\beta}=\frac{n\pi}{L}
\end{align*}
 and we obtain 
\[
\chi=C_{2}e^{-i\frac{\omega_{n}A\left(\omega_{d}-\Omega\right)}{r^{3}C^{-2}\left(\Omega\right)}x}\sin\left(\frac{n\pi}{L}x\right)
\]

Method 2: 
\[
\gamma_{2}\partial^{2}_{x}\chi+\gamma_{1}\partial_{x}\chi+\gamma_{0}\chi=0
\]
 where 
\begin{align*}
\gamma_{2} & =-\frac{r^{2}}{\Delta}C^{-2}\left(\Omega\right)\\
\gamma_{1} & =2i\omega_{n}\frac{A}{\Delta r}\left(\Omega-\omega_{d}\right)\\
\gamma_{0} & =-\frac{A}{\Delta r^{2}}\omega^{2}_{n}+\frac{\Delta}{r^{2}}k^{2}_{y}+k^{2}_{z}
\end{align*}
 Let 
\[
\chi=e^{-i\beta x}\tilde{\chi}
\]
 Then 
\begin{align*}
\partial_{x}\chi & =\left(-i\beta\right)e^{-i\beta x}\tilde{\chi}+e^{-i\beta x}\partial_{x}\tilde{\chi}\\
 & =e^{-i\beta x}\left(-i\beta\tilde{\chi}+\partial_{x}\tilde{\chi}\right)
\end{align*}
 {}
\begin{align*}
\partial^{2}_{x}\chi & =\partial_{x}\left[e^{-i\beta x}\left(-i\beta\tilde{\chi}+\partial_{x}\tilde{\chi}\right)\right]\\
 & =\left(-i\beta\right)e^{-i\beta x}\left(-i\beta\tilde{\chi}+\partial_{x}\tilde{\chi}\right)+e^{-i\beta x}\left(-i\beta\partial_{x}\tilde{\chi}+\partial^{2}_{x}\tilde{\chi}\right)\\
 & =e^{-i\beta x}\left[-\beta^{2}\tilde{\chi}-i2\beta\partial_{x}\tilde{\chi}+\partial^{2}_{x}\tilde{\chi}\right]
\end{align*}
 Substitution into the original equation gives 
\begin{align*}
\gamma_{2}e^{-i\beta x}\left[-\beta^{2}\tilde{\chi}-i2\beta\partial_{x}\tilde{\chi}+\partial^{2}_{x}\tilde{\chi}\right]+\gamma_{1}e^{-i\beta x}\left(-i\beta\tilde{\chi}+\partial_{x}\tilde{\chi}\right)+\gamma_{0}e^{-i\beta x}\tilde{\chi} & =0\\
\gamma_{2}\left[-\beta^{2}\tilde{\chi}-i2\beta\partial_{x}\tilde{\chi}+\partial^{2}_{x}\tilde{\chi}\right]+\gamma_{1}\left(-i\beta\tilde{\chi}+\partial_{x}\tilde{\chi}\right)+\gamma_{0}\tilde{\chi} & =0\\
\gamma_{2}\partial^{2}_{x}\tilde{\chi}+\left(\gamma_{1}-i2\beta\gamma_{2}\right)\partial_{x}\tilde{\chi}+\left(\gamma_{0}-\beta^{2}\gamma_{2}-i\beta\gamma_{1}\right)\tilde{\chi} & =0
\end{align*}
 Set 
\[
\gamma_{1}-i2\beta\gamma_{2}=0
\]
 Then 
\begin{align*}
\beta & =-\frac{\gamma_{1}}{2\gamma_{2}}i\\
 & =-\frac{2i\omega_{n}\frac{A}{\Delta r}\left(\Omega-\omega_{d}\right)}{2\left(-\frac{r^{2}}{\Delta}C^{-2}\left(\Omega\right)\right)}i\\
 & =\frac{\omega_{n}C^{2}\left(\Omega\right)A}{r^{3}}\left(\omega_{d}-\Omega\right)\\
 & =\frac{b_{n}A}{r^{3}}\left(\omega_{d}-\Omega\right)
\end{align*}
 where we have defined 
\[
b_{n}=\omega_{n}C^{2}\left(\Omega\right)
\]
 We have 
\begin{align*}
\gamma_{2}\partial^{2}_{x}\tilde{\chi}+\left(\gamma_{0}-\beta^{2}\gamma_{2}-i\beta\gamma_{1}\right)\tilde{\chi} & =0\\
\gamma_{2}\partial^{2}_{x}\tilde{\chi}+\left[\gamma_{0}-\beta^{2}\gamma_{2}-i\beta\left(i2\beta\gamma_{2}\right)\right]\tilde{\chi} & =0\\
\gamma_{2}\partial^{2}_{x}\tilde{\chi}+\left(\gamma_{0}+\beta^{2}\gamma_{2}\right)\tilde{\chi} & =0\\
\partial^{2}_{x}\tilde{\chi}+\frac{\left(\gamma_{0}+\beta^{2}\gamma_{2}\right)}{\gamma_{2}}\tilde{\chi} & =0
\end{align*}
 Seek a solution of this equation in the form 
\[
\tilde{\chi}=e^{\lambda x}
\]
 (corresponding to $\Psi=e^{-i\omega_{n}t}e^{i\vec{k}_{\bot}\cdot\vec{x}_{\bot}}e^{-i\beta x}e^{\lambda x}$).
Then 
\[
\left[\lambda^{2}+\frac{\left(\gamma_{0}+\beta^{2}\gamma_{2}\right)}{\gamma_{2}}\right]e^{\lambda x}=0
\]
 and hence 
\[
\lambda^{2}+\frac{\left(\gamma_{0}+\beta^{2}\gamma_{2}\right)}{\gamma_{2}}=0
\]
 {}
\[
\lambda_{1,2}=\pm\sqrt{-\frac{\left(\gamma_{0}+\beta^{2}\gamma_{2}\right)}{\gamma_{2}}}
\]
 Thus, 
\begin{align*}
 & \chi=e^{-i\beta x}\biggl(\tilde{C}_{1}e^{\lambda_{1}x}+\tilde{C}_{2}e^{\lambda_{2}x}\biggr)\\
 & =e^{-i\beta x}\biggl(\tilde{C}_{1}e^{\sqrt{-\frac{\biggl(\gamma_{0}+\beta^{2}\gamma_{2}\biggr)}{\gamma_{2}}}x}+\tilde{C}_{2}e^{-\sqrt{-\frac{\biggl(\gamma_{0}+\beta^{2}\gamma_{2}\biggr)}{\gamma_{2}}}x}\biggr)\\
 & =e^{-i\beta x}\biggl(\tilde{C}_{1}e^{i\sqrt{\frac{-\omega^{2}_{n}C^{2}\biggl(\Omega\biggr)+\frac{\Delta}{r^{2}}k^{2}_{y}+k^{2}_{z}}{\gamma_{2}}}x}+\tilde{C}_{2}e^{-i\sqrt{\frac{-\omega^{2}_{n}C^{2}\biggl(\Omega\biggr)+\frac{\Delta}{r^{2}}k^{2}_{y}+k^{2}_{z}}{\gamma_{2}}}x}\biggr)\\
 & =e^{-i\beta x}\biggl[\tilde{C}_{1}\biggl(\cos\sqrt{\frac{-\omega^{2}_{n}C^{2}\biggl(\Omega\biggr)+\frac{\Delta}{r^{2}}k^{2}_{y}+k^{2}_{z}}{\gamma_{2}}}x+i\sin\sqrt{\frac{-\omega^{2}_{n}C^{2}\biggl(\Omega\biggr)+\frac{\Delta}{r^{2}}k^{2}_{y}+k^{2}_{z}}{\gamma_{2}}}x\biggr)+\tilde{C}_{2}\\
 & \qquad\biggl(\cos\sqrt{\frac{-\omega^{2}_{n}C^{2}\biggl(\Omega\biggr)+\frac{\Delta}{r^{2}}k^{2}_{y}+k^{2}_{z}}{\gamma_{2}}}x-i\sin\sqrt{\frac{-\omega^{2}_{n}C^{2}\biggl(\Omega\biggr)+\frac{\Delta}{r^{2}}k^{2}_{y}+k^{2}_{z}}{\gamma_{2}}}x\biggr)\biggr]\\
 & =e^{-i\beta x}\biggl[\biggl(\tilde{C}_{1}+\tilde{C}_{2}\biggr)\cos\sqrt{\frac{-\omega^{2}_{n}C^{2}\biggl(\Omega\biggr)+\frac{\Delta}{r^{2}}k^{2}_{y}+k^{2}_{z}}{\gamma_{2}}}x+i\biggl(\tilde{C}_{1}-\tilde{C}_{2}\biggr)\sin\sqrt{\frac{-\omega^{2}_{n}C^{2}\biggl(\Omega\biggr)+\frac{\Delta}{r^{2}}k^{2}_{y}+k^{2}_{z}}{\gamma_{2}}}x\biggr]\\
 & =e^{-i\frac{\omega_{n}A\biggl(\omega_{d}-\Omega\biggr)}{r^{3}C^{-2}\biggl(\Omega\biggr)}x}\biggl(C_{1}\cos\sqrt{\frac{-\omega^{2}_{n}C^{2}\biggl(\Omega\biggr)+\frac{\Delta}{r^{2}}k^{2}_{y}+k^{2}_{z}}{\gamma_{2}}}x+C_{2}\sin\sqrt{\frac{-\omega^{2}_{n}C^{2}\biggl(\Omega\biggr)+\frac{\Delta}{r^{2}}k^{2}_{y}+k^{2}_{z}}{\gamma_{2}}}x\biggr)
\end{align*}
 where 
\begin{align*}
\gamma_{0}+\beta^{2}\gamma_{2} & =-\frac{A}{\Delta r^{2}}\omega^{2}_{n}+\frac{\Delta}{r^{2}}k^{2}_{y}+k^{2}_{z}+\left[\frac{\omega_{n}C^{2}\left(\Omega\right)A}{r^{3}}\left(\omega_{d}-\Omega\right)\right]^{2}\left[-\frac{r^{2}}{\Delta}C^{-2}\left(\Omega\right)\right]\\
 & =-\omega^{2}_{n}\left[\frac{A}{\Delta r^{2}}+\frac{C^{2}\left(\Omega\right)A^{2}}{r^{6}}\left(\omega_{d}-\Omega\right)^{2}\frac{r^{2}}{\Delta}\right]+\frac{\Delta}{r^{2}}k^{2}_{y}+k^{2}_{z}\\
 & =-\omega^{2}_{n}C^{2}\left(\Omega\right)\left[C^{-2}\left(\Omega\right)\frac{A}{\Delta r^{2}}+\frac{A^{2}}{\Delta r^{4}}\left(\omega_{d}-\Omega\right)^{2}\right]+\frac{\Delta}{r^{2}}k^{2}_{y}+k^{2}_{z}\\
 & =-\omega^{2}_{n}C^{2}\left(\Omega\right)\left(\frac{\Delta r^{2}}{A}\left[1-\frac{A^{2}}{\Delta r^{4}}\left(\Omega-\omega_{d}\right)^{2}\right]\frac{A}{\Delta r^{2}}+\frac{A^{2}}{\Delta r^{4}}\left(\omega_{d}-\Omega\right)^{2}\right)+\frac{\Delta}{r^{2}}k^{2}_{y}+k^{2}_{z}\\
 & =-\omega^{2}_{n}C^{2}\left(\Omega\right)+\frac{\Delta}{r^{2}}k^{2}_{y}+k^{2}_{z}
\end{align*}
 The coefficients $\tilde{C}_{1}$ and $\tilde{C}_{2}$ are arbitrary
complex constants. Define $C_{1}=\tilde{C}_{1}+\tilde{C}_{2}$ and
$C_{2}=i\left(\tilde{C}_{1}-\tilde{C}_{2}\right)$. We may choose
$\tilde{C}_{1}=\left(C_{1}-iC_{2}\right)/2$ and $\tilde{C}_{2}=\left(C_{1}+iC_{2}\right)/2$
so that $C_{1}$ and $C_{2}$ are arbitrary real constants. The Dirichlet
boundary conditions 
\[
\Psi\left(t,x=0,y,z\right)=\Psi\left(t,x=L,y,z\right)=0
\]
 give 
\begin{align*}
\chi\left(x=0\right) & =e^{-i\frac{\omega_{n}A\left(\omega_{d}-\Omega\right)}{r^{3}C^{-2}\left(\Omega\right)}\cdot0}\left(C_{1}\cos\sqrt{\frac{-\omega^{2}_{n}C^{2}\left(\Omega\right)+\frac{\Delta}{r^{2}}k^{2}_{y}+k^{2}_{z}}{\gamma_{2}}}\cdot0+C_{2}\sin\sqrt{\frac{-\omega^{2}_{n}C^{2}\left(\Omega\right)+\frac{\Delta}{r^{2}}k^{2}_{y}+k^{2}_{z}}{\gamma_{2}}}\cdot0\right)\\
 & =C_{1}=0
\end{align*}
 {}
\begin{align*}
\chi\left(x=L\right) & =e^{-i\frac{\omega_{n}A\left(\omega_{d}-\Omega\right)}{r^{3}C^{-2}\left(\Omega\right)}\cdot L}\left(0\cdot\cos\sqrt{\frac{-\omega^{2}_{n}C^{2}\left(\Omega\right)+\frac{\Delta}{r^{2}}k^{2}_{y}+k^{2}_{z}}{\gamma_{2}}}\cdot L+C_{2}\sin\sqrt{\frac{-\omega^{2}_{n}C^{2}\left(\Omega\right)+\frac{\Delta}{r^{2}}k^{2}_{y}+k^{2}_{z}}{\gamma_{2}}}\cdot L\right)\\
 & =C_{2}e^{-i\frac{\omega_{n}A\left(\omega_{d}-\Omega\right)}{r^{3}C^{-2}\left(\Omega\right)}L}\sin\sqrt{\frac{-\omega^{2}_{n}C^{2}\left(\Omega\right)+\frac{\Delta}{r^{2}}k^{2}_{y}+k^{2}_{z}}{\gamma_{2}}}L\\
 & =0
\end{align*}
 To obtain a nontrivial solution, that is, a solution that is not
identically zero, we require 
\[
\sin\sqrt{\frac{-\omega^{2}_{n}C^{2}\left(\Omega\right)+\frac{\Delta}{r^{2}}k^{2}_{y}+k^{2}_{z}}{\gamma_{2}}}L=0
\]
 Hence, 
\begin{align*}
\sqrt{\frac{-\omega^{2}_{n}C^{2}\left(\Omega\right)+\frac{\Delta}{r^{2}}k^{2}_{y}+k^{2}_{z}}{\gamma_{2}}}L & =n\pi\\
\Rightarrow & \sqrt{\frac{-\omega^{2}_{n}C^{2}\left(\Omega\right)+\frac{\Delta}{r^{2}}k^{2}_{y}+k^{2}_{z}}{\gamma_{2}}}=\frac{n\pi}{L}
\end{align*}
 and we obtain 
\[
\chi=C_{2}e^{-i\frac{\omega_{n}A\left(\omega_{d}-\Omega\right)}{r^{3}C^{-2}\left(\Omega\right)}x}\sin\left(\frac{n\pi}{L}x\right)
\]

Solving gives the following relation for $\omega_{n}$: 
\begin{align*}
\sqrt{\frac{-\omega^{2}_{n}C^{2}\left(\Omega\right)+\frac{\Delta}{r^{2}}k^{2}_{y}+k^{2}_{z}}{\gamma_{2}}} & =\frac{n\pi}{L}\\
\frac{-\omega^{2}_{n}C^{2}\left(\Omega\right)+\frac{\Delta}{r^{2}}k^{2}_{y}+k^{2}_{z}}{-\frac{r^{2}}{\Delta}C^{-2}\left(\Omega\right)} & =\left(\frac{n\pi}{L}\right)^{2}\\
\omega^{2}_{n}\frac{\Delta}{r^{2}}C^{4}\left(\Omega\right)-\frac{\Delta}{r^{2}}C^{2}\left(\Omega\right)\left(\frac{\Delta}{r^{2}}k^{2}_{y}+k^{2}_{z}\right) & =\left(\frac{n\pi}{L}\right)^{2}\\
\omega^{2}_{n} & =\frac{r^{2}}{\Delta}C^{-4}\left(\Omega\right)\left(\frac{n\pi}{L}\right)^{2}+C^{-2}\left(\Omega\right)\left(\frac{\Delta}{r^{2}}k^{2}_{y}+k^{2}_{z}\right)\\
 & =\frac{r^{2}}{\Delta}C^{-4}\left(\Omega\right)\left[\left(\frac{n\pi}{L}\right)^{2}+\frac{\Delta}{r^{2}}C^{2}\left(\Omega\right)\left(\frac{\Delta}{r^{2}}k^{2}_{y}+k^{2}_{z}\right)\right]\\
\omega_{n} & =\frac{r}{\sqrt{\Delta}}C^{-2}\left(\Omega\right)\sqrt{\left(\frac{n\pi}{L}\right)^{2}+\frac{\Delta}{r^{2}}C^{2}\left(\Omega\right)\left(\frac{\Delta}{r^{2}}k^{2}_{y}+k^{2}_{z}\right)}
\end{align*}

The complete solution is 
\[
\Psi=C_{2}e^{-i\omega_{n}t}e^{i\vec{k}_{\bot}\cdot\vec{x}_{\bot}}e^{-i\frac{\omega_{n}A\left(\omega_{d}-\Omega\right)}{r^{3}C^{-2}\left(\Omega\right)}x}\sin\left(\frac{n\pi}{L}x\right)
\]

In fact, $\Psi$ is the solution for mode $n$, rather than the full
field. We use $u_{n}$ in place of $\Psi$ to denote this mode solution:
\[
u_{n}=C_{2}e^{-i\omega_{n}t}e^{i\vec{k}_{\bot}\cdot\vec{x}_{\bot}}e^{-i\frac{\omega_{n}A\left(\omega_{d}-\Omega\right)}{r^{3}C^{-2}\left(\Omega\right)}x}\sin\left(\frac{n\pi}{L}x\right)
\]
 or 
\[
u_{n}=N_{n}e^{-i\omega_{n}t}e^{i\vec{k}_{\bot}\cdot\vec{x}_{\bot}}e^{-i\beta_{n}x}\sin\left(\frac{n\pi}{L}x\right)
\]
 where $C_{2},\beta=\frac{\omega_{n}C^{2}\left(\Omega\right)A}{r^{3}}\left(\omega_{d}-\Omega\right)$
correspond to $N_{n},\beta_{n}$, respectively.

We now use $\Psi$ to denote the full field again. The quantized field
$\hat{\Psi}$ is a linear combination of all modes, with a discrete
sum in the $x$ direction and continuous integration in the $y,z$
directions: 
\[
\hat{\Psi}=\sum_{n}\int\frac{d^{2}k_{\perp}}{(2\pi)^{2}}\left(\hat{a}_{n,\vec{k}_{\perp}}u_{n}+\hat{a}^{\dagger}_{n,\vec{k}_{\perp}}u^{*}_{n}\right)
\]
 However, with the $\delta$-function normalization used here, the
mode functions satisfy $\langle u_{n}\vert u_{m}\rangle=\delta^{(2)}\left(\vec{k}_{\bot,n}-\vec{k}_{\bot,m}\right)\delta_{mn}$,
so the normalization coefficient $N_{n}$ already includes a factor
of $\frac{1}{2\pi}$. The momentum integral should therefore be $\int d^{2}k_{\perp}$,
without an additional division by $2\pi$. Thus, 
\[
\hat{\Psi}=\sum_{n}\int d^{2}k_{\perp}\left(\hat{a}_{n,\vec{k}_{\perp}}u_{n}+\hat{a}^{\dagger}_{n,\vec{k}_{\perp}}u^{*}_{n}\right)
\]

\Needspace{110pt}\par\addvspace{20pt}\noindent\begin{minipage}{\linewidth}\begingroup\fontsize{12}{15}\selectfont\setlength{\abovedisplayskip}{0pt}\setlength{\belowdisplayskip}{0pt}\setlength{\abovedisplayshortskip}{0pt}\setlength{\belowdisplayshortskip}{0pt}\phantomsection\addcontentsline{toc}{section}{IV. The Klein-Gordon Inner Product, Normalization Conditions, and Integration Constants}
\[
\begin{gathered}\text{\bfseries IV. The Klein-Gordon Inner Product,}\\
\text{\bfseries Normalization Conditions, and Integration Constants}
\end{gathered}
\]
\endgroup\end{minipage}\par\nopagebreak[4]\vspace{20pt}\nopagebreak[4]

\[
\langle u_{n},u_{m}\rangle=i\int_{S}\left[\left(\partial_{\mu}u_{n}\right)u^{*}_{m}-u_{n}\left(\partial_{\mu}u^{*}_{m}\right)\right]\sqrt{-\hat{g}_{S}}n^{\mu}dS
\]

1. The hypersurface $S$ ($dt=0$, $t=\text{constant}$): the unit
normal vector to $S$ and its dual; the induced metric matrix on $S$
and its determinant. The induced metric is given by 
\[
h_{\mu\nu}=g_{\mu\nu}-n_{\mu}n_{\nu}\hfill(n^{\mu}n_{\mu}=1)
\]

$S$ is the hypersurface $t=\text{constant}$, whose normal direction
is $n_{\mu}\propto\partial_{\mu}t=\left(1,0,0,0\right)$. Let the
future-directed unit normal to $S$ in the dual space (the normal
covector, where a covector is a dual vector) be 
\[
n_{\mu}=\left(C,0,0,0\right),
\]
 Here, $C$ is a normalization coefficient to be determined. The corresponding
normal vector in the tangent space is 
\[
n^{\mu}=\hat{g}^{\mu\nu}n_{\nu}.
\]
 We have 
\begin{align*}
n^{t} & =\hat{g}^{tt}n_{t}+\hat{g}^{tx}n_{x}+\hat{g}^{ty}n_{y}+\hat{g}^{tz}n_{z}\\
 & =\hat{g}^{tt}n_{t}+0\\
 & =\frac{A}{\Delta r^{2}}C,
\end{align*}
 {}
\[
n^{x}=\hat{g}^{xt}n_{t}=\frac{A}{\Delta r}\left(\omega_{d}-\Omega\right)C,
\]
 {}
\[
n^{y}=\hat{g}^{yt}n_{t}=0,
\]
 {}
\[
n^{z}=\hat{g}^{zt}n_{t}=0.
\]
 A unit vector has norm $1$ and therefore satisfies 
\begin{align*}
n_{\mu}n^{\mu} & =1\\
n_{t}n^{t}+n_{x}n^{x}+n_{y}n^{y}+n_{z}n^{z} & =1\\
\frac{A}{\Delta r^{2}}C^{2} & =1\\
C & =\sqrt{\frac{\Delta r^{2}}{A}}\\
 & =r\sqrt{\frac{\Delta}{A}},
\end{align*}
 This gives 
\[
n^{\mu}=r\sqrt{\frac{\Delta}{A}}\left(\frac{A}{\Delta r^{2}},\frac{A}{\Delta r}\left(\omega_{d}-\Omega\right),0,0\right).
\]
 (Both $n_{t}$ and $n^{t}$ are greater than $0$, so the normal
is future-directed.)

From 
\begin{align*}
\hat{g}_{S,\mu\nu} & =\hat{g}_{\mu\nu}-n_{\mu}n_{\nu}\\
 & =\hat{g}_{\mu\nu}-\left(r\sqrt{\frac{\Delta}{A}}\right)^{2}\delta_{\mu t}\delta_{\nu t}
\end{align*}
 we obtain 
\[
\hat{g}_{S,ij}=\hat{g}_{ij}.
\]
 Alternatively, starting from 
\[
ds^{2}=\hat{g}_{tt}dt^{2}+2\hat{g}_{tx}dtdx+\hat{g}_{yy}dy^{2}+\hat{g}_{zz}dz^{2}+\hat{g}_{xx}dx^{2},
\]
 and setting 
\[
dt=0,
\]
 gives 
\[
ds^{'2}=\hat{g}_{xx}dx^{2}+\hat{g}_{yy}dy^{2}+\hat{g}_{zz}dz^{2}.
\]
 The induced metric on this hypersurface ($t=\text{constant}$) is
\begin{align*}
\hat{g}_{S,ij} & =\hat{g}_{ij}\\
 & =\begin{pmatrix}\hat{g}_{xx} & 0 & 0\\
0 & \hat{g}_{yy} & 0\\
0 & 0 & \hat{g}_{zz}
\end{pmatrix},
\end{align*}
 and its determinant is 
\begin{align*}
\hat{g}_{S} & =\hat{g}_{xx}\cdot\hat{g}_{yy}\cdot\hat{g}_{zz}=-\frac{A}{r^{4}}\cdot\left(-\frac{r^{2}}{\Delta}\right)\cdot\left(-1\right)\\
 & =-\frac{A}{\Delta r^{2}}.
\end{align*}

2. The hypersurface $\tilde{S}$ ($d\tilde{t}=0$, $\tilde{t}=\text{constant}$):
the unit normal vector to $\tilde{S}$ and its dual; the induced metric
matrix on $\tilde{S}$ and its determinant.

Introduce a coordinate transformation to diagonalize the metric matrix.
The four-dimensional spacetime line element is 
\begin{align*}
ds^{2} & =\hat{g}_{tt}dt^{2}+2\hat{g}_{tx}dtdx+\hat{g}_{yy}dy^{2}+\hat{g}_{zz}dz^{2}+\hat{g}_{xx}dx^{2}\\
 & =\hat{g}_{tt}\left(dt^{2}+2\frac{\hat{g}_{tx}}{\hat{g}_{tt}}dtdx+\frac{\hat{g}^{2}_{tx}}{\hat{g}^{2}_{tt}}dx^{2}\right)-\hat{g}_{tt}\frac{\hat{g}^{2}_{tx}}{\hat{g}^{2}_{tt}}dx^{2}+\hat{g}_{yy}dy^{2}+\hat{g}_{zz}dz^{2}+\hat{g}_{xx}dx^{2}\\
 & =\hat{g}_{tt}\left(dt+\frac{\hat{g}_{tx}}{\hat{g}_{tt}}dx\right)^{2}+\left(\hat{g}_{xx}-\frac{\hat{g}^{2}_{tx}}{\hat{g}_{tt}}\right)dx^{2}+\hat{g}_{yy}dy^{2}+\hat{g}_{zz}dz^{2},
\end{align*}
 The coordinate transformation $d\tilde{t}=dt+\frac{\hat{g}_{tx}}{\hat{g}_{tt}}dx$
diagonalizes the metric matrix. Setting 
\[
d\tilde{t}=0,
\]
 gives the induced metric matrix on the hypersurface $\tilde{S}$
($\tilde{t}=\text{constant}$): 
\[
\hat{g}_{\tilde{S},ij}=\begin{pmatrix}\hat{g}_{xx}-\frac{\hat{g}^{2}_{tx}}{\hat{g}_{tt}} & 0 & 0\\
0 & \hat{g}_{yy} & 0\\
0 & 0 & \hat{g}_{zz}
\end{pmatrix},
\]
 Its determinant is 
\begin{align*}
\hat{g}_{\tilde{S}} & =\left(\hat{g}_{xx}-\frac{\hat{g}^{2}_{tx}}{\hat{g}_{tt}}\right)\times\hat{g}_{yy}\times\hat{g}_{zz}\\
 & =\frac{\hat{g}}{\hat{g}_{tt}}=\frac{-1}{C^{-2}\left(\Omega\right)}=-C^{2}\left(\Omega\right).
\end{align*}
 (using $\hat{g}=\det\left(\hat{g}_{\mu\nu}\right)=\left(\hat{g}_{tt}\hat{g}_{xx}-\hat{g}^{2}_{tx}\right)\times\hat{g}_{yy}\times\hat{g}_{zz}=-1$).

The use of $\hat{g}_{S}=-\frac{\hat{g}}{\hat{g}_{tt}}$ in the paper
can be explained as follows. The metric signature is $(+,-,-,-)$.
For the spatial part, $\left(-1\right)\times\left(-1\right)\times\left(-1\right)=-1$,
so the determinant of the spatial metric, $\hat{g}_{S}$, should be
$<0$. Since the quantity under the square root in the spatial volume
element $\sqrt{\hat{g}_{S}}dS$ cannot be negative, a minus sign must
be included. Thus, the spatial volume element used here is $\sqrt{-\hat{g}_{S}}dS$.

Method 2 for finding the induced metric on the hypersurface $\tilde{S}$
($\tilde{t}=\text{constant}$): the metric 
\begin{align*}
ds^{2} & =C^{-2}\left(\Omega\right)dt^{2}-2\frac{A}{r^{3}}\left(\Omega-\omega_{d}\right)dtdx-\frac{r^{2}}{\Delta}dy^{2}-dz^{2}-\frac{A}{r^{4}}dx^{2}\\
 & =\hat{g}_{tt}dt^{2}+2\hat{g}_{tx}dtdx+\hat{g}_{yy}dy^{2}+\hat{g}_{zz}dz^{2}+\hat{g}_{xx}dx^{2}
\end{align*}
 contains the cross term $dtdx$, so the time and spatial coordinate
directions are not orthogonal. Construct a direction orthogonal to
the time direction, such that the inner product of a vector in this
direction with a vector in the time direction is $0$. The basis vector
in the time direction is 
\[
\partial_{t}=\left(\partial_{t}\right)^{\mu}\partial_{\mu}
\]
 From 
\[
\partial_{t}f\left(x^{\mu}\right)=\frac{\partial f}{\partial t}=\frac{\partial f}{\partial x^{\mu}}\frac{\partial x^{\mu}}{\partial t}
\]
 we obtain 
\[
\partial_{t}=\frac{\partial x^{\mu}}{\partial t}\frac{\partial}{\partial x^{\mu}}=\frac{\partial x^{\mu}}{\partial t}\partial_{\mu}
\]
 and hence 
\[
\left(\partial_{t}\right)^{\mu}=\frac{\partial x^{\mu}}{\partial t}=(1,0,0,0)=\delta^{\mu}_{\!t}
\]
 For an infinitesimal displacement 
\[
d\boldsymbol{x}=dx^{\nu}\partial_{\nu}
\]
 require the inner product to vanish: 
\begin{align*}
\hat{g}_{\mu\nu}\left(\partial_{t}\right)^{\mu}dx^{\nu} & =\hat{g}_{\mu\nu}\delta^{\mu}_{\!t}dx^{\nu}\\
 & =\hat{g}_{t\nu}dx^{\nu}\\
 & =\hat{g}_{tt}dt+\hat{g}_{ti}dx^{i}\\
 & =0
\end{align*}
 This gives 
\[
dt=-\frac{\hat{g}_{ti}}{\hat{g}_{tt}}dx^{i}
\]
 The line element on this hypersurface $\tilde{S}$, which is orthogonal
to the time direction, is 
\begin{align*}
d\tilde{s}^{2} & =\hat{g}_{tt}dt^{2}+2\hat{g}_{tx}dtdx+\hat{g}_{yy}dy^{2}+\hat{g}_{zz}dz^{2}+\hat{g}_{xx}dx^{2}\\
 & =\hat{g}_{tt}\left(-\frac{\hat{g}_{ti}}{\hat{g}_{tt}}dx^{i}\right)^{2}+2\hat{g}_{tx}\left(-\frac{\hat{g}_{ti}}{\hat{g}_{tt}}dx^{i}\right)dx+\hat{g}_{yy}dy^{2}+\hat{g}_{zz}dz^{2}+\hat{g}_{xx}dx^{2}\\
 & =\hat{g}_{tt}\left(-\frac{\hat{g}_{tx}}{\hat{g}_{tt}}dx\right)^{2}+2\hat{g}_{tx}\left(-\frac{\hat{g}_{tx}}{\hat{g}_{tt}}dx\right)dx+\hat{g}_{yy}dy^{2}+\hat{g}_{zz}dz^{2}+\hat{g}_{xx}dx^{2}\\
 & =\frac{\hat{g}^{2}_{tx}}{\hat{g}_{tt}}dx^{2}-2\frac{\hat{g}^{2}_{tx}}{\hat{g}_{tt}}dx^{2}+\hat{g}_{yy}dy^{2}+\hat{g}_{zz}dz^{2}+\hat{g}_{xx}dx^{2}\\
 & =\left(\hat{g}_{xx}-\frac{\hat{g}^{2}_{tx}}{\hat{g}_{tt}}\right)dx^{2}+\hat{g}_{yy}dy^{2}+\hat{g}_{zz}dz^{2}
\end{align*}

We now find the unit normal vector to $\tilde{S}$ and its dual. The
metric components are treated as constants, so $d\tilde{t}=dt+\frac{\hat{g}_{tx}}{\hat{g}_{tt}}dx$
is integrable, giving $\tilde{t}=t+\frac{\hat{g}_{tx}}{\hat{g}_{tt}}x+\text{const}$.
The hypersurface $\tilde{S}$ is defined by $\tilde{t}=\text{constant}$,
and its normal direction is $n_{\mu}\propto\partial_{\mu}\tilde{t}=\left(1,\frac{\hat{g}_{tx}}{\hat{g}_{tt}},0,0\right)$.
Let the future-directed unit normal to $\tilde{S}$ in the dual space
(the normal covector, where a covector is a dual vector; see Liang
Canbin, p. 99) be 
\[
n_{\mu}=\tilde{C}\left(1,\frac{\hat{g}_{tx}}{\hat{g}_{tt}},0,0\right).
\]
 Here, $\tilde{C}$ is a normalization coefficient to be determined.
\[
n^{\mu}=\hat{g}^{\mu\nu}n_{\nu},
\]
 {}
\[
\hat{g}_{\mu\nu}=\begin{pmatrix}C^{-2}\left(\Omega\right) & -\frac{A}{r^{3}}\left(\Omega-\omega_{d}\right) & 0 & 0\\
-\frac{A}{r^{3}}\left(\Omega-\omega_{d}\right) & -\frac{A}{r^{4}} & 0 & 0\\
0 & 0 & -\frac{r^{2}}{\Delta} & 0\\
0 & 0 & 0 & -1
\end{pmatrix},
\]
 {}
\[
\hat{g}^{\mu\nu}=\begin{pmatrix}\frac{A}{\Delta r^{2}} & \frac{A}{\Delta r}\left(\omega_{d}-\Omega\right) & 0 & 0\\
\frac{A}{\Delta r}\left(\omega_{d}-\Omega\right) & -\frac{r^{2}}{\Delta}C^{-2}\left(\Omega\right) & 0 & 0\\
0 & 0 & -\frac{\Delta}{r^{2}} & 0\\
0 & 0 & 0 & -1
\end{pmatrix},
\]
 {}
\[
C\left(\Omega\right)=\left\{ \frac{\Delta r^{2}}{A}\left[1-\frac{A^{2}}{\Delta r^{4}}\left(\Omega-\omega_{d}\right)^{2}\right]\right\} ^{-1/2},
\]
 We have 
\begin{align*}
n^{t} & =\hat{g}^{tt}n_{t}+\hat{g}^{tx}n_{x}+\hat{g}^{ty}n_{y}+\hat{g}^{tz}n_{z}\\
 & =\frac{A}{\Delta r^{2}}\cdot\tilde{C}+\frac{A}{\Delta r}\left(\omega_{d}-\Omega\right)\cdot\frac{-\frac{A}{r^{3}}\left(\Omega-\omega_{d}\right)}{C^{-2}\left(\Omega\right)}\tilde{C}+0+0\\
 & =\frac{A}{\Delta r^{2}}\tilde{C}+\frac{A^{2}}{\Delta r^{4}}\left(\Omega-\omega_{d}\right)^{2}\frac{1}{\frac{\Delta r^{2}}{A}\left[1-\frac{A^{2}}{\Delta r^{4}}\left(\Omega-\omega_{d}\right)^{2}\right]}\tilde{C}\\
 & =\frac{A}{\Delta r^{2}}\tilde{C}+\frac{A}{\Delta r^{2}}\left[-1+\frac{1}{1-\frac{A^{2}}{\Delta r^{4}}\left(\Omega-\omega_{d}\right)^{2}}\right]\tilde{C}\\
 & =\frac{A}{\Delta r^{2}}\frac{1}{1-\frac{A^{2}}{\Delta r^{4}}\left(\Omega-\omega_{d}\right)^{2}}\tilde{C},
\end{align*}
 where 
\[
\frac{a}{1-a}=-1+\frac{1}{1-a},\hfill a=\frac{A^{2}}{\Delta r^{4}}\left(\Omega-\omega_{d}\right)^{2}.
\]
 {}
\begin{align*}
n^{x} & =\hat{g}^{xt}n_{t}+\hat{g}^{xx}n_{x}+\hat{g}^{xy}n_{y}+\hat{g}^{xz}n_{z}\\
 & =\frac{A}{\Delta r}\left(\omega_{d}-\Omega\right)\cdot\tilde{C}-\frac{r^{2}}{\Delta}C^{-2}\left(\Omega\right)\cdot\frac{-\frac{A}{r^{3}}\left(\Omega-\omega_{d}\right)}{C^{-2}\left(\Omega\right)}\tilde{C}+0+0\\
 & =\frac{A}{\Delta r}\left(\omega_{d}-\Omega\right)\tilde{C}+\frac{r^{2}}{\Delta}\frac{A}{r^{3}}\left(\Omega-\omega_{d}\right)\tilde{C}\\
 & =-\frac{A}{\Delta r}\left(\Omega-\omega_{d}\right)\tilde{C}+\frac{A}{\Delta r}\left(\Omega-\omega_{d}\right)\tilde{C}\\
 & =0,
\end{align*}
 {}
\[
n^{y}=n^{z}=0,\hfill\left(\hat{g}^{yt}=\hat{g}^{yt}=\hat{g}^{zt}=\hat{g}^{zt}=0,n_{y}=n_{z}=0\right).
\]
 A unit vector has norm $1$ and therefore satisfies 
\begin{align*}
n_{\mu}n^{\mu} & =1\\
n_{t}n^{t}+n_{x}n^{x}+n_{y}n^{y}+n_{z}n^{z} & =1\\
\frac{A}{\Delta r^{2}}\frac{1}{1-\frac{A^{2}}{\Delta r^{4}}\left(\Omega-\omega_{d}\right)^{2}}\tilde{C}^{2} & =1\\
\tilde{C} & =\sqrt{\frac{\Delta r^{2}}{A}\left[1-\frac{A^{2}}{\Delta r^{4}}\left(\Omega-\omega_{d}\right)^{2}\right]}=C^{-1}\left(\Omega\right),
\end{align*}
 This gives 
\begin{align*}
n^{\mu} & =\left(\frac{A}{\Delta r^{2}}\frac{1}{1-\frac{A^{2}}{\Delta r^{4}}\left(\Omega-\omega_{d}\right)^{2}}\cdot\sqrt{\frac{\Delta r^{2}}{A}\left[1-\frac{A^{2}}{\Delta r^{4}}\left(\Omega-\omega_{d}\right)^{2}\right]},0,0,0\right)\\
 & =\left(\frac{1}{\left(\sqrt{\frac{\Delta r^{2}}{A}\left[1-\frac{A^{2}}{\Delta r^{4}}\left(\Omega-\omega_{d}\right)^{2}\right]}\right)^{2}}\cdot\sqrt{\frac{\Delta r^{2}}{A}\left[1-\frac{A^{2}}{\Delta r^{4}}\left(\Omega-\omega_{d}\right)^{2}\right]},0,0,0\right)\\
 & =\left(\left\{ \frac{\Delta r^{2}}{A}\left[1-\frac{A^{2}}{\Delta r^{4}}\left(\Omega-\omega_{d}\right)^{2}\right]\right\} ^{-1/2},0,0,0\right)\\
 & =C\left(\Omega\right)\left(1,0,0,0\right).
\end{align*}
 (Both $n_{t}$ and $n^{t}$ are greater than $0$, so the normal
is future-directed.)

Summary:

1. The hypersurface $S$ ($dt=0$, $t=\text{constant}$). The unit
normal vector to $S$ and its dual are 
\begin{align*}
n^{\mu} & =r\sqrt{\frac{\Delta}{A}}\left(\frac{A}{\Delta r^{2}},\frac{A}{\Delta r}\left(\omega_{d}-\Omega\right),0,0\right),\\
n_{\mu} & =\left(r\sqrt{\frac{\Delta}{A}},0,0,0\right).
\end{align*}
 The induced metric matrix on $S$ and its determinant are 
\begin{align*}
\hat{g}_{S,ij} & =\hat{g}_{ij}\\
 & =\begin{pmatrix}\hat{g}_{xx} & 0 & 0\\
0 & \hat{g}_{yy} & 0\\
0 & 0 & \hat{g}_{zz}
\end{pmatrix},\\
\hat{g}_{S} & =-\frac{A}{\Delta r^{2}}.
\end{align*}

2. The hypersurface $\tilde{S}$ ($d\tilde{t}=0$, $\tilde{t}=\text{constant}$).
The unit normal vector to $\tilde{S}$ and its dual are 
\begin{align*}
n^{\mu} & =C\left(\Omega\right)\left(1,0,0,0\right),\\
n_{\mu} & =C^{-1}\left(\Omega\right)\left(1,\frac{\hat{g}_{tx}}{\hat{g}_{tt}},0,0\right).
\end{align*}
 The induced metric matrix on $\tilde{S}$ and its determinant are
\begin{align*}
\hat{g}_{\tilde{S},ij} & =\begin{pmatrix}\hat{g}_{xx}-\frac{\hat{g}^{2}_{tx}}{\hat{g}_{tt}} & 0 & 0\\
0 & \hat{g}_{yy} & 0\\
0 & 0 & \hat{g}_{zz}
\end{pmatrix},\\
\hat{g}_{\tilde{S}} & =-C^{2}\left(\Omega\right).
\end{align*}
 where $C\left(\Omega\right)=\left\{ \frac{\Delta r^{2}}{A}\left[1-\frac{A^{2}}{\Delta r^{4}}\left(\Omega-\omega_{d}\right)^{2}\right]\right\} ^{-1/2},\hat{g}_{tx}=-\frac{A}{r^{3}}\left(\Omega-\omega_{d}\right),\hat{g}_{tt}=C^{-2}\left(\Omega\right)$.

The quantum field modes satisfy the normalization condition 
\[
\langle u_{n},u_{m}\rangle=i\int_{S}\left[\left(\partial_{\mu}u_{n}\right)u^{*}_{m}-u_{n}\left(\partial_{\mu}u^{*}_{m}\right)\right]\sqrt{-\hat{g}_{S}}n^{\mu}dS=\delta^{(2)}\left(\vec{k}_{\bot,n}-\vec{k}_{\bot,m}\right)\delta_{mn}.
\]

1. The hypersurface $S$ ($dt=0$, $t=\text{constant}$): 
\begin{align*}
 & \langle u_{n},u_{m}\rangle=i\int_{S}\biggl[\biggl(\partial_{\mu}u_{n}\biggr)u^{*}_{m}-u_{n}\biggl(\partial_{\mu}u^{*}_{m}\biggr)\biggr]\sqrt{-\hat{g}_{S}}n^{\mu}dS\\
 & =i\int_{S}\biggl[\biggl(\partial_{t}u_{n}\biggr)u^{*}_{m}-u_{n}\biggl(\partial_{t}u^{*}_{m}\biggr)\biggr]\sqrt{-\hat{g}_{S}}n^{t}dS+i\int_{S}\biggl[\biggl(\partial_{x}u_{n}\biggr)u^{*}_{m}-u_{n}\biggl(\partial_{x}u^{*}_{m}\biggr)\biggr]\sqrt{-\hat{g}_{S}}n^{x}dS\\
 & =i\int_{S}\biggl[\biggl(-i\omega_{n}u_{n}\biggr)u^{*}_{m}-u_{n}\biggl(i\omega_{m}u^{*}_{m}\biggr)\biggr]\sqrt{-\hat{g}_{S}}n^{t}dS+i\int_{S}\biggl[\biggl(\partial_{x}u_{n}\biggr)u^{*}_{m}-u_{n}\biggl(\partial_{x}u^{*}_{m}\biggr)\biggr]\sqrt{-\hat{g}_{S}}n^{x}dS\\
 & =i\int_{S}\biggl[-i\biggl(\omega_{n}+\omega_{m}\biggr)u_{n}u^{*}_{m}\biggr]\sqrt{-\hat{g}_{S}}n^{t}dS+i\int_{S}\biggl[\biggl(-i\beta_{n}u_{n}+\frac{n\pi}{L}\cos\biggl(\frac{n\pi}{L}x\biggr)F_{n}\biggr)u^{*}_{m}\\
 & \qquad-u_{n}\biggl(i\beta_{m}u^{*}_{m}+\frac{m\pi}{L}\cos\biggl(\frac{m\pi}{L}x\biggr)F^{*}_{m}\biggr)\biggr]\sqrt{-\hat{g}_{S}}n^{x}dS\\
 & =i\int_{S}\biggl[-i\biggl(\omega_{n}+\omega_{m}\biggr)u_{n}u^{*}_{m}\biggr]\sqrt{-\hat{g}_{S}}n^{t}dS+i\int_{S}\biggl[\biggl(-i\beta_{n}u_{n}u^{*}_{m}+\frac{n\pi}{L}\cos\biggl(\frac{n\pi}{L}x\biggr)F_{n}u^{*}_{m}\\
 & \qquad\biggr)-\biggl(i\beta_{m}u_{n}u^{*}_{m}+\frac{m\pi}{L}\cos\biggl(\frac{m\pi}{L}x\biggr)F^{*}_{m}u_{n}\biggr)\biggr]\sqrt{-\hat{g}_{S}}n^{x}dS\\
 & =i\int_{S}\biggl[-i\biggl(\omega_{n}+\omega_{m}\biggr)u_{n}u^{*}_{m}\biggr]\sqrt{-\hat{g}_{S}}n^{t}dS+i\int_{S}\biggl[-i\biggl(\beta_{n}+\beta_{m}\biggr)u_{n}u^{*}_{m}\biggr]\sqrt{-\hat{g}_{S}}n^{x}dS\\
 & \qquad+i\int_{S}\biggl[\frac{n\pi}{L}\cos\biggl(\frac{n\pi}{L}x\biggr)F_{n}u^{*}_{m}-\frac{m\pi}{L}\cos\biggl(\frac{m\pi}{L}x\biggr)F^{*}_{m}u_{n}\biggr]\sqrt{-\hat{g}_{S}}n^{x}dS
\end{align*}
 where 
\begin{align*}
\partial_{x}u_{n} & =\partial_{x}\left[N_{n}e^{-i\omega_{n}t}e^{i\vec{k}_{\bot,n}\cdot\vec{x}_{\bot}}e^{-i\beta_{n}x}\sin\left(\frac{n\pi}{L}x\right)\right]\\
 & =-i\beta_{n}u_{n}+\frac{n\pi}{L}\cos\left(\frac{n\pi}{L}x\right)N_{n}e^{-i\omega_{n}t}e^{i\vec{k}_{\bot,n}\cdot\vec{x}_{\bot}}e^{-i\beta_{n}x}\\
 & =-i\beta_{n}u_{n}+\frac{n\pi}{L}\cos\left(\frac{n\pi}{L}x\right)F_{n}
\end{align*}
 {}
\begin{align*}
\partial_{x}u^{*}_{m} & =\partial_{x}\left[N_{m}e^{i\omega_{m}t}e^{-i\vec{k}_{\bot,m}\cdot\vec{x}_{\bot}}e^{i\beta_{m}x}\sin\left(\frac{m\pi}{L}x\right)\right]\\
 & =i\beta_{m}u^{*}_{m}+\frac{m\pi}{L}\cos\left(\frac{m\pi}{L}x\right)F^{*}_{m}
\end{align*}
 and 
\begin{align*}
 & i\int_{S}\biggl[-i\biggl(\omega_{n}+\omega_{m}\biggr)u_{n}u^{*}_{m}\biggr]\sqrt{-\hat{g}_{S}}n^{t}dS=i\int_{S}\biggl\{-i\biggl(\omega_{n}+\omega_{m}\biggr)\biggl[N_{n}e^{-i\omega_{n}t}e^{i\vec{k}_{\bot,n}\cdot\vec{x}_{\bot}}\\
 & \qquad e^{-i\beta_{n}x}\sin\biggl(\frac{n\pi}{L}x\biggr)\biggr]\biggl[N_{m}e^{i\omega_{m}t}e^{-i\vec{k}_{\bot,m}\cdot\vec{x}_{\bot}}e^{i\beta_{m}x}\sin\biggl(\frac{m\pi}{L}x\biggr)\biggr]\biggr\}\sqrt{-\hat{g}_{S}}n^{t}dxdydz\\
 & =i\int_{S}\biggl\{-i\biggl(\omega_{n}+\omega_{m}\biggr)\biggl[N_{n}N_{m}e^{-i\biggl(\omega_{n}-\omega_{m}\biggr)t}e^{i\biggl(\vec{k}_{\bot,n}-\vec{k}_{\bot,m}\biggr)\cdot\vec{x}_{\bot}}e^{-i\biggl(\beta_{n}-\beta_{m}\biggr)x}\sin\biggl(\frac{n\pi}{L}x\biggr)\\
 & \qquad\sin\biggl(\frac{m\pi}{L}x\biggr)\biggr]\biggr\}\sqrt{-\hat{g}_{S}}n^{t}dxdydz\\
 & =\biggl(\omega_{n}+\omega_{m}\biggr)\sqrt{-\hat{g}_{S}}n^{t}N_{n}N_{m}e^{-i\biggl(\omega_{n}-\omega_{m}\biggr)t}\int^{L}_{0}e^{-i\biggl(\beta_{n}-\beta_{m}\biggr)x}\sin\biggl(\frac{n\pi}{L}x\biggr)\sin\biggl(\frac{m\pi}{L}x\biggr)\\
 & \qquad dx\int^{\infty}_{-\infty}\int^{\infty}_{-\infty}e^{i\biggl(\vec{k}_{\bot,n}-\vec{k}_{\bot,m}\biggr)\cdot\vec{x}_{\bot}}dydz
\end{align*}
 Similarly, 
\begin{align*}
 & i\int_{S}\biggl[-i\biggl(\beta_{n}+\beta_{m}\biggr)u_{n}u^{*}_{m}\biggr]\sqrt{-\hat{g}_{S}}n^{x}dS=\biggl(\beta_{n}+\beta_{m}\biggr)\sqrt{-\hat{g}_{S}}n^{x}N_{n}N_{m}e^{-i\biggl(\omega_{n}-\omega_{m}\biggr)t}\\
 & \qquad\int^{L}_{0}e^{-i\biggl(\beta_{n}-\beta_{m}\biggr)x}\sin\biggl(\frac{n\pi}{L}x\biggr)\sin\biggl(\frac{m\pi}{L}x\biggr)dx\int^{\infty}_{-\infty}\int^{\infty}_{-\infty}e^{i\biggl(\vec{k}_{\bot,n}-\vec{k}_{\bot,m}\biggr)\cdot\vec{x}_{\bot}}dydz
\end{align*}
 Also, 
\[
\int^{\infty}_{-\infty}\int^{\infty}_{-\infty}e^{i\left(\vec{k}_{\bot,n}-\vec{k}_{\bot,m}\right)\cdot\vec{x}_{\bot}}dydz=\left(2\pi\right)^{2}\delta^{(2)}\left(\vec{k}_{\bot,n}-\vec{k}_{\bot,m}\right)
\]
 Imposing $\langle u_{n},u_{m}\rangle=\delta^{(2)}\left(\vec{k}_{\bot,n}-\vec{k}_{\bot,m}\right)\delta_{mn}$
gives 
\[
2\omega_{n}\sqrt{-\hat{g}_{S}}n^{t}N^{2}_{n}\frac{L}{2}\left(2\pi\right)^{2}+2\beta_{n}\sqrt{-\hat{g}_{S}}n^{x}N^{2}_{n}\frac{L}{2}\left(2\pi\right)^{2}=1
\]
 (For $n=m$, and only in this case, $\delta_{mn}=1$; when $n\neq m$,
we have $\delta_{mn}=0$ and correspondingly $\langle u_{n},u_{m}\rangle=0$.
Also, $\omega_{n}+\omega_{m}\rightarrow2\omega_{n}$; $e^{-i\left(\omega_{n}-\omega_{m}\right)t}\rightarrow e^{-i\cdot0\cdot t}=1$;
$\int^{L}_{0}e^{-i\left(\beta_{n}-\beta_{m}\right)x}\sin\left(\frac{n\pi}{L}x\right)\sin\left(\frac{m\pi}{L}x\right)dx\rightarrow\int^{L}_{0}\sin^{2}\left(\frac{n\pi}{L}x\right)dx=L/2$.)
We obtain 
\begin{align*}
N^{2}_{n} & =\frac{1}{2\omega_{n}\sqrt{-\hat{g}_{S}}n^{t}\frac{L}{2}\left(2\pi\right)^{2}+2\beta_{n}\sqrt{-\hat{g}_{S}}n^{x}\frac{L}{2}\left(2\pi\right)^{2}}\\
 & =\frac{1}{\left(2\pi\right)^{2}L\sqrt{-\hat{g}_{S}}\left(\omega_{n}n^{t}+\beta_{n}n^{x}\right)}\\
 & =\frac{1}{\left(2\pi\right)^{2}L\sqrt{\frac{A}{\Delta r^{2}}}\left\{ \omega_{n}r\sqrt{\frac{\Delta}{A}}\frac{A}{\Delta r^{2}}+\left[\frac{\omega_{n}C^{2}\left(\Omega\right)A}{r^{3}}\left(\omega_{d}-\Omega\right)\right]r\sqrt{\frac{\Delta}{A}}\frac{A}{\Delta r}\left(\omega_{d}-\Omega\right)\right\} }\\
 & =\frac{1}{\left(2\pi\right)^{2}L\sqrt{\frac{A}{\Delta r^{2}}}\omega_{n}\left\{ \sqrt{\frac{A}{\Delta}}\frac{1}{r}+\frac{C^{2}\left(\Omega\right)A}{r^{3}}\sqrt{\frac{A}{\Delta}}\left(\omega_{d}-\Omega\right)^{2}\right\} }\\
 & =\frac{1}{\left(2\pi\right)^{2}L\sqrt{\frac{A}{\Delta r^{2}}}\omega_{n}\left\{ \sqrt{\frac{A}{\Delta}}\frac{1}{r}+\frac{\left\{ \frac{\Delta r^{2}}{A}\left[1-\frac{A^{2}}{\Delta r^{4}}\left(\Omega-\omega_{d}\right)^{2}\right]\right\} ^{-1}A}{r^{3}}\sqrt{\frac{A}{\Delta}}\left(\omega_{d}-\Omega\right)^{2}\right\} }\\
 & =\frac{1}{\left(2\pi\right)^{2}L\sqrt{\frac{A}{\Delta r^{2}}}\omega_{n}\frac{1}{r}\sqrt{\frac{A}{\Delta}}\left\{ 1+\frac{1}{\left[1-\frac{A^{2}}{\Delta r^{4}}\left(\Omega-\omega_{d}\right)^{2}\right]}\frac{A^{2}}{\Delta r^{4}}\left(\omega_{d}-\Omega\right)^{2}\right\} }\\
 & =\frac{1}{\left(2\pi\right)^{2}L\sqrt{\frac{A}{\Delta r^{2}}}\omega_{n}\frac{1}{r}\sqrt{\frac{A}{\Delta}}\left\{ 1+\frac{1}{\frac{A}{\Delta r^{2}}C^{-2}\left(\Omega\right)}\left[1-\frac{A}{\Delta r^{2}}C^{-2}\left(\Omega\right)\right]\right\} }\\
 & =\frac{1}{\left(2\pi\right)^{2}L\sqrt{\frac{A}{\Delta r^{2}}}\omega_{n}\frac{1}{r}\sqrt{\frac{A}{\Delta}}\left\{ 1+\frac{\Delta r^{2}}{A}C^{2}\left(\Omega\right)\left[1-\frac{A}{\Delta r^{2}}C^{-2}\left(\Omega\right)\right]\right\} }\\
 & =\frac{1}{\left(2\pi\right)^{2}L\sqrt{\frac{A}{\Delta r^{2}}}\omega_{n}\frac{1}{r}\sqrt{\frac{A}{\Delta}}\left[\frac{\Delta r^{2}}{A}C^{2}\left(\Omega\right)\right]}\\
 & =\frac{1}{\left(2\pi\right)^{2}L\frac{A}{\Delta}\frac{1}{r}\omega_{n}\frac{1}{r}\left[\frac{\Delta r^{2}}{A}C^{2}\left(\Omega\right)\right]}\\
 & =\frac{1}{\left(2\pi\right)^{2}L\omega_{n}C^{2}\left(\Omega\right)}
\end{align*}
 or equivalently, 
\[
N^{2}_{n}=\frac{1}{\left(2\pi\right)^{2}L\omega_{n}C^{2}\left(\Omega\right)}
\]
 Here, 
\[
C\left(\Omega\right)=\left\{ \frac{\Delta r^{2}}{A}\left[1-\frac{A^{2}}{\Delta r^{4}}\left(\Omega-\omega_{d}\right)^{2}\right]\right\} ^{-1/2}
\]
 {}
\[
C^{-2}\left(\Omega\right)=\frac{\Delta r^{2}}{A}\left[1-\frac{A^{2}}{\Delta r^{4}}\left(\Omega-\omega_{d}\right)^{2}\right]
\]
 which give 
\[
1-\frac{A^{2}}{\Delta r^{4}}\left(\Omega-\omega_{d}\right)^{2}=\frac{A}{\Delta r^{2}}C^{-2}\left(\Omega\right)
\]
 {}
\[
\frac{A^{2}}{\Delta r^{4}}\left(\Omega-\omega_{d}\right)^{2}=1-\frac{A}{\Delta r^{2}}C^{-2}\left(\Omega\right)
\]
 We also have 
\[
\beta_{n}=\frac{b_{n}A}{r^{3}}\left(\omega_{d}-\Omega\right)=\frac{\omega_{n}C^{2}\left(\Omega\right)A}{r^{3}}\left(\omega_{d}-\Omega\right)
\]
 and 
\begin{align*}
\int^{L}_{0}e^{-i\left(\beta_{n}-\beta_{n}\right)x}\sin\left(\frac{n\pi}{L}x\right)\sin\left(\frac{n\pi}{L}x\right)dx & =\int^{L}_{0}\sin^{2}\left(\frac{n\pi}{L}x\right)dx\\
 & =\int^{L}_{0}\frac{1-\cos\left(2\frac{n\pi}{L}x\right)}{2}dx\\
 & =\frac{L}{2}\hfill,\text{where }\sin2n\pi=0
\end{align*}

\[
F_{n}u^{*}_{n}=F^{*}_{n}u_{n}
\]
\[
\frac{n\pi}{L}\cos\left(\frac{n\pi}{L}x\right)F_{n}u^{*}_{n}-\frac{n\pi}{L}\cos\left(\frac{n\pi}{L}x\right)F^{*}_{n}u_{n}=0
\]

Note that the original paper gives 
\begin{align*}
N^{2}_{n} & =\frac{1}{\left(2\pi\right)^{2}L\omega_{n}}\frac{1}{r}\sqrt{\frac{A}{\Delta}}C^{-3}\left(\Omega\right)\\
 & =\frac{1}{\left(2\pi\right)^{2}L\omega_{n}C^{2}\left(\Omega\right)}\frac{1}{r}\sqrt{\frac{A}{\Delta}}C^{-1}\left(\Omega\right)\\
 & =\frac{1}{\left(2\pi\right)^{2}L\omega_{n}C^{2}\left(\Omega\right)}\frac{1}{r}\sqrt{\frac{A}{\Delta}}\left\{ \frac{\Delta r^{2}}{A}\left[1-\frac{A^{2}}{\Delta r^{4}}\left(\Omega-\omega_{d}\right)^{2}\right]\right\} ^{1/2}\\
 & =\frac{1}{\left(2\pi\right)^{2}L\omega_{n}C^{2}\left(\Omega\right)}\frac{1}{r}\sqrt{\frac{A}{\Delta}}\sqrt{\frac{\Delta}{A}}r\left[1-\frac{A^{2}}{\Delta r^{4}}\left(\Omega-\omega_{d}\right)^{2}\right]^{1/2}\\
 & =\frac{1}{\left(2\pi\right)^{2}L\omega_{n}C^{2}\left(\Omega\right)}\left[1-\frac{A^{2}}{\Delta r^{4}}\left(\Omega-\omega_{d}\right)^{2}\right]^{1/2}
\end{align*}
 where 
\[
C\left(\Omega\right)=\left\{ \frac{\Delta r^{2}}{A}\left[1-\frac{A^{2}}{\Delta r^{4}}\left(\Omega-\omega_{d}\right)^{2}\right]\right\} ^{-1/2}
\]

2. The hypersurface $\tilde{S}$ ($d\tilde{t}=0$, $\tilde{t}=\text{constant}$):
\[
n^{\mu}=C\left(\Omega\right)\left(1,0,0,0\right),
\]
 {}
\begin{align*}
\langle u_{n},u_{m}\rangle & =i\int_{\tilde{S}}\left[\left(\partial_{\mu}u_{n}\right)u^{*}_{m}-u_{n}\left(\partial_{\mu}u^{*}_{m}\right)\right]\sqrt{-\hat{g}_{\tilde{S}}}n^{\mu}d\tilde{S}\\
 & =i\int_{\tilde{S}}\left[\left(\partial_{t}u_{n}\right)u^{*}_{m}-u_{n}\left(\partial_{t}u^{*}_{m}\right)\right]\sqrt{-\hat{g}_{\tilde{S}}}n^{t}d\tilde{S}\\
 & =i\int_{\tilde{S}}\left[\left(-i\omega_{n}u_{n}\right)u^{*}_{m}-u_{n}\left(i\omega_{m}u^{*}_{m}\right)\right]\sqrt{-\hat{g}_{\tilde{S}}}n^{t}d\tilde{S}\\
 & =i\int_{\tilde{S}}\left[-i\left(\omega_{n}+\omega_{m}\right)u_{n}u^{*}_{m}\right]\sqrt{-\hat{g}_{\tilde{S}}}n^{t}d\tilde{S}
\end{align*}
 and 
\begin{align*}
 & i\int_{\tilde{S}}\biggl[-i\biggl(\omega_{n}+\omega_{m}\biggr)u_{n}u^{*}_{m}\biggr]\sqrt{-\hat{g}_{\tilde{S}}}n^{t}d\tilde{S}=i\int_{\tilde{S}}\biggl\{-i\biggl(\omega_{n}+\omega_{m}\biggr)\biggl[N_{n}e^{-i\omega_{n}t}e^{i\vec{k}_{\bot,n}\cdot\vec{x}_{\bot}}\\
 & \qquad e^{-i\beta_{n}x}\sin\biggl(\frac{n\pi}{L}x\biggr)\biggr]\biggl[N_{m}e^{i\omega_{m}t}e^{-i\vec{k}_{\bot,m}\cdot\vec{x}_{\bot}}e^{i\beta_{m}x}\sin\biggl(\frac{m\pi}{L}x\biggr)\biggr]\biggr\}\sqrt{-\hat{g}_{\tilde{S}}}n^{t}dxdydz\\
 & =i\int_{\tilde{S}}\biggl\{-i\biggl(\omega_{n}+\omega_{m}\biggr)\biggl[N_{n}N_{m}e^{-i\biggl(\omega_{n}-\omega_{m}\biggr)t}e^{i\biggl(\vec{k}_{\bot,n}-\vec{k}_{\bot,m}\biggr)\cdot\vec{x}_{\bot}}e^{-i\biggl(\beta_{n}-\beta_{m}\biggr)x}\sin\biggl(\frac{n\pi}{L}x\biggr)\\
 & \qquad\sin\biggl(\frac{m\pi}{L}x\biggr)\biggr]\biggr\}\sqrt{-\hat{g}_{\tilde{S}}}n^{t}dxdydz\\
 & =\biggl(\omega_{n}+\omega_{m}\biggr)\sqrt{-\hat{g}_{\tilde{S}}}n^{t}N_{n}N_{m}\int^{L}_{0}e^{-i\biggl(\omega_{n}-\omega_{m}\biggr)t(x)}e^{-i\biggl(\beta_{n}-\beta_{m}\biggr)x}\sin\biggl(\frac{n\pi}{L}x\biggr)\sin\biggl(\frac{m\pi}{L}x\biggr)\\
 & \qquad dx\int^{\infty}_{-\infty}\int^{\infty}_{-\infty}e^{i\biggl(\vec{k}_{\bot,n}-\vec{k}_{\bot,m}\biggr)\cdot\vec{x}_{\bot}}dydz
\end{align*}
 where $t(x)=\tilde{t}-\frac{\hat{g}_{tx}}{\hat{g}_{tt}}x-\text{const}$.
Also, 
\[
\int^{\infty}_{-\infty}\int^{\infty}_{-\infty}e^{i\left(\vec{k}_{\bot,n}-\vec{k}_{\bot,m}\right)\cdot\vec{x}_{\bot}}dydz=\left(2\pi\right)^{2}\delta^{(2)}\left(\vec{k}_{\bot,n}-\vec{k}_{\bot,m}\right)
\]
 Imposing $\langle u_{n},u_{m}\rangle=\delta^{(2)}\left(\vec{k}_{\bot,n}-\vec{k}_{\bot,m}\right)\delta_{mn}$
gives 
\[
2\omega_{n}\sqrt{-\hat{g}_{\tilde{S}}}n^{t}N^{2}_{n}\frac{L}{2}\left(2\pi\right)^{2}=1
\]
 (For $n=m$, and only in this case, $\delta_{mn}=1$; when $n\neq m$,
we have $\delta_{mn}=0$ and correspondingly $\langle u_{n},u_{m}\rangle=0$.
Also, $\omega_{n}+\omega_{m}\rightarrow2\omega_{n}$; $e^{-i\left(\omega_{n}-\omega_{m}\right)t(x)}\rightarrow e^{-i\cdot0\cdot t(x)}=1$;
$\int^{L}_{0}e^{-i\left(\beta_{n}-\beta_{m}\right)x}\sin\left(\frac{n\pi}{L}x\right)\sin\left(\frac{m\pi}{L}x\right)dx\rightarrow\int^{L}_{0}\sin^{2}\left(\frac{n\pi}{L}x\right)dx=L/2$.)
This gives 
\begin{align*}
N^{2}_{n} & =\frac{1}{2\omega_{n}\sqrt{-\hat{g}_{\tilde{S}}}n^{t}\frac{L}{2}\left(2\pi\right)^{2}}\\
 & =\frac{1}{\left(2\pi\right)^{2}L\omega_{n}\sqrt{-\hat{g}_{\tilde{S}}}n^{t}}\\
 & =\frac{1}{\left(2\pi\right)^{2}L\omega_{n}C\left(\Omega\right)\cdot C\left(\Omega\right)}\\
 & =\frac{1}{\left(2\pi\right)^{2}L\omega_{n}C^{2}\left(\Omega\right)}
\end{align*}

\Needspace{110pt}\par\addvspace{20pt}\noindent\begin{minipage}{\linewidth}\begingroup\fontsize{12}{15}\selectfont\setlength{\abovedisplayskip}{0pt}\setlength{\belowdisplayskip}{0pt}\setlength{\abovedisplayshortskip}{0pt}\setlength{\belowdisplayshortskip}{0pt}\phantomsection\addcontentsline{toc}{section}{V. Mean Vacuum Energy Density}
\[
\text{\bfseries V. Mean Vacuum Energy Density}
\]
\endgroup\end{minipage}\par\nopagebreak[4]\vspace{20pt}\nopagebreak[4]

\begin{align*}
\langle\epsilon_{\text{vac}}\rangle & =\frac{1}{V_{p}}\int_{V}\sqrt{-\hat{g}_{S}}dV\langle0\vert T_{\mu\nu}\hat{w}^{\mu}\hat{w}^{\nu}\vert0\rangle\\
 & =C^{2}\left(\Omega\right)\frac{1}{V_{p}}\int_{V}\sqrt{-\hat{g}_{S}}dV\langle0\vert T_{tt}\vert0\rangle\\
 & =C^{2}\left(\Omega\right)\frac{1}{\sqrt{-\hat{g}_{S}}\int^{L}_{0}dx\int^{+\infty}_{-\infty}\int^{+\infty}_{-\infty}dydz}\sqrt{-\hat{g}_{S}}\int_{V}dV\langle0\vert T_{tt}\vert0\rangle\\
 & =C^{2}\left(\Omega\right)\frac{1}{L\int^{+\infty}_{-\infty}\int^{+\infty}_{-\infty}dydz}\int_{V}dV\langle0\vert T_{tt}\vert0\rangle\\
 & =C^{2}\left(\Omega\right)\frac{1}{L\int^{+\infty}_{-\infty}\int^{+\infty}_{-\infty}dydz}\int^{L}_{0}dx\int^{+\infty}_{-\infty}\int^{+\infty}_{-\infty}dydz\langle0\vert T_{tt}\vert0\rangle
\end{align*}
 where 
\[
T_{\mu\nu}\hat{w}^{\mu}\hat{w}^{\nu}=T_{tt}\hat{w}^{t}\hat{w}^{t}=C^{2}\left(\Omega\right)T_{tt}
\]
 {}
\[
\hat{w}^{\mu}=C\left(\Omega\right)\left(1,0,0,0\right)
\]
 {}
\[
V_{p}=\int_{V}\sqrt{-\hat{g}_{S}}dV=\sqrt{-\hat{g}_{S}}\int^{L}_{0}dx\int^{+\infty}_{-\infty}\int^{+\infty}_{-\infty}dydz
\]

We have 
\begin{align*}
T_{\mu\nu} & =2\frac{\partial\mathcal{L}}{\partial g^{\mu\nu}}-g_{\mu\nu}\mathcal{L}\\
 & =2\frac{\partial\left(\frac{1}{2}g^{\alpha\beta}\partial_{\alpha}\Psi\partial_{\beta}\Psi\right)}{\partial g^{\mu\nu}}-g_{\mu\nu}\left(\frac{1}{2}g^{\alpha\beta}\partial_{\alpha}\Psi\partial_{\beta}\Psi\right)\\
 & =\frac{\partial g^{\alpha\beta}}{\partial g^{\mu\nu}}\partial_{\alpha}\Psi\partial_{\beta}\Psi-\frac{1}{2}g_{\mu\nu}\left(g^{\alpha\beta}\partial_{\alpha}\Psi\partial_{\beta}\Psi\right)\\
 & =\frac{1}{2}\left(\delta^{\alpha}_{\!\mu}\delta^{\beta}_{\!\nu}+\delta^{\alpha}_{\!\nu}\delta^{\beta}_{\!\mu}\right)\partial_{\alpha}\Psi\partial_{\beta}\Psi-\frac{1}{2}g_{\mu\nu}\left(g^{\alpha\beta}\partial_{\alpha}\Psi\partial_{\beta}\Psi\right)\\
 & =\frac{1}{2}\left(\partial_{\mu}\Psi\partial_{\nu}\Psi+\partial_{\nu}\Psi\partial_{\mu}\Psi\right)-\frac{1}{2}g_{\mu\nu}\left(g^{\alpha\beta}\partial_{\alpha}\Psi\partial_{\beta}\Psi\right)\\
 & =\partial_{\mu}\Psi\partial_{\nu}\Psi-\frac{1}{2}g_{\mu\nu}\left(g^{\alpha\beta}\partial_{\alpha}\Psi\partial_{\beta}\Psi\right)
\end{align*}
 where, for a massless scalar field in vacuum in curved spacetime,
the Lagrangian density is 
\[
\mathcal{L}=\frac{1}{2}g^{\alpha\beta}\partial_{\alpha}\Psi\partial_{\beta}\Psi
\]

(With the metric signature $(-,+,+,+)$, the kinetic term in the Lagrangian
density must be defined as 
\[
-\frac{1}{2}g^{\mu\nu}\partial_{\mu}\Psi\partial_{\nu}\Psi
\]
 to obtain positive kinetic energy. To ensure that the energy density
measured by a local observer is positive, that is, $T_{\mu\nu}U^{\mu}U^{\nu}>0$,
the energy-momentum tensor must be defined as $T_{\mu\nu}=-\frac{2}{\sqrt{-g}}\frac{\delta S_{m}}{\delta g^{\mu\nu}}$,
where $S_{m}=\int d^{4}x\sqrt{-g}\mathcal{L}_{m}$ is the matter action.
This gives $T_{\mu\nu}=-2\frac{\partial\mathcal{L}_{m}}{\partial g^{\mu\nu}}+g_{\mu\nu}\mathcal{L}_{m}$.
With the metric signature $(+,-,-,-)$, the kinetic term in the Lagrangian
density must instead be defined as $\frac{1}{2}g^{\mu\nu}\partial_{\mu}\Psi\partial_{\nu}\Psi$
to obtain positive kinetic energy. To ensure that the energy density
measured by a local observer is positive, that is, $T_{\mu\nu}U^{\mu}U^{\nu}>0$,
the energy-momentum tensor must be defined as $T_{\mu\nu}=\frac{2}{\sqrt{-g}}\frac{\delta S_{m}}{\delta g^{\mu\nu}}$,
giving $T_{\mu\nu}=2\frac{\partial\mathcal{L}_{m}}{\partial g^{\mu\nu}}-g_{\mu\nu}\mathcal{L}_{m}$.)
We also use 
\[
\ensuremath{\frac{\partial g^{\alpha\beta}}{\partial g^{\mu\nu}}=\frac{1}{2}\left(\delta^{\alpha}_{\!\mu}\delta^{\beta}_{\!\nu}+\delta^{\alpha}_{\!\nu}\delta^{\beta}_{\!\mu}\right)}
\]

Then 
\begin{align*}
 & \hat{T}_{tt}=\partial_{t}\hat{\Psi}\partial_{t}\hat{\Psi}-\frac{1}{2}\hat{g}_{tt}\biggl(\hat{g}^{\alpha\beta}\partial_{\alpha}\hat{\Psi}\partial_{\beta}\hat{\Psi}\biggr)\\
 & =\partial_{t}\sum_{n}\int d^{2}k_{\perp,n}\biggl(\hat{a}_{n,\boldsymbol{k}_{\perp,n}}u_{n}+\hat{a}^{\dagger}_{n,\boldsymbol{k}_{\perp,n}}u^{*}_{n}\biggr)\partial_{t}\sum_{m}\int d^{2}k_{\perp,m}\biggl(\hat{a}_{m,\boldsymbol{k}_{\perp,m}}u_{m}+\hat{a}^{\dagger}_{m,\boldsymbol{k}_{\perp,m}}u^{*}_{m}\biggr)\\
 & \qquad-\frac{1}{2}\hat{g}_{tt}\biggl[\hat{g}^{\mu\nu}\partial_{\mu}\sum_{n}\int d^{2}k_{\perp,n}\biggl(\hat{a}_{n,\boldsymbol{k}_{\perp,n}}u_{n}+\hat{a}^{\dagger}_{n,\boldsymbol{k}_{\perp,n}}u^{*}_{n}\biggr)\partial_{\nu}\sum_{m}\int d^{2}k_{\perp,m}\biggl(\hat{a}_{m,\boldsymbol{k}_{\perp,m}}u_{m}\\
 & \qquad+\hat{a}^{\dagger}_{m,\boldsymbol{k}_{\perp,m}}u^{*}_{m}\biggr)\biggr]\\
 & =\sum_{n}\int d^{2}k_{\perp,n}\biggl(\hat{a}_{n,\boldsymbol{k}_{\perp,n}}\partial_{t}u_{n}+\hat{a}^{\dagger}_{n,\boldsymbol{k}_{\perp,n}}\partial_{t}u^{*}_{n}\biggr)\sum_{m}\int d^{2}k_{\perp,m}\biggl(\hat{a}_{m,\boldsymbol{k}_{\perp,m}}\partial_{t}u_{m}+\hat{a}^{\dagger}_{m,\boldsymbol{k}_{\perp,m}}\partial_{t}\\
 & \qquad u^{*}_{m}\biggr)-\frac{1}{2}\hat{g}_{tt}\biggl[\hat{g}^{\mu\nu}\sum_{n}\int d^{2}k_{\perp,n}\biggl(\hat{a}_{n,\boldsymbol{k}_{\perp,n}}\partial_{\mu}u_{n}+\hat{a}^{\dagger}_{n,\boldsymbol{k}_{\perp,n}}\partial_{\mu}u^{*}_{n}\biggr)\sum_{m}\int d^{2}k_{\perp,m}\biggl(\hat{a}_{m,\boldsymbol{k}_{\perp,m}}\partial_{\nu}\\
 & \qquad u_{m}+\hat{a}^{\dagger}_{m,\boldsymbol{k}_{\perp,m}}\partial_{\nu}u^{*}_{m}\biggr)\biggr]\\
 & =\sum_{n,m}\int d^{2}k_{\perp,n}\int d^{2}k_{\perp,m}\biggl(\hat{a}_{n,\boldsymbol{k}_{\perp,n}}\partial_{t}u_{n}\cdot\hat{a}_{m,\boldsymbol{k}_{\perp,m}}\partial_{t}u_{m}+\hat{a}_{n,\boldsymbol{k}_{\perp,n}}\partial_{t}u_{n}\cdot\hat{a}^{\dagger}_{m,\boldsymbol{k}_{\perp,m}}\partial_{t}u^{*}_{m}\\
 & \qquad+\hat{a}^{\dagger}_{n,\boldsymbol{k}_{\perp,n}}\partial_{t}u^{*}_{n}\cdot\hat{a}_{m,\boldsymbol{k}_{\perp,m}}\partial_{t}u_{m}+\hat{a}^{\dagger}_{n,\boldsymbol{k}_{\perp,n}}\partial_{t}u^{*}_{n}\cdot\hat{a}^{\dagger}_{m,\boldsymbol{k}_{\perp,m}}\partial_{t}u^{*}_{m}\biggr)-\frac{1}{2}\hat{g}_{tt}\biggl[\hat{g}^{\mu\nu}\sum_{n,m}\int d^{2}k_{\perp,n}\\
 & \qquad\int d^{2}k_{\perp,m}\biggl(\hat{a}_{n,\boldsymbol{k}_{\perp,n}}\partial_{\mu}u_{n}\cdot\hat{a}_{m,\boldsymbol{k}_{\perp,m}}\partial_{\nu}u_{m}+\hat{a}_{n,\boldsymbol{k}_{\perp,n}}\partial_{\mu}u_{n}\cdot\hat{a}^{\dagger}_{m,\boldsymbol{k}_{\perp,m}}\partial_{\nu}u^{*}_{m}+\hat{a}^{\dagger}_{n,\boldsymbol{k}_{\perp,n}}\partial_{\mu}u^{*}_{n}\\
 & \qquad\cdot\hat{a}_{m,\boldsymbol{k}_{\perp,m}}\partial_{\nu}u_{m}+\hat{a}^{\dagger}_{n,\boldsymbol{k}_{\perp,n}}\partial_{\mu}u^{*}_{n}\cdot\hat{a}^{\dagger}_{m,\boldsymbol{k}_{\perp,m}}\partial_{\nu}u^{*}_{m}\biggr)\biggr]
\end{align*}
 where 
\[
\hat{\Psi}=\sum_{n}\int d^{2}k_{\perp}\left(\hat{a}_{n,\boldsymbol{k}_{\perp}}u_{n}+\hat{a}^{\dagger}_{n,\boldsymbol{k}_{\perp}}u^{*}_{n}\right)
\]
 {}
\begin{align*}
 & \langle0\vert\hat{T}_{tt}\vert0\rangle=\sum_{n,m}\int d^{2}k_{\perp,n}\int d^{2}k_{\perp,m}\partial_{t}u_{n}\partial_{t}u^{*}_{m}\langle0\vert\hat{a}_{n,\boldsymbol{k}_{\perp,n}}\hat{a}^{\dagger}_{m,\boldsymbol{k}_{\perp,m}}\vert0\rangle-\frac{1}{2}\hat{g}_{tt}\biggl[\hat{g}^{\mu\nu}\sum_{n,m}\int d^{2}\\
 & \qquad k_{\perp,n}\int d^{2}k_{\perp,m}\partial_{\mu}u_{n}\partial_{\nu}u^{*}_{m}\langle0\vert\hat{a}_{n,\boldsymbol{k}_{\perp,n}}\hat{a}^{\dagger}_{m,\boldsymbol{k}_{\perp,m}}\vert0\rangle\biggr]\\
 & =\sum_{n,m}\int d^{2}k_{\perp,n}\int d^{2}k_{\perp,m}\partial_{t}u_{n}\partial_{t}u^{*}_{m}\delta^{(2)}\biggl(\boldsymbol{k}_{\perp,n}-\boldsymbol{k}_{\perp,m}\biggr)\delta_{nm}-\frac{1}{2}\hat{g}_{tt}\biggl[\hat{g}^{\mu\nu}\sum_{n,m}\int d^{2}k_{\perp,n}\\
 & \qquad\int d^{2}k_{\perp,m}\partial_{\mu}u_{n}\partial_{\nu}u^{*}_{m}\delta^{(2)}\biggl(\boldsymbol{k}_{\perp,n}-\boldsymbol{k}_{\perp,m}\biggr)\delta_{nm}\biggr]\\
 & =\sum_{n}\int d^{2}k_{\perp,n}\partial_{t}u_{n}\partial_{t}u^{*}_{n}-\frac{1}{2}\hat{g}_{tt}\biggl(\hat{g}^{\mu\nu}\sum_{n}\int d^{2}k_{\perp,n}\partial_{\mu}u_{n}\partial_{\nu}u^{*}_{n}\biggr)\\
 & =\sum_{n}\int d^{2}k_{\perp,n}\biggl[\partial_{t}u_{n}\partial_{t}u^{*}_{n}-\frac{1}{2}\hat{g}_{tt}\biggl(\hat{g}^{\mu\nu}\partial_{\mu}u_{n}\partial_{\nu}u^{*}_{n}\biggr)\biggr]
\end{align*}
 Here, 
\[
\langle0\vert\hat{a}_{n,\boldsymbol{k}_{\perp,n}}\partial_{t}u_{n}\cdot\hat{a}_{m,\boldsymbol{k}_{\perp,m}}\partial_{t}u_{m}\vert0\rangle=\partial_{t}u_{n}\partial_{t}u_{m}\langle0\vert\hat{a}_{n,\boldsymbol{k}_{\perp,n}}\hat{a}_{m,\boldsymbol{k}_{\perp,m}}\vert0\rangle=0
\]
 {}
\begin{align*}
\langle0\vert\hat{a}_{n,\boldsymbol{k}_{\perp,n}}\partial_{t}u_{n}\cdot\hat{a}^{\dagger}_{m,\boldsymbol{k}_{\perp,m}}\partial_{t}u^{*}_{m}\vert0\rangle & =\partial_{t}u_{n}\partial_{t}u^{*}_{m}\langle0\vert\hat{a}_{n,\boldsymbol{k}_{\perp,n}}\hat{a}^{\dagger}_{m,\boldsymbol{k}_{\perp,m}}\vert0\rangle\\
 & =\partial_{t}u_{n}\partial_{t}u^{*}_{m}\langle0\vert\delta^{(2)}\left(\boldsymbol{k}_{\perp,n}-\boldsymbol{k}_{\perp,m}\right)\delta_{nm}+\hat{a}^{\dagger}_{m,\boldsymbol{k}_{\perp,m}}\hat{a}_{n,\boldsymbol{k}_{\perp,n}}\vert0\rangle\\
 & =\partial_{t}u_{n}\partial_{t}u^{*}_{m}\langle0\vert\delta^{(2)}\left(\boldsymbol{k}_{\perp,n}-\boldsymbol{k}_{\perp,m}\right)\delta_{nm}\vert0\rangle+\langle0\vert\hat{a}^{\dagger}_{m,\boldsymbol{k}_{\perp,m}}\hat{a}_{n,\boldsymbol{k}_{\perp,n}}\vert0\rangle\\
 & =\partial_{t}u_{n}\partial_{t}u^{*}_{m}\delta^{(2)}\left(\boldsymbol{k}_{\perp,n}-\boldsymbol{k}_{\perp,m}\right)\delta_{nm}\langle0\vert0\rangle+0\\
 & =\partial_{t}u_{n}\partial_{t}u^{*}_{m}\delta^{(2)}\left(\boldsymbol{k}_{\perp,n}-\boldsymbol{k}_{\perp,m}\right)\delta_{nm}
\end{align*}
 {}
\[
\langle0\vert\hat{a}^{\dagger}_{n,\boldsymbol{k}_{\perp,n}}\partial_{t}u^{*}_{n}\cdot\hat{a}_{m,\boldsymbol{k}_{\perp,m}}\partial_{t}u_{m}\vert0\rangle=\partial_{t}u^{*}_{n}\partial_{t}u_{m}\langle0\vert\hat{a}^{\dagger}_{n,\boldsymbol{k}_{\perp,n}}\hat{a}_{m,\boldsymbol{k}_{\perp,m}}\vert0\rangle=0
\]
 {}
\[
\langle0\vert\hat{a}^{\dagger}_{n,\boldsymbol{k}_{\perp,n}}\partial_{t}u^{*}_{n}\cdot\hat{a}^{\dagger}_{m,\boldsymbol{k}_{\perp,m}}\partial_{t}u^{*}_{m}\vert0\rangle=\partial_{t}u^{*}_{n}\partial_{t}u^{*}_{m}\langle0\vert\hat{a}^{\dagger}_{n,\boldsymbol{k}_{\perp,n}}\hat{a}^{\dagger}_{m,\boldsymbol{k}_{\perp,m}}\vert0\rangle=0
\]
 These results follow from 
\[
\hat{a}_{m,\boldsymbol{k}_{\perp,m}}\vert0\rangle=0
\]
 (the definition of the vacuum state) and 
\[
\langle0\vert\hat{a}^{\dagger}_{n,\boldsymbol{k}_{\perp,n}}\hat{a}^{\dagger}_{m,\boldsymbol{k}_{\perp,m}}\vert0\rangle=\langle0\vert1_{n,\boldsymbol{k}_{\perp,n}},1_{m,\boldsymbol{k}_{\perp,m}}\rangle=0
\]
 (the two states are orthogonal), or equivalently, 
\[
\langle0\vert\hat{a}^{\dagger}_{n,\boldsymbol{k}_{\perp,n}}\hat{a}^{\dagger}_{m,\boldsymbol{k}_{\perp,m}}\vert0\rangle=\left(\hat{a}_{m,\boldsymbol{k}_{\perp,m}}\hat{a}_{n,\boldsymbol{k}_{\perp,n}}\vert0\rangle\right)^{\dagger}\vert0\rangle=0^{\dagger}\cdot\vert0\rangle=0
\]
 From 
\[
\left[\hat{a}_{n,\boldsymbol{k}_{\perp,n}},\hat{a}^{\dagger}_{m,\boldsymbol{k}_{\perp,m}}\right]=\hat{a}_{n,\boldsymbol{k}_{\perp,n}}\hat{a}^{\dagger}_{m,\boldsymbol{k}_{\perp,m}}-\hat{a}^{\dagger}_{m,\boldsymbol{k}_{\perp,m}}\hat{a}_{n,\boldsymbol{k}_{\perp,n}}=\delta^{(2)}\left(\boldsymbol{k}_{\perp,n}-\boldsymbol{k}_{\perp,m}\right)\delta_{nm}
\]
 we obtain 
\[
\hat{a}_{n,\boldsymbol{k}_{\perp,n}}\hat{a}^{\dagger}_{m,\boldsymbol{k}_{\perp,m}}=\delta^{(2)}\left(\boldsymbol{k}_{\perp,n}-\boldsymbol{k}_{\perp,m}\right)\delta_{nm}+\hat{a}^{\dagger}_{m,\boldsymbol{k}_{\perp,m}}\hat{a}_{n,\boldsymbol{k}_{\perp,n}}
\]

We previously obtained 
\[
\langle0\vert\hat{T}_{tt}\vert0\rangle=\sum_{n}\int d^{2}k_{\perp}\left[\partial_{t}u_{n}\partial_{t}u^{*}_{n}-\frac{1}{2}\hat{g}_{tt}\left(\hat{g}^{\mu\nu}\partial_{\mu}u_{n}\partial_{\nu}u^{*}_{n}\right)\right]
\]
 where $\int d^{2}k_{\perp,n}$ is denoted by $\int d^{2}k_{\perp}$.
We also have 
\[
u_{n}=N_{n}e^{-i\omega_{n}t}e^{i\boldsymbol{k}_{\bot}\cdot\boldsymbol{r}_{\bot}}e^{-i\beta_{n}x}\sin\left(\frac{n\pi}{L}x\right)
\]
 {}
\[
e^{i\boldsymbol{k}_{\bot}\cdot\boldsymbol{r}_{\bot}}=e^{i\left(k_{y}y+k_{z}z\right)}
\]
 {}
\[
\hat{g}_{\mu\nu}=\begin{pmatrix}C^{-2}\left(\Omega\right) & -\frac{A}{r^{3}}\left(\Omega-\omega_{d}\right) & 0 & 0\\
-\frac{A}{r^{3}}\left(\Omega-\omega_{d}\right) & -\frac{A}{r^{4}} & 0 & 0\\
0 & 0 & -\frac{r^{2}}{\Delta} & 0\\
0 & 0 & 0 & -1
\end{pmatrix},
\]
 {}
\[
\hat{g}^{\mu\nu}=\begin{pmatrix}\frac{A}{\Delta r^{2}} & \frac{A}{\Delta r}\left(\omega_{d}-\Omega\right) & 0 & 0\\
\frac{A}{\Delta r}\left(\omega_{d}-\Omega\right) & -\frac{r^{2}}{\Delta}C^{-2}\left(\Omega\right) & 0 & 0\\
0 & 0 & -\frac{\Delta}{r^{2}} & 0\\
0 & 0 & 0 & -1
\end{pmatrix},
\]
 {}
\[
C\left(\Omega\right)=\left\{ \frac{\Delta r^{2}}{A}\left[1-\frac{A^{2}}{\Delta r^{4}}\left(\Omega-\omega_{d}\right)^{2}\right]\right\} ^{-1/2},
\]
 {}
\[
N^{2}_{n}=\frac{1}{\left(2\pi\right)^{2}L\omega_{n}C^{2}\left(\Omega\right)}
\]
 {}
\[
\omega_{n}=\frac{r}{\sqrt{\Delta}}C^{-2}\left(\Omega\right)\sqrt{\left(\frac{n\pi}{L}\right)^{2}+\frac{\Delta}{r^{2}}C^{2}\left(\Omega\right)\left(\frac{\Delta}{r^{2}}k^{2}_{y}+k^{2}_{z}\right)}
\]
 {}
\[
\beta_{n}=\frac{b_{n}A}{r^{3}}\left(\omega_{d}-\Omega\right)
\]
 {}
\[
b_{n}=\omega_{n}C^{2}\left(\Omega\right)
\]
 Thus, 
\begin{align*}
\partial_{t}u_{n}\partial_{t}u^{*}_{n} & =\partial_{t}N_{n}e^{-i\omega_{n}t}e^{i\boldsymbol{k}_{\bot}\cdot\boldsymbol{r}_{\bot}}e^{-i\beta_{n}x}\sin\left(\frac{n\pi}{L}x\right)\partial_{t}N_{n}e^{i\omega_{n}t}e^{-i\boldsymbol{k}_{\bot}\cdot\boldsymbol{r}_{\bot}}e^{i\beta_{n}x}\sin\left(\frac{n\pi}{L}x\right)\\
 & =\left(-i\omega_{n}\right)i\omega_{n}N_{n}e^{-i\omega_{n}t}e^{i\boldsymbol{k}_{\bot}\cdot\boldsymbol{r}_{\bot}}e^{-i\beta_{n}x}\sin\left(\frac{n\pi}{L}x\right)N_{n}e^{i\omega_{n}t}e^{-i\boldsymbol{k}_{\bot}\cdot\boldsymbol{r}_{\bot}}e^{i\beta_{n}x}\sin\left(\frac{n\pi}{L}x\right)\\
 & =\omega^{2}_{n}N^{2}_{n}\sin^{2}\left(\frac{n\pi}{L}x\right)
\end{align*}
 {}
\begin{align*}
 & \partial_{t}u_{n}\partial_{x}u^{*}_{n}=\partial_{t}N_{n}e^{-i\omega_{n}t}e^{i\boldsymbol{k}_{\bot}\cdot\boldsymbol{r}_{\bot}}e^{-i\beta_{n}x}\sin\biggl(\frac{n\pi}{L}x\biggr)\partial_{x}N_{n}e^{i\omega_{n}t}e^{-i\boldsymbol{k}_{\bot}\cdot\boldsymbol{r}_{\bot}}e^{i\beta_{n}x}\sin\biggl(\frac{n\pi}{L}x\biggr)\\
 & =\biggl(-i\omega_{n}\biggr)N_{n}e^{-i\omega_{n}t}e^{i\boldsymbol{k}_{\bot}\cdot\boldsymbol{r}_{\bot}}e^{-i\beta_{n}x}\sin\biggl(\frac{n\pi}{L}x\biggr)N_{n}e^{i\omega_{n}t}e^{-i\boldsymbol{k}_{\bot}\cdot\boldsymbol{r}_{\bot}}\biggl[i\beta_{n}e^{i\beta_{n}x}\sin\biggl(\frac{n\pi}{L}x\biggr)\\
 & \qquad+e^{i\beta_{n}x}\biggl(\frac{n\pi}{L}\biggr)\cos\biggl(\frac{n\pi}{L}x\biggr)\biggr]\\
 & =\biggl(-i\omega_{n}\biggr)i\beta_{n}N_{n}e^{-i\omega_{n}t}e^{i\boldsymbol{k}_{\bot}\cdot\boldsymbol{r}_{\bot}}e^{-i\beta_{n}x}\sin\biggl(\frac{n\pi}{L}x\biggr)N_{n}e^{i\omega_{n}t}e^{-i\boldsymbol{k}_{\bot}\cdot\boldsymbol{r}_{\bot}}e^{i\beta_{n}x}\sin\biggl(\frac{n\pi}{L}x\biggr)+\biggl(\\
 & \qquad-i\omega_{n}\biggr)N_{n}e^{-i\omega_{n}t}e^{i\boldsymbol{k}_{\bot}\cdot\boldsymbol{r}_{\bot}}e^{-i\beta_{n}x}\sin\biggl(\frac{n\pi}{L}x\biggr)N_{n}e^{i\omega_{n}t}e^{-i\boldsymbol{k}_{\bot}\cdot\boldsymbol{r}_{\bot}}e^{i\beta_{n}x}\biggl(\frac{n\pi}{L}\biggr)\cos\biggl(\frac{n\pi}{L}x\biggr)\\
 & =\omega_{n}\beta_{n}N^{2}_{n}\sin^{2}\biggl(\frac{n\pi}{L}x\biggr)+\biggl(-i\omega_{n}\biggr)\frac{n\pi}{L}N^{2}_{n}\sin\biggl(\frac{n\pi}{L}x\biggr)\cos\biggl(\frac{n\pi}{L}x\biggr)
\end{align*}
 {}
\begin{align*}
 & \partial_{x}u_{n}\partial_{t}u^{*}_{n}=\partial_{x}N_{n}e^{-i\omega_{n}t}e^{i\boldsymbol{k}_{\bot}\cdot\boldsymbol{r}_{\bot}}e^{-i\beta_{n}x}\sin\biggl(\frac{n\pi}{L}x\biggr)\partial_{t}N_{n}e^{i\omega_{n}t}e^{-i\boldsymbol{k}_{\bot}\cdot\boldsymbol{r}_{\bot}}e^{i\beta_{n}x}\sin\biggl(\frac{n\pi}{L}x\biggr)\\
 & =N_{n}e^{-i\omega_{n}t}e^{i\boldsymbol{k}_{\bot}\cdot\boldsymbol{r}_{\bot}}\biggl[\biggl(-i\beta_{n}\biggr)e^{-i\beta_{n}x}\sin\biggl(\frac{n\pi}{L}x\biggr)+e^{-i\beta_{n}x}\biggl(\frac{n\pi}{L}\biggr)\cos\biggl(\frac{n\pi}{L}x\biggr)\biggr]i\omega_{n}N_{n}e^{i\omega_{n}t}\\
 & \qquad e^{-i\boldsymbol{k}_{\bot}\cdot\boldsymbol{r}_{\bot}}e^{i\beta_{n}x}\sin\biggl(\frac{n\pi}{L}x\biggr)\\
 & =\biggl(-i\beta_{n}\biggr)i\omega_{n}N_{n}e^{-i\omega_{n}t}e^{i\boldsymbol{k}_{\bot}\cdot\boldsymbol{r}_{\bot}}e^{-i\beta_{n}x}\sin\biggl(\frac{n\pi}{L}x\biggr)N_{n}e^{i\omega_{n}t}e^{-i\boldsymbol{k}_{\bot}\cdot\boldsymbol{r}_{\bot}}e^{i\beta_{n}x}\sin\biggl(\frac{n\pi}{L}x\biggr)+i\omega_{n}\\
 & \qquad N_{n}e^{-i\omega_{n}t}e^{i\boldsymbol{k}_{\bot}\cdot\boldsymbol{r}_{\bot}}e^{-i\beta_{n}x}\biggl(\frac{n\pi}{L}\biggr)\cos\biggl(\frac{n\pi}{L}x\biggr)N_{n}e^{i\omega_{n}t}e^{-i\boldsymbol{k}_{\bot}\cdot\boldsymbol{r}_{\bot}}e^{i\beta_{n}x}\sin\biggl(\frac{n\pi}{L}x\biggr)\\
 & =\beta_{n}\omega_{n}N^{2}_{n}\sin^{2}\biggl(\frac{n\pi}{L}x\biggr)+i\omega_{n}\frac{n\pi}{L}N^{2}_{n}\cos\biggl(\frac{n\pi}{L}x\biggr)\sin\biggl(\frac{n\pi}{L}x\biggr)
\end{align*}
 {}
\begin{align*}
 & \partial_{x}u_{n}\partial_{x}u^{*}_{n}=\partial_{x}N_{n}e^{-i\omega_{n}t}e^{i\boldsymbol{k}_{\bot}\cdot\boldsymbol{r}_{\bot}}e^{-i\beta_{n}x}\sin\biggl(\frac{n\pi}{L}x\biggr)\partial_{x}N_{n}e^{i\omega_{n}t}e^{-i\boldsymbol{k}_{\bot}\cdot\boldsymbol{r}_{\bot}}e^{i\beta_{n}x}\sin\biggl(\frac{n\pi}{L}x\biggr)\\
 & =N_{n}e^{-i\omega_{n}t}e^{i\boldsymbol{k}_{\bot}\cdot\boldsymbol{r}_{\bot}}\biggl[\biggl(-i\beta_{n}\biggr)e^{-i\beta_{n}x}\sin\biggl(\frac{n\pi}{L}x\biggr)+e^{-i\beta_{n}x}\biggl(\frac{n\pi}{L}\biggr)\cos\biggl(\frac{n\pi}{L}x\biggr)\biggr]N_{n}e^{i\omega_{n}t}\\
 & \qquad e^{-i\boldsymbol{k}_{\bot}\cdot\boldsymbol{r}_{\bot}}\biggl[i\beta_{n}e^{i\beta_{n}x}\sin\biggl(\frac{n\pi}{L}x\biggr)+e^{i\beta_{n}x}\biggl(\frac{n\pi}{L}\biggr)\cos\biggl(\frac{n\pi}{L}x\biggr)\biggr]\\
 & =N^{2}_{n}\biggl[\biggl(-i\beta_{n}\biggr)e^{-i\beta_{n}x}\sin\biggl(\frac{n\pi}{L}x\biggr)i\beta_{n}e^{i\beta_{n}x}\sin\biggl(\frac{n\pi}{L}x\biggr)+\biggl(-i\beta_{n}\biggr)e^{-i\beta_{n}x}\sin\biggl(\frac{n\pi}{L}x\biggr)e^{i\beta_{n}x}\\
 & \qquad\biggl(\frac{n\pi}{L}\biggr)\cos\biggl(\frac{n\pi}{L}x\biggr)+e^{-i\beta_{n}x}\biggl(\frac{n\pi}{L}\biggr)\cos\biggl(\frac{n\pi}{L}x\biggr)i\beta_{n}e^{i\beta_{n}x}\sin\biggl(\frac{n\pi}{L}x\biggr)+e^{-i\beta_{n}x}\biggl(\frac{n\pi}{L}\biggr)\cos\biggl(\frac{n\pi}{L}\\
 & \qquad x\biggr)e^{i\beta_{n}x}\biggl(\frac{n\pi}{L}\biggr)\cos\biggl(\frac{n\pi}{L}x\biggr)\biggr]\\
 & =N^{2}_{n}\biggl[\beta^{2}_{n}\sin^{2}\biggl(\frac{n\pi}{L}x\biggr)+\biggl(-i\beta_{n}\biggr)\frac{n\pi}{L}\sin\biggl(\frac{n\pi}{L}x\biggr)\cos\biggl(\frac{n\pi}{L}x\biggr)+i\beta_{n}\frac{n\pi}{L}\cos\biggl(\frac{n\pi}{L}x\biggr)\sin\biggl(\frac{n\pi}{L}\\
 & \qquad x\biggr)+\biggl(\frac{n\pi}{L}\biggr)^{2}\cos^{2}\biggl(\frac{n\pi}{L}x\biggr)\biggr]\\
 & =N^{2}_{n}\biggl[\beta^{2}_{n}\sin^{2}\biggl(\frac{n\pi}{L}x\biggr)+\biggl(\frac{n\pi}{L}\biggr)^{2}\cos^{2}\biggl(\frac{n\pi}{L}x\biggr)\biggr]
\end{align*}
 {}
\begin{align*}
\partial_{y}u_{n}\partial_{y}u^{*}_{n} & =\partial_{y}N_{n}e^{-i\omega_{n}t}e^{i\left(k_{y}y+k_{z}z\right)}e^{-i\beta_{n}x}\sin\left(\frac{n\pi}{L}x\right)\partial_{y}N_{n}e^{i\omega_{n}t}e^{-i\left(k_{y}y+k_{z}z\right)}e^{i\beta_{n}x}\sin\left(\frac{n\pi}{L}x\right)\\
 & =ik_{y}\left(-ik_{y}\right)N_{n}e^{-i\omega_{n}t}e^{i\left(k_{y}y+k_{z}z\right)}e^{-i\beta_{n}x}\sin\left(\frac{n\pi}{L}x\right)N_{n}e^{i\omega_{n}t}e^{-i\left(k_{y}y+k_{z}z\right)}e^{i\beta_{n}x}\sin\left(\frac{n\pi}{L}x\right)\\
 & =k^{2}_{y}N^{2}_{n}\sin^{2}\left(\frac{n\pi}{L}x\right)
\end{align*}
 {}
\begin{align*}
\partial_{z}u_{n}\partial_{z}u^{*}_{n} & =\partial_{z}N_{n}e^{-i\omega_{n}t}e^{i\left(k_{y}y+k_{z}z\right)}e^{-i\beta_{n}x}\sin\left(\frac{n\pi}{L}x\right)\partial_{z}N_{n}e^{i\omega_{n}t}e^{-i\left(k_{y}y+k_{z}z\right)}e^{i\beta_{n}x}\sin\left(\frac{n\pi}{L}x\right)\\
 & =ik_{z}\left(-ik_{z}\right)N_{n}e^{-i\omega_{n}t}e^{i\left(k_{y}y+k_{z}z\right)}e^{-i\beta_{n}x}\sin\left(\frac{n\pi}{L}x\right)N_{n}e^{i\omega_{n}t}e^{-i\left(k_{y}y+k_{z}z\right)}e^{i\beta_{n}x}\sin\left(\frac{n\pi}{L}x\right)\\
 & =k^{2}_{z}N^{2}_{n}\sin^{2}\left(\frac{n\pi}{L}x\right)
\end{align*}
 {}
\begin{align*}
 & \frac{1}{2}\hat{g}_{tt}\biggl(\hat{g}^{\mu\nu}\partial_{\mu}u_{n}\partial_{\nu}u^{*}_{n}\biggr)=\frac{1}{2}\hat{g}_{tt}\biggl(\hat{g}^{tt}\partial_{t}u_{n}\partial_{t}u^{*}_{n}+\hat{g}^{tx}\partial_{t}u_{n}\partial_{x}u^{*}_{n}+\hat{g}^{xt}\partial_{x}u_{n}\partial_{t}u^{*}_{n}\\
 & \qquad+\hat{g}^{xx}\partial_{x}u_{n}\partial_{x}u^{*}_{n}+\hat{g}^{yy}\partial_{y}u_{n}\partial_{y}u^{*}_{n}+\hat{g}^{zz}\partial_{z}u_{n}\partial_{z}u^{*}_{n}\biggr)\\
 & =\frac{1}{2}\hat{g}_{tt}\biggl\{\hat{g}^{tt}\omega^{2}_{n}N^{2}_{n}\sin^{2}\biggl(\frac{n\pi}{L}x\biggr)+\hat{g}^{tx}\biggl[\omega_{n}\beta_{n}N^{2}_{n}\sin^{2}\biggl(\frac{n\pi}{L}x\biggr)+\biggl(-i\omega_{n}\biggr)\frac{n\pi}{L}N^{2}_{n}\sin\biggl(\frac{n\pi}{L}\\
 & \qquad x\biggr)\cos\biggl(\frac{n\pi}{L}x\biggr)\biggr]+\hat{g}^{xt}\biggl[\beta_{n}\omega_{n}N^{2}_{n}\sin^{2}\biggl(\frac{n\pi}{L}x\biggr)+i\omega_{n}\frac{n\pi}{L}N^{2}_{n}\cos\biggl(\frac{n\pi}{L}x\biggr)\sin\biggl(\frac{n\pi}{L}x\biggr)\biggr]+\hat{g}^{xx}\\
 & \qquad N^{2}_{n}\biggl[\beta^{2}_{n}\sin^{2}\biggl(\frac{n\pi}{L}x\biggr)+\biggl(\frac{n\pi}{L}\biggr)^{2}\cos^{2}\biggl(\frac{n\pi}{L}x\biggr)\biggr]+\hat{g}^{yy}k^{2}_{y}N^{2}_{n}\sin^{2}\biggl(\frac{n\pi}{L}x\biggr)+\hat{g}^{zz}k^{2}_{z}N^{2}_{n}\sin^{2}\\
 & \qquad\biggl(\frac{n\pi}{L}x\biggr)\biggr\}\\
 & =\frac{1}{2}\hat{g}_{tt}\biggl\{\hat{g}^{tt}\omega^{2}_{n}N^{2}_{n}\sin^{2}\biggl(\frac{n\pi}{L}x\biggr)+\hat{g}^{tx}\omega_{n}\beta_{n}N^{2}_{n}\sin^{2}\biggl(\frac{n\pi}{L}x\biggr)+\hat{g}^{xt}\beta_{n}\omega_{n}N^{2}_{n}\sin^{2}\biggl(\frac{n\pi}{L}x\biggr)\\
 & \qquad+\hat{g}^{xx}N^{2}_{n}\biggl[\beta^{2}_{n}\sin^{2}\biggl(\frac{n\pi}{L}x\biggr)+\biggl(\frac{n\pi}{L}\biggr)^{2}\cos^{2}\biggl(\frac{n\pi}{L}x\biggr)\biggr]+\hat{g}^{yy}k^{2}_{y}N^{2}_{n}\sin^{2}\biggl(\frac{n\pi}{L}x\biggr)+\hat{g}^{zz}k^{2}_{z}N^{2}_{n}\\
 & \qquad\sin^{2}\biggl(\frac{n\pi}{L}x\biggr)\biggr\}\\
 & =\frac{1}{2}\hat{g}_{tt}\biggl\{\hat{g}^{tt}\omega^{2}_{n}N^{2}_{n}\sin^{2}\biggl(\frac{n\pi}{L}x\biggr)+2\hat{g}^{tx}\omega_{n}\beta_{n}N^{2}_{n}\sin^{2}\biggl(\frac{n\pi}{L}x\biggr)+\hat{g}^{xx}N^{2}_{n}\biggl[\beta^{2}_{n}\sin^{2}\biggl(\frac{n\pi}{L}x\biggr)\\
 & \qquad+\biggl(\frac{n\pi}{L}\biggr)^{2}\cos^{2}\biggl(\frac{n\pi}{L}x\biggr)\biggr]+\hat{g}^{yy}k^{2}_{y}N^{2}_{n}\sin^{2}\biggl(\frac{n\pi}{L}x\biggr)+\hat{g}^{zz}k^{2}_{z}N^{2}_{n}\sin^{2}\biggl(\frac{n\pi}{L}x\biggr)\biggr\}
\end{align*}
 Hence, 
\begin{align*}
 & \partial_{t}u_{n}\partial_{t}u^{*}_{n}-\frac{1}{2}\hat{g}_{tt}\biggl(\hat{g}^{\mu\nu}\partial_{\mu}u_{n}\partial_{\nu}u^{*}_{n}\biggr)=\omega^{2}_{n}N^{2}_{n}\sin^{2}\biggl(\frac{n\pi}{L}x\biggr)-\frac{1}{2}\hat{g}_{tt}\biggl\{\hat{g}^{tt}\omega^{2}_{n}N^{2}_{n}\sin^{2}\biggl(\frac{n\pi}{L}\\
 & \qquad x\biggr)+2\hat{g}^{tx}\omega_{n}\beta_{n}N^{2}_{n}\sin^{2}\biggl(\frac{n\pi}{L}x\biggr)+\hat{g}^{xx}N^{2}_{n}\biggl[\beta^{2}_{n}\sin^{2}\biggl(\frac{n\pi}{L}x\biggr)+\biggl(\frac{n\pi}{L}\biggr)^{2}\cos^{2}\biggl(\frac{n\pi}{L}x\biggr)\biggr]+\hat{g}^{yy}\\
 & \qquad k^{2}_{y}N^{2}_{n}\sin^{2}\biggl(\frac{n\pi}{L}x\biggr)+\hat{g}^{zz}k^{2}_{z}N^{2}_{n}\sin^{2}\biggl(\frac{n\pi}{L}x\biggr)\biggr\}
\end{align*}
 Write $k_{n}=\frac{n\pi}{L}$, $S_{x}=\sin^{2}\left(k_{n}x\right)$,
and $C_{x}=\cos^{2}\left(k_{n}x\right)$. Then 
\begin{align*}
 & \partial_{t}u_{n}\partial_{t}u^{*}_{n}-\frac{1}{2}\hat{g}_{tt}\biggl(\hat{g}^{\mu\nu}\partial_{\mu}u_{n}\partial_{\nu}u^{*}_{n}\biggr)=\omega^{2}_{n}N^{2}_{n}S_{x}-\frac{1}{2}\hat{g}_{tt}\biggl[\hat{g}^{tt}\omega^{2}_{n}N^{2}_{n}S_{x}+2\hat{g}^{tx}\omega_{n}\beta_{n}N^{2}_{n}\\
 & \qquad S_{x}+\hat{g}^{xx}N^{2}_{n}\biggl(\beta^{2}_{n}S_{x}+k^{2}_{n}C_{x}\biggr)+\hat{g}^{yy}k^{2}_{y}N^{2}_{n}S_{x}+\hat{g}^{zz}k^{2}_{z}N^{2}_{n}S_{x}\biggr]\\
 & =\omega^{2}_{n}N^{2}_{n}S_{x}-\frac{1}{2}\hat{g}_{tt}\biggl[\biggl(\hat{g}^{tt}\omega^{2}_{n}+2\hat{g}^{tx}\omega_{n}\beta_{n}+\hat{g}^{xx}\beta^{2}_{n}\biggr)N^{2}_{n}S_{x}+\hat{g}^{xx}N^{2}_{n}k^{2}_{n}C_{x}+\hat{g}^{yy}k^{2}_{y}N^{2}_{n}S_{x}+\hat{g}^{zz}k^{2}_{z}N^{2}_{n}S_{x}\biggr]\\
 & =\omega^{2}_{n}N^{2}_{n}S_{x}-\frac{1}{2}\hat{g}_{tt}\biggl[\omega^{2}_{n}C^{2}\biggl(\Omega\biggr)N^{2}_{n}S_{x}+\hat{g}^{xx}N^{2}_{n}k^{2}_{n}C_{x}+\hat{g}^{yy}k^{2}_{y}N^{2}_{n}S_{x}+\hat{g}^{zz}k^{2}_{z}N^{2}_{n}S_{x}\biggr]\\
 & =\omega^{2}_{n}N^{2}_{n}S_{x}-\frac{1}{2}C^{-2}\biggl(\Omega\biggr)\biggl\{\omega^{2}_{n}C^{2}\biggl(\Omega\biggr)N^{2}_{n}S_{x}+\biggl[-\frac{r^{2}}{\Delta}C^{-2}\biggl(\Omega\biggr)\biggr]N^{2}_{n}k^{2}_{n}C_{x}+\biggl(-\frac{\Delta}{r^{2}}\biggr)\\
 & \qquad k^{2}_{y}N^{2}_{n}S_{x}+\biggl(-1\biggr)k^{2}_{z}N^{2}_{n}S_{x}\biggr\}\\
 & =\omega^{2}_{n}N^{2}_{n}S_{x}-\frac{1}{2}N^{2}_{n}\biggl\{\omega^{2}_{n}S_{x}-\frac{r^{2}}{\Delta}C^{-4}\biggl(\Omega\biggr)k^{2}_{n}C_{x}-C^{-2}\biggl(\Omega\biggr)\frac{\Delta}{r^{2}}k^{2}_{y}S_{x}-C^{-2}\biggl(\Omega\biggr)k^{2}_{z}S_{x}\biggr\}\\
 & =N^{2}_{n}\biggl[\frac{1}{2}\omega^{2}_{n}S_{x}+\frac{1}{2}\frac{r^{2}}{\Delta}C^{-4}\biggl(\Omega\biggr)k^{2}_{n}C_{x}+\frac{1}{2}C^{-2}\biggl(\Omega\biggr)\biggl(\frac{\Delta}{r^{2}}k^{2}_{y}+k^{2}_{z}\biggr)S_{x}\biggr]\\
 & =N^{2}_{n}\biggl[\frac{1}{2}\omega^{2}_{n}S_{x}+\frac{1}{2}\frac{r^{2}}{\Delta}C^{-4}\biggl(\Omega\biggr)k^{2}_{n}C_{x}+\frac{1}{2}\biggl[\omega^{2}_{n}-\frac{r^{2}}{\Delta}C^{-4}\biggl(\Omega\biggr)k^{2}_{n}\biggr]S_{x}\biggr]\\
 & =N^{2}_{n}\biggl[\omega^{2}_{n}S_{x}+\frac{1}{2}\frac{r^{2}}{\Delta}C^{-4}\biggl(\Omega\biggr)k^{2}_{n}C_{x}-\frac{1}{2}\frac{r^{2}}{\Delta}C^{-4}\biggl(\Omega\biggr)k^{2}_{n}S_{x}\biggr]\\
 & =N^{2}_{n}\biggl[\omega^{2}_{n}S_{x}+\frac{1}{2}\frac{r^{2}}{\Delta}C^{-4}\biggl(\Omega\biggr)k^{2}_{n}\biggl(C_{x}-S_{x}\biggr)\biggr]\\
 & =N^{2}_{n}\biggl[\omega^{2}_{n}S_{x}+\frac{1}{2}\frac{r^{2}}{\Delta}C^{-4}\biggl(\Omega\biggr)k^{2}_{n}\cos\biggl(2k_{n}x\biggr)\biggr]\\
 & =N^{2}_{n}\biggl[\omega^{2}_{n}\sin^{2}\biggl(\frac{n\pi}{L}x\biggr)+\frac{1}{2}\frac{r^{2}}{\Delta}C^{-4}\biggl(\Omega\biggr)\biggl(\frac{n\pi}{L}\biggr)^{2}\cos\biggl(\frac{2n\pi}{L}x\biggr)\biggr]
\end{align*}
 where 
\begin{align*}
 & \hat{g}^{tt}\omega^{2}_{n}+2\hat{g}^{tx}\omega_{n}\beta_{n}+\hat{g}^{xx}\beta^{2}_{n}=\frac{A}{\Delta r^{2}}\omega^{2}_{n}+2\frac{A}{\Delta r}\biggl(\omega_{d}-\Omega\biggr)\omega_{n}\cdot\frac{\omega_{n}C^{2}\biggl(\Omega\biggr)A}{r^{3}}\biggl(\omega_{d}-\Omega\biggr)+\biggl[-\frac{r^{2}}{\Delta}\\
 & \qquad C^{-2}\biggl(\Omega\biggr)\biggr]\biggl[\frac{\omega_{n}C^{2}\biggl(\Omega\biggr)A}{r^{3}}\biggl(\omega_{d}-\Omega\biggr)\biggr]^{2}\\
 & =\frac{A}{\Delta r^{2}}\omega^{2}_{n}+2\frac{A^{2}}{\Delta r^{4}}\biggl(\omega_{d}-\Omega\biggr)^{2}\omega^{2}_{n}C^{2}\biggl(\Omega\biggr)-\frac{r^{2}}{\Delta}C^{-2}\biggl(\Omega\biggr)\frac{\omega^{2}_{n}C^{4}\biggl(\Omega\biggr)A^{2}}{r^{6}}\biggl(\omega_{d}-\Omega\biggr)^{2}\\
 & =\frac{A}{\Delta r^{2}}\omega^{2}_{n}+2\frac{A^{2}}{\Delta r^{4}}\biggl(\omega_{d}-\Omega\biggr)^{2}\omega^{2}_{n}C^{2}\biggl(\Omega\biggr)-\frac{A^{2}}{\Delta r^{4}}\biggl(\omega_{d}-\Omega\biggr)^{2}\omega^{2}_{n}C^{2}\biggl(\Omega\biggr)\\
 & =\frac{A}{\Delta r^{2}}\omega^{2}_{n}+\frac{A^{2}}{\Delta r^{4}}\biggl(\omega_{d}-\Omega\biggr)^{2}\omega^{2}_{n}C^{2}\biggl(\Omega\biggr)
\end{align*}
 and 
\begin{align*}
\frac{A}{\Delta r^{2}}\omega^{2}_{n}+\frac{A^{2}}{\Delta r^{4}}\left(\omega_{d}-\Omega\right)^{2}\omega^{2}_{n}C^{2}\left(\Omega\right) & =\left[\frac{A}{\Delta r^{2}}C^{-2}\left(\Omega\right)+\frac{A^{2}}{\Delta r^{4}}\left(\omega_{d}-\Omega\right)^{2}\right]\omega^{2}_{n}C^{2}\left(\Omega\right)\\
 & =\left\{ \frac{A}{\Delta r^{2}}\cdot\frac{\Delta r^{2}}{A}\left[1-\frac{A^{2}}{\Delta r^{4}}\left(\Omega-\omega_{d}\right)^{2}\right]+\frac{A^{2}}{\Delta r^{4}}\left(\omega_{d}-\Omega\right)^{2}\right\} \omega^{2}_{n}C^{2}\left(\Omega\right)\\
 & =\left[1-\frac{A^{2}}{\Delta r^{4}}\left(\Omega-\omega_{d}\right)^{2}+\frac{A^{2}}{\Delta r^{4}}\left(\omega_{d}-\Omega\right)^{2}\right]\omega^{2}_{n}C^{2}\left(\Omega\right)\\
 & =\omega^{2}_{n}C^{2}\left(\Omega\right)
\end{align*}
 This gives 
\[
\hat{g}^{tt}\omega^{2}_{n}+2\hat{g}^{tx}\omega_{n}\beta_{n}+\hat{g}^{xx}\beta^{2}_{n}=\omega^{2}_{n}C^{2}\left(\Omega\right)
\]
 Also, 
\begin{align*}
\omega^{2}_{n} & =\frac{r^{2}}{\Delta}C^{-4}\left(\Omega\right)\left[k^{2}_{n}+\frac{\Delta}{r^{2}}C^{2}\left(\Omega\right)\left(\frac{\Delta}{r^{2}}k^{2}_{y}+k^{2}_{z}\right)\right]\\
 & =\frac{r^{2}}{\Delta}C^{-4}\left(\Omega\right)k^{2}_{n}+C^{-2}\left(\Omega\right)\left(\frac{\Delta}{r^{2}}k^{2}_{y}+k^{2}_{z}\right)
\end{align*}
 which gives 
\[
C^{-2}\left(\Omega\right)\left(\frac{\Delta}{r^{2}}k^{2}_{y}+k^{2}_{z}\right)=\omega^{2}_{n}-\frac{r^{2}}{\Delta}C^{-4}\left(\Omega\right)k^{2}_{n}
\]
 Finally, 
\[
C_{x}-S_{x}=\cos^{2}\left(k_{n}x\right)-\sin^{2}\left(k_{n}x\right)=\cos\left(2k_{n}x\right)
\]
 by the double-angle identity $\cos2\alpha=\cos^{2}\alpha-\sin^{2}\alpha$.

Alternatively, 
\begin{align*}
\partial_{t}u_{n}\partial_{t}u^{*}_{n}-\frac{1}{2}\hat{g}_{tt}\left(\hat{g}^{\mu\nu}\partial_{\mu}u_{n}\partial_{\nu}u^{*}_{n}\right) & =N^{2}_{n}\left[\omega^{2}_{n}\sin^{2}\left(\frac{n\pi}{L}x\right)+\frac{1}{2}\frac{r^{2}}{\Delta}C^{-4}\left(\Omega\right)\left(\frac{n\pi}{L}\right)^{2}\cos\left(\frac{2n\pi}{L}x\right)\right]\\
 & =N^{2}_{n}\left\{ \omega^{2}_{n}\sin^{2}\left(k_{n}x\right)+\frac{1}{2}\frac{r^{2}}{\Delta}C^{-4}\left(\Omega\right)k^{2}_{n}\left[\cos^{2}\left(k_{n}x\right)-\sin^{2}\left(k_{n}x\right)\right]\right\} \\
 & =N^{2}_{n}\left\{ \left[\omega^{2}_{n}-\frac{1}{2}\frac{r^{2}}{\Delta}C^{-4}\left(\Omega\right)k^{2}_{n}\right]\sin^{2}\left(k_{n}x\right)+\frac{1}{2}\frac{r^{2}}{\Delta}C^{-4}\left(\Omega\right)k^{2}_{n}\cos^{2}\left(k_{n}x\right)\right\} \\
 & =\frac{1}{2}N^{2}_{n}\left\{ \left[2\omega^{2}_{n}-\frac{r^{2}}{\Delta}C^{-4}\left(\Omega\right)k^{2}_{n}\right]\sin^{2}\left(k_{n}x\right)+\frac{r^{2}}{\Delta}C^{-4}\left(\Omega\right)k^{2}_{n}\cos^{2}\left(k_{n}x\right)\right\} \\
 & =\frac{1}{2}N^{2}_{n}\left\{ \left[\omega^{2}_{n}+\frac{\omega_{n}}{b_{n}}\left(\frac{\Delta}{r^{2}}k^{2}_{y}+k^{2}_{z}\right)\right]\sin^{2}\left(k_{n}x\right)+\frac{\omega^{2}_{n}r^{2}}{b^{2}_{n}\Delta}k^{2}_{n}\cos^{2}\left(k_{n}x\right)\right\} \\
 & =\frac{1}{2}N^{2}_{n}\left\{ \mathcal{F}_{n}\sin^{2}\left(k_{n}x\right)+\mathcal{G}_{n}\cos^{2}\left(k_{n}x\right)\right\} 
\end{align*}
 where we define 
\[
\mathcal{F}_{n}=\omega^{2}_{n}+\frac{\omega_{n}}{b_{n}}\left(\frac{\Delta}{r^{2}}k^{2}_{y}+k^{2}_{z}\right)
\]
 {}
\[
\mathcal{G}_{n}=\frac{\omega^{2}_{n}r^{2}}{b^{2}_{n}\Delta}k^{2}_{n}
\]
 We also have 
\[
\omega_{n}=\frac{r}{\sqrt{\Delta}}C^{-2}\left(\Omega\right)\sqrt{k^{2}_{n}+\frac{\Delta}{r^{2}}C^{2}\left(\Omega\right)\left(\frac{\Delta}{r^{2}}k^{2}_{y}+k^{2}_{z}\right)}
\]
 {}
\[
\omega^{2}_{n}=\frac{r^{2}}{\Delta}C^{-4}\left(\Omega\right)k^{2}_{n}+C^{-2}\left(\Omega\right)\left(\frac{\Delta}{r^{2}}k^{2}_{y}+k^{2}_{z}\right)
\]
 {}
\[
\beta_{n}=\frac{b_{n}A}{r^{3}}\left(\omega_{d}-\Omega\right)
\]
 {}
\[
b_{n}=\omega_{n}C^{2}\left(\Omega\right)
\]
 {}
\[
k_{n}=\frac{n\pi}{L}
\]
 {}
\begin{align*}
2\omega^{2}_{n}-\frac{r^{2}}{\Delta}C^{-4}\left(\Omega\right)k^{2}_{n} & =\omega^{2}_{n}+\omega^{2}_{n}-\frac{r^{2}}{\Delta}C^{-4}\left(\Omega\right)k^{2}_{n}\\
 & =\omega^{2}_{n}+\left[\frac{r^{2}}{\Delta}C^{-4}\left(\Omega\right)k^{2}_{n}+C^{-2}\left(\Omega\right)\left(\frac{\Delta}{r^{2}}k^{2}_{y}+k^{2}_{z}\right)\right]-\frac{r^{2}}{\Delta}C^{-4}\left(\Omega\right)k^{2}_{n}\\
 & =\omega^{2}_{n}+C^{-2}\left(\Omega\right)\left(\frac{\Delta}{r^{2}}k^{2}_{y}+k^{2}_{z}\right)\\
 & =\omega^{2}_{n}+\frac{\omega_{n}}{\omega_{n}C^{2}\left(\Omega\right)}\left(\frac{\Delta}{r^{2}}k^{2}_{y}+k^{2}_{z}\right)\\
 & =\omega^{2}_{n}+\frac{\omega_{n}}{b_{n}}\left(\frac{\Delta}{r^{2}}k^{2}_{y}+k^{2}_{z}\right)
\end{align*}
 {}
\begin{align*}
\frac{r^{2}}{\Delta}C^{-4}\left(\Omega\right)k^{2}_{n} & =\frac{\omega^{2}_{n}r^{2}}{\omega^{2}_{n}C^{4}\left(\Omega\right)\Delta}k^{2}_{n}\\
 & =\frac{\omega^{2}_{n}r^{2}}{b^{2}_{n}\Delta}k^{2}_{n}
\end{align*}

We obtain 
\begin{align*}
 & \langle\epsilon_{\text{vac}}\rangle=C^{2}\biggl(\Omega\biggr)\frac{1}{L\int^{+\infty}_{-\infty}\int^{+\infty}_{-\infty}dydz}\int^{L}_{0}dx\int^{+\infty}_{-\infty}\int^{+\infty}_{-\infty}dydz\langle0\vert T_{tt}\vert0\rangle\\
 & =C^{2}\biggl(\Omega\biggr)\frac{1}{L\int^{+\infty}_{-\infty}\int^{+\infty}_{-\infty}dydz}\int^{L}_{0}dx\int^{+\infty}_{-\infty}\int^{+\infty}_{-\infty}dydz\sum_{n}\int d^{2}k_{\perp}\biggl[\partial_{t}u_{n}\partial_{t}u^{*}_{n}-\frac{1}{2}\hat{g}_{tt}\biggl(\hat{g}^{\mu\nu}\partial_{\mu}u_{n}\partial_{\nu}u^{*}_{n}\biggr)\biggr]\\
 & =C^{2}\biggl(\Omega\biggr)\frac{1}{L\int^{+\infty}_{-\infty}\int^{+\infty}_{-\infty}dydz}\int^{L}_{0}dx\int^{+\infty}_{-\infty}\int^{+\infty}_{-\infty}dydz\sum_{n}\int d^{2}k_{\perp}\biggl\{ N^{2}_{n}\biggl[\omega^{2}_{n}\sin^{2}\biggl(\frac{n\pi}{L}x\biggr)+\frac{1}{2}\\
 & \qquad\frac{r^{2}}{\Delta}C^{-4}\biggl(\Omega\biggr)\biggl(\frac{n\pi}{L}\biggr)^{2}\cos\biggl(\frac{2n\pi}{L}x\biggr)\biggr]\biggr\}\\
 & =\frac{C^{2}\biggl(\Omega\biggr)}{L}\int^{L}_{0}dx\sum_{n}\int d^{2}k_{\perp}N^{2}_{n}\biggl[\omega^{2}_{n}\sin^{2}\biggl(\frac{n\pi}{L}x\biggr)+\frac{1}{2}\frac{r^{2}}{\Delta}C^{-4}\biggl(\Omega\biggr)\biggl(\frac{n\pi}{L}\biggr)^{2}\cos\biggl(\frac{2n\pi}{L}x\biggr)\biggr]\\
 & =\frac{C^{2}\biggl(\Omega\biggr)}{L}\sum_{n}\int d^{2}k_{\perp}N^{2}_{n}\biggl[\omega^{2}_{n}\int^{L}_{0}\sin^{2}\biggl(\frac{n\pi}{L}x\biggr)dx+\frac{1}{2}\frac{r^{2}}{\Delta}C^{-4}\biggl(\Omega\biggr)\biggl(\frac{n\pi}{L}\biggr)^{2}\int^{L}_{0}\cos\biggl(\frac{2n\pi}{L}x\biggr)dx\biggr]\\
 & =\frac{C^{2}\biggl(\Omega\biggr)}{L}\sum_{n}\int d^{2}k_{\perp}N^{2}_{n}\omega^{2}_{n}\frac{L}{2}\\
 & =\frac{C^{2}\biggl(\Omega\biggr)}{2}\sum_{n}\int d^{2}k_{\perp}N^{2}_{n}\omega^{2}_{n}
\end{align*}
 Here, $\langle0\vert T_{tt}\vert0\rangle$ is independent of $y,z$,
so the integral $\int^{+\infty}_{-\infty}\int^{+\infty}_{-\infty}dydz$
can be factored out. We have $\int^{L}_{0}\sin^{2}\left(\frac{n\pi}{L}x\right)dx=\int^{L}_{0}\frac{1-\cos\left(2\frac{n\pi}{L}x\right)}{2}dx=\frac{L}{2}$
and $\int^{L}_{0}\cos\left(\frac{2n\pi}{L}x\right)dx=\left.\frac{L}{2n\pi}\sin\left(\frac{2n\pi}{L}x\right)\right|^{L}_{0}=\frac{L}{2n\pi}\left[\sin(2n\pi)-0\right]=0$
(since $\sin2n\pi=0$). Also, $N^{2}_{n}=\frac{1}{\left(2\pi\right)^{2}L\omega_{n}C^{2}\left(\Omega\right)}$.
\begin{align*}
\langle\epsilon_{\text{vac}}\rangle & =\frac{C^{2}\left(\Omega\right)}{2}\sum_{n}\int d^{2}k_{\perp}N^{2}_{n}\omega^{2}_{n}\\
 & =\frac{C^{2}\left(\Omega\right)}{2}\sum_{n}\int d^{2}k_{\perp}\frac{1}{\left(2\pi\right)^{2}L\omega_{n}C^{2}\left(\Omega\right)}\omega^{2}_{n}\\
 & =\frac{1}{2\left(2\pi\right)^{2}L}\sum_{n}\int d^{2}k_{\perp}\omega_{n}\\
 & =\frac{1}{2\left(2\pi\right)^{2}L}\sum_{n}\int d^{2}k_{\perp}\left[\frac{r}{\sqrt{\Delta}}C^{-2}\left(\Omega\right)\sqrt{\left(\frac{n\pi}{L}\right)^{2}+\frac{\Delta}{r^{2}}C^{2}\left(\Omega\right)\left(\frac{\Delta}{r^{2}}k^{2}_{y}+k^{2}_{z}\right)}\right]\\
 & =\frac{1}{2\left(2\pi\right)^{2}L}\sum_{n}\int d^{2}k_{\perp}\left\{ \frac{r}{\sqrt{\Delta}}C^{-2}\left(\Omega\right)\sqrt{\frac{\Delta}{r^{2}}C^{2}\left(\Omega\right)\left[\frac{r^{2}}{\Delta}C^{-2}\left(\Omega\right)\left(\frac{n\pi}{L}\right)^{2}+\left(\frac{\Delta}{r^{2}}k^{2}_{y}+k^{2}_{z}\right)\right]}\right\} \\
 & =\frac{1}{2\left(2\pi\right)^{2}L}\sum_{n}\int d^{2}k_{\perp}\left[\frac{r}{\sqrt{\Delta}}C^{-2}\left(\Omega\right)\cdot\frac{\sqrt{\Delta}}{r}C\left(\Omega\right)\sqrt{\frac{r^{2}}{\Delta}C^{-2}\left(\Omega\right)\left(\frac{n\pi}{L}\right)^{2}+\frac{\Delta}{r^{2}}k^{2}_{y}+k^{2}_{z}}\right]\\
 & =\frac{1}{2\left(2\pi\right)^{2}L}\sum_{n}\int d^{2}k_{\perp}\left[C^{-1}\left(\Omega\right)\sqrt{\frac{r^{2}}{\Delta}C^{-2}\left(\Omega\right)\left(\frac{n\pi}{L}\right)^{2}+\frac{\Delta}{r^{2}}k^{2}_{y}+k^{2}_{z}}\right]
\end{align*}
 Define 
\[
s=\frac{r^{2}}{\Delta}C^{-2}\left(\Omega\right)\left(\frac{n\pi}{L}\right)^{2}
\]
 {}
\[
\tilde{k}_{y}=\frac{\sqrt{\Delta}}{r}k_{y}
\]
 Then 
\begin{align*}
\sum_{n}\int d^{2}k_{\perp}\omega_{n} & =\sum_{n}\int d^{2}k_{\perp}\left[C^{-1}\left(\Omega\right)\sqrt{\frac{r^{2}}{\Delta}C^{-2}\left(\Omega\right)\left(\frac{n\pi}{L}\right)^{2}+\frac{\Delta}{r^{2}}k^{2}_{y}+k^{2}_{z}}\right]\\
 & =C^{-1}\left(\Omega\right)\sum_{n}\int d^{2}k_{\perp}\sqrt{s+\tilde{k}^{2}_{y}+k^{2}_{z}}\\
 & =\frac{r}{C\left(\Omega\right)\sqrt{\Delta}}\sum_{n}\int d^{2}\tilde{k}_{\perp}\left(s+\tilde{k}^{2}_{y}+k^{2}_{z}\right)^{1/2}
\end{align*}
 where $k_{y}=\frac{r}{\sqrt{\Delta}}\tilde{k}_{y}$, $dk_{y}=d\left(\frac{r}{\sqrt{\Delta}}\tilde{k}_{y}\right)$,
and $d^{2}k_{\perp}=dk_{y}dk_{z}=d\left(\frac{r}{\sqrt{\Delta}}\tilde{k}_{y}\right)dk_{z}=\frac{r}{\sqrt{\Delta}}d\tilde{k}_{y}dk_{z}$;
also, $d^{2}\tilde{k}_{\perp}=d\tilde{k}_{y}dk_{z}$. Apply the Schwinger
proper-time method to evaluate the momentum-space integral: 
\begin{align*}
\int d^{2}\tilde{k}_{\perp}\left(s+\tilde{k}^{2}_{y}+k^{2}_{z}\right)^{1/2} & =\int d^{2}\tilde{k}_{\perp}\frac{1}{\Gamma\left(-1/2\right)}\int^{\infty}_{0}dtt^{-\frac{1}{2}-1}e^{-t\left(s+\tilde{k}^{2}_{y}+k^{2}_{z}\right)}\\
 & =\int d^{2}\tilde{k}_{\perp}\left[-\frac{1}{2\sqrt{\pi}}\int^{\infty}_{0}dtt^{-\frac{3}{2}}e^{-t\left(s+\tilde{k}^{2}_{y}+k^{2}_{z}\right)}\right]\\
 & =-\frac{1}{2\sqrt{\pi}}\int^{\infty}_{0}dtt^{-\frac{3}{2}}\int^{\infty}_{-\infty}e^{-t\tilde{k}^{2}_{y}}d\tilde{k}_{y}\int^{\infty}_{-\infty}e^{-tk^{2}_{z}}dk_{z}e^{-ts}\\
 & =-\frac{1}{2\sqrt{\pi}}\int^{\infty}_{0}dtt^{-\frac{3}{2}}\frac{\pi}{t}e^{-ts}\\
 & =-\frac{\sqrt{\pi}}{2}\int^{\infty}_{0}dtt^{-\frac{3}{2}-1}e^{-ts}\\
 & =-\frac{\sqrt{\pi}}{2}s^{3/2}\Gamma\left(-\frac{3}{2}\right)\\
 & =-\frac{\sqrt{\pi}}{2}s^{3/2}\frac{4\sqrt{\pi}}{3}\\
 & =-\frac{2\pi}{3}s^{3/2}
\end{align*}
 Here, after analytic continuation, $A^{-\alpha}=\frac{1}{\Gamma\left(\alpha\right)}\int^{\infty}_{0}dtt^{\alpha-1}e^{-tA}$;
$\frac{1}{\Gamma\left(-1/2\right)}=-\frac{1}{2\sqrt{\pi}}$; $\int^{\infty}_{-\infty}e^{-t\tilde{k}^{2}_{y}}d\tilde{k}_{y}=\int^{\infty}_{-\infty}e^{-tk^{2}_{z}}dk_{z}=\sqrt{\frac{\pi}{t}}$;
$\int^{\infty}_{0}t^{z-1}e^{-ts}dt=s^{-z}\Gamma\left(z\right)$; and
$\Gamma\left(-\frac{3}{2}\right)=\frac{4\sqrt{\pi}}{3}$. Hence, 
\begin{align*}
\sum_{n}\int d^{2}k_{\perp}\omega_{n} & =\frac{r}{C\left(\Omega\right)\sqrt{\Delta}}\sum_{n}\int d^{2}\tilde{k}_{\perp}\left(s+\tilde{k}^{2}_{y}+k^{2}_{z}\right)^{1/2}\\
 & =\frac{r}{C\left(\Omega\right)\sqrt{\Delta}}\sum_{n}\left(-\frac{2\pi}{3}s^{3/2}\right)\\
 & =\frac{r}{C\left(\Omega\right)\sqrt{\Delta}}\sum_{n}\left\{ -\frac{2\pi}{3}\left[\frac{r^{2}}{\Delta}C^{-2}\left(\Omega\right)\left(\frac{n\pi}{L}\right)^{2}\right]^{3/2}\right\} \\
 & =\frac{r^{4}}{C^{4}\left(\Omega\right)\Delta^{2}}\sum_{n}\left[-\frac{2\pi}{3}\left(\frac{n\pi}{L}\right)^{3}\right]\\
 & =-\frac{2\pi^{4}}{3L^{3}}\frac{r^{4}}{C^{4}\left(\Omega\right)\Delta^{2}}\sum_{n}n^{3}
\end{align*}
 Applying Riemann $\zeta$-function regularization, $\sum^{\infty}_{n=1}n^{3}\rightarrow\zeta\left(-3\right)=\frac{1}{120}$.
We obtain 
\begin{align*}
\sum_{n}\int d^{2}k_{\perp}\omega_{n} & =-\frac{2\pi^{4}}{3L^{3}}\frac{r^{4}}{C^{4}\left(\Omega\right)\Delta^{2}}\cdot\frac{1}{120}\\
 & =-\frac{\pi^{4}}{180L^{3}}\frac{r^{4}}{C^{4}\left(\Omega\right)\Delta^{2}}
\end{align*}
 and consequently 
\begin{align*}
\langle\epsilon_{\text{vac}}\rangle & =\frac{1}{2\left(2\pi\right)^{2}L}\sum_{n}\int d^{2}k_{\perp}\omega_{n}\\
 & =\frac{1}{8\pi^{2}L}\left(-\frac{\pi^{4}}{180L^{3}}\frac{r^{4}}{C^{4}\left(\Omega\right)\Delta^{2}}\right)\\
 & =-\frac{\pi^{2}}{1440L^{4}}\frac{r^{4}}{C^{4}\left(\Omega\right)\Delta^{2}}
\end{align*}
 The proper length of the Casimir cavity is obtained as follows. For
$ds^{2}=\hat{g}_{tt}d\tilde{t}^{2}+\left(\hat{g}_{xx}-\frac{\hat{g}^{2}_{tx}}{\hat{g}_{tt}}\right)dx^{2}+\hat{g}_{yy}dy^{2}+\hat{g}_{zz}dz^{2}$
with $d\tilde{t}=dy=dz=0$, a spacelike interval is negative under
the signature $(+,-,-,-)$, so the absolute value must be taken under
the square root: $ds=\sqrt{\left|\hat{g}_{xx}-\frac{\hat{g}^{2}_{tx}}{\hat{g}_{tt}}\right|}dx$.
Thus, 
\[
L_{p}=\int^{L}_{0}ds=\int^{L}_{0}\sqrt{\left|\hat{g}_{xx}-\frac{\hat{g}^{2}_{tx}}{\hat{g}_{tt}}\right|}dx
\]
 With the metric approximated as constant, 
\begin{align*}
L_{p} & =\sqrt{\left|\hat{g}_{xx}-\frac{\hat{g}^{2}_{tx}}{\hat{g}_{tt}}\right|}\int^{L}_{0}dx\\
 & =L\sqrt{\left|\hat{g}_{xx}-\frac{\hat{g}^{2}_{tx}}{\hat{g}_{tt}}\right|}\\
 & =L\sqrt{\left|-\frac{\Delta}{r^{2}}C^{2}\left(\Omega\right)\right|}\\
 & =L\sqrt{\frac{\Delta}{r^{2}}C^{2}\left(\Omega\right)}\\
 & =L\frac{\sqrt{\Delta}}{r}C\left(\Omega\right)
\end{align*}
 where 
\begin{align*}
\hat{g}_{xx}-\frac{\hat{g}^{2}_{tx}}{\hat{g}_{tt}} & =-\frac{A}{r^{4}}-\frac{\left[-\frac{A}{r^{3}}\left(\Omega-\omega_{d}\right)\right]^{2}}{C^{-2}\left(\Omega\right)}\\
 & =-\frac{A}{r^{4}}-\frac{\frac{A^{2}}{r^{6}}\left(\Omega-\omega_{d}\right)^{2}}{C^{-2}\left(\Omega\right)}\\
 & =-\frac{A}{r^{4}}-\frac{\frac{A^{2}}{r^{6}}\left(\Omega-\omega_{d}\right)^{2}}{\frac{\Delta r^{2}}{A}\left[1-\frac{A^{2}}{\Delta r^{4}}\left(\Omega-\omega_{d}\right)^{2}\right]}\\
 & =-\frac{A}{r^{4}}-\frac{A}{r^{4}}\frac{\frac{A^{2}}{\Delta r^{4}}\left(\Omega-\omega_{d}\right)^{2}}{1-\frac{A^{2}}{\Delta r^{4}}\left(\Omega-\omega_{d}\right)^{2}}\\
 & =\frac{A}{r^{4}}\left[-1+\frac{-\frac{A^{2}}{\Delta r^{4}}\left(\Omega-\omega_{d}\right)^{2}}{1-\frac{A^{2}}{\Delta r^{4}}\left(\Omega-\omega_{d}\right)^{2}}\right]\\
 & =\frac{A}{r^{4}}\left[-1+\left(1-\frac{1}{1-\frac{A^{2}}{\Delta r^{4}}\left(\Omega-\omega_{d}\right)^{2}}\right)\right]\\
 & =-\frac{A}{r^{4}}\frac{1}{1-\frac{A^{2}}{\Delta r^{4}}\left(\Omega-\omega_{d}\right)^{2}}\\
 & =-\frac{\Delta}{r^{2}}\frac{1}{\frac{\Delta r^{2}}{A}\left[1-\frac{A^{2}}{\Delta r^{4}}\left(\Omega-\omega_{d}\right)^{2}\right]}\\
 & =-\frac{\Delta}{r^{2}}\frac{1}{C^{-2}\left(\Omega\right)}\\
 & =-\frac{\Delta}{r^{2}}C^{2}\left(\Omega\right)
\end{align*}
 Here, 
\begin{align*}
\frac{-\frac{A^{2}}{\Delta r^{4}}\left(\Omega-\omega_{d}\right)^{2}}{1-\frac{A^{2}}{\Delta r^{4}}\left(\Omega-\omega_{d}\right)^{2}} & =\frac{1-\frac{A^{2}}{\Delta r^{4}}\left(\Omega-\omega_{d}\right)^{2}}{1-\frac{A^{2}}{\Delta r^{4}}\left(\Omega-\omega_{d}\right)^{2}}-\frac{1}{1-\frac{A^{2}}{\Delta r^{4}}\left(\Omega-\omega_{d}\right)^{2}}\\
 & =1-\frac{1}{1-\frac{A^{2}}{\Delta r^{4}}\left(\Omega-\omega_{d}\right)^{2}}
\end{align*}
 {}
\[
C^{-2}\left(\Omega\right)=\frac{\Delta r^{2}}{A}\left[1-\frac{A^{2}}{\Delta r^{4}}\left(\Omega-\omega_{d}\right)^{2}\right]
\]
 We obtain 
\[
L=L_{p}\frac{r}{C\left(\Omega\right)\sqrt{\Delta}}
\]
 {}
\[
L^{4}=L^{4}_{p}\frac{r^{4}}{C^{4}\left(\Omega\right)\Delta^{2}}
\]
 The Casimir energy per unit volume is therefore 
\begin{align*}
\langle\epsilon_{\text{vac}}\rangle & =-\frac{\pi^{2}}{1440L^{4}}\frac{r^{4}}{C^{4}\left(\Omega\right)\Delta^{2}}\\
 & =-\frac{\pi^{2}}{1440\left[L^{4}_{p}\frac{r^{4}}{C^{4}\left(\Omega\right)\Delta^{2}}\right]}\frac{r^{4}}{C^{4}\left(\Omega\right)\Delta^{2}}\\
 & =-\frac{\pi^{2}}{1440L^{4}_{p}}\frac{C^{4}\left(\Omega\right)\Delta^{2}}{r^{4}}\frac{r^{4}}{C^{4}\left(\Omega\right)\Delta^{2}}\\
 & =-\frac{\pi^{2}}{1440L^{4}_{p}}.
\end{align*}
 This is precisely the expression for the Casimir energy density in
flat spacetime (Minkowski spacetime).

\Needspace{110pt}\par\addvspace{20pt}\noindent\begin{minipage}{\linewidth}\begingroup\fontsize{12}{15}\selectfont\setlength{\abovedisplayskip}{0pt}\setlength{\belowdisplayskip}{0pt}\setlength{\abovedisplayshortskip}{0pt}\setlength{\belowdisplayshortskip}{0pt}\phantomsection\addcontentsline{toc}{section}{VI. Further Perspectives: Geometry and Boundary Conditions in the Constant-Metric Approximation}
\[
\begin{gathered}\text{\bfseries VI. Further Perspectives: Geometry and Boundary Conditions}\\
\text{\bfseries in the Constant-Metric Approximation}
\end{gathered}
\]
\endgroup\end{minipage}\par\nopagebreak[4]\vspace{20pt}\nopagebreak[4]

Once the metric is approximated as constant, the spacetime is in fact
Minkowski spacetime. A linear coordinate transformation with constant
coefficients brings the line element into the standard Minkowski form.
Let $\bar{g}_{\mu\nu}$ denote the metric components after this approximation.
Completing the square gives 
\[
ds^{2}=\bar{g}_{tt}\left(dt+\frac{\bar{g}_{tx}}{\bar{g}_{tt}}dx\right)^{2}+\left(\bar{g}_{xx}-\frac{\bar{g}^{2}_{tx}}{\bar{g}_{tt}}\right)dx^{2}+\bar{g}_{yy}dy^{2}+\bar{g}_{zz}dz^{2},
\]
 where all metric coefficients on the right-hand side are fixed constants.
Set 
\begin{align*}
dT & =\sqrt{\bar{g}_{tt}}\left(dt+\frac{\bar{g}_{tx}}{\bar{g}_{tt}}dx\right),\\
dX & =\sqrt{-\left(\bar{g}_{xx}-\frac{\bar{g}^{2}_{tx}}{\bar{g}_{tt}}\right)}dx,\\
dY & =\sqrt{-\bar{g}_{yy}}dy,\\
dZ & =\sqrt{-\bar{g}_{zz}}dz,
\end{align*}
 or apply the affine transformation 
\begin{align*}
T-T_{0} & =\sqrt{\bar{g}_{tt}}\left[\left(t-t_{0}\right)+\frac{\bar{g}_{tx}}{\bar{g}_{tt}}\left(x-x_{0}\right)\right],\\
X-X_{0} & =\sqrt{-\left(\bar{g}_{xx}-\frac{\bar{g}^{2}_{tx}}{\bar{g}_{tt}}\right)}\left(x-x_{0}\right),\\
Y-Y_{0} & =\sqrt{-\bar{g}_{yy}}\left(y-y_{0}\right),\\
Z-Z_{0} & =\sqrt{-\bar{g}_{zz}}\left(z-z_{0}\right),
\end{align*}
 Here, all radicands are greater than $0$ when the worldline of the
cavity center is timelike and $r_{0}>r_{+}$, $\theta_{0}=\pi/2$.
All coordinate values with subscript $0$ are constants and can be
set to $0$ by suitable choices of the respective coordinate origins.
We then obtain 
\[
ds^{2}=dT^{2}-dX^{2}-dY^{2}-dZ^{2}.
\]
 In this coordinate system, the Dirichlet boundary conditions become
\[
\Psi\left(T,X=0,Y,Z\right)=\Psi\left(T,X=\frac{\sqrt{\Delta}}{r}C\left(\Omega\right)L,Y,Z\right)=0,
\]
 When 
\[
x=L
\]
 we have 
\begin{align*}
X & =\sqrt{-\left(\hat{g}_{xx}-\frac{\hat{g}^{2}_{tx}}{\hat{g}_{tt}}\right)}\cdot L\\
 & =\frac{\sqrt{\Delta}}{r}C\left(\Omega\right)L.
\end{align*}
 This also explains why, after transforming the line element into
the standard Minkowski form, the positions of the two parallel plates
are no longer both expressed in the same form as $X=0$ and $X=L$.

From the viewpoint of spacetime curvature, approximating the metric
components as constants gives 
\[
\partial_{\lambda}\bar{g}_{\mu\nu}=0,
\]
 {}
\[
\bar{\Gamma}^{\sigma}_{\!\mu\nu}=\frac{1}{2}\bar{g}^{\sigma\rho}\left(\partial_{\mu}\bar{g}_{\nu\rho}+\partial_{\nu}\bar{g}_{\rho\mu}-\partial_{\rho}\bar{g}_{\mu\nu}\right)=0,
\]
 {}
\[
\bar{R}^{\rho}_{\!\sigma\mu\nu}=\partial_{\mu}\bar{\Gamma}^{\rho}_{\!\nu\sigma}-\partial_{\nu}\bar{\Gamma}^{\rho}_{\!\mu\sigma}+\bar{\Gamma}^{\rho}_{\!\mu\lambda}\bar{\Gamma}^{\lambda}_{\!\nu\sigma}-\bar{\Gamma}^{\rho}_{\!\nu\lambda}\bar{\Gamma}^{\lambda}_{\!\mu\sigma}=0,
\]
 Thus, every component of the Riemann curvature tensor is $0$. This
is precisely what it means for spacetime to be flat, and its physical
properties are naturally those of flat spacetime. Unless otherwise
necessary, the corresponding physical problem can therefore be studied
entirely in flat spacetime.

For parallel plates in flat spacetime undergoing circular motion about
an external center, the Casimir energy need not be the same as when
the plates are at rest in an inertial frame, because their motion
is noninertial. In the present analysis and in the original paper
under discussion, however, this motion is effectively approximated
as inertial. This can be understood in two ways for parallel plates
in circular motion. First, in a coordinate system at rest relative
to the plates, the metric depends on position. Second, in standard
inertial Minkowski coordinates, the boundary conditions depend on
time. In the present analysis, adopting the constant-metric approximation
allows a transformation to the standard Minkowski form of the metric,
while time-independent boundary conditions similar to those for plates
at rest in an inertial frame are still imposed. This clearly describes
the physical situation of stationary parallel plates, rather than
that of parallel plates in circular motion.
\end{document}
